% Réduction des risques liés aux catastrophes naturelles — SVT Tle D — Cours signé OnBuch+
% Programme officiel MINESEC. Compiler avec : tectonic N15-catastrophes-naturelles-mouvements-lithosphere.tex
\newcommand{\DOCMATIERE}{SVT}
\newcommand{\DOCNIVEAU}{Tle D}
\newcommand{\DOCMODULE}{Module III : L'Éducation à l'Environnement et au Développement Durable — Les mouvements de la lithosphère et leurs conséquences}
\newcommand{\DOCLECON}{N15}
\newcommand{\DOCTITRE}{Réduction des risques liés aux catastrophes naturelles}
% OnBuch+ — préambule « cours signés » ENRICHI (compilé avec tectonic / XeLaTeX)
% Système visuel « Pop » d'OnBuch+ : fond crème (jamais blanc), bordures encre
% épaisses, ombres « dures » décalées, titres Archivo Black en capitales.
% Version MATHÉMATIQUES : graphes (pgfplots/tikz), boîtes pédagogiques
% (définition, propriété, méthode, attention, le savais-tu, exercice résolu,
% exercice, TP, bilan), page de garde et signature OnBuch+ ancrée en bas.
%
% Macros attendues dans main.tex AVANT \input{preamble} :
%   \DOCMATIERE  \DOCNIVEAU  \DOCMODULE  \DOCLECON (numéro)  \DOCTITRE
\documentclass[11pt]{article}

\usepackage{fontspec}
\usepackage[french]{babel}
\usepackage[a4paper,margin=2cm,top=3.4cm,bottom=2.8cm,headheight=1.4cm,footskip=1.3cm]{geometry}
\usepackage{amsmath,amssymb,amsfonts}
\usepackage{mathtools} % \xmapsto, \coloneqq... souvent utilisés par les modèles
\usepackage{cancel} % simplifications barrées (bilans de forces)
% \abs{z} (module, valeur absolue) et \norm{u} (norme), utilisés par les modèles
\providecommand{\abs}{}\let\abs\relax\DeclarePairedDelimiter{\abs}{\lvert}{\rvert}
\providecommand{\norm}{}\let\norm\relax\DeclarePairedDelimiter{\norm}{\lVert}{\rVert}
\usepackage[table]{xcolor} % [table] : fournit aussi \rowcolors (alternance de lignes)
\usepackage{pagecolor}
\usepackage{tikz}
\usetikzlibrary{arrows.meta,positioning,calc,shapes.geometric,shapes.misc,shapes.arrows,decorations.pathmorphing,decorations.pathreplacing,decorations.markings,patterns,shadows,mindmap,trees,fit,backgrounds,matrix,intersections,angles,quotes}
\usepackage{pgfplots}
\usepackage[version=4]{mhchem} % équations nucléaires \ce{^{235}_{92}U}
\usepackage[european,siunitx]{circuitikz} % schémas électriques
\pgfplotsset{compat=1.18}
\usepgfplotslibrary{fillbetween,groupplots} % aires entre courbes (calcul intégral)
\usepackage{enumitem}
\usepackage{fancyhdr}
% [most] charge toutes les bibliothèques tcolorbox (listings, minted, raster,
% documentation...) et gonfle inutilement la mémoire TeX sur un document long
% (~300 boîtes) : on ne charge que ce qui est réellement utilisé.
\usepackage[skins,breakable]{tcolorbox}
\usepackage{graphicx}
\usepackage{colortbl}
\usepackage{booktabs}
\usepackage{tabularx}
\usepackage{diagbox} % case d'en-tête diagonale des tableaux à double entrée
% Colonnes extensibles centrée / gauche / droite (C, L, R), souvent utilisées par les modèles
\newcolumntype{C}{>{\centering\arraybackslash}X}
\newcolumntype{L}{>{\raggedright\arraybackslash}X}
\newcolumntype{R}{>{\raggedleft\arraybackslash}X}
\usepackage{array}
\usepackage{multicol}
\usepackage{siunitx}
% parse-numbers=false : accepte aussi \SI{2,5 \times 10^{-3}}{...}
\sisetup{locale=FR,output-decimal-marker={,},group-minimum-digits=4,parse-numbers=false}
\DeclareSIUnit\ohm{\ensuremath{\symup{\Omega}}} % U+2126 absent de la police
\DeclareSIUnit\division{div} % réglages d'oscilloscope (ms/div, V/div)
\DeclareSIUnit\tr{tr} % tours (vitesse de rotation, ex. \SI{1500}{\tr\per\minute})
\DeclareSIUnit\molar{M} % concentration molaire (ex. \SI{3}{\molar}, \SI{1}{\micro\molar})
\DeclareSIUnit\an{an} % année (le modèle confond parfois avec \year, un compteur TeX) — ex. \SI{5}{\kilo\meter\per\an}
\DeclareSIUnit\FCFA{FCFA} % franc CFA (ex. \SI{500}{\FCFA\per\kilo\gram})
\DeclareSIUnit\degreeBrix{\textdegree Brix} % teneur en sucre d'un sirop/jus (réfractométrie)
\DeclareSIUnit\month{mois} % le modèle utilise parfois \month, un compteur TeX, comme unité de durée
\usepackage{qrcode}
% \degree : commande fréquemment utilisée par les modèles pour un angle ou une
% température en dehors de \SI{}{} (non fournie par les paquets chargés).
\providecommand{\degree}{^{\circ}}
% \qed (fin de démonstration), utilisé par les modèles sans amsthm
\providecommand{\qed}{\hfill$\square$}
% \barre : double barre des valeurs interdites dans un tableau de variations (tkz-tab)
\providecommand{\barre}{\ensuremath{\|}}
% environnement proof (amsthm non chargé) utilisé par les modèles
\ifdefined\proof\else\newenvironment{proof}[1][Démonstration]{\par\noindent\textit{#1.}\ }{\qed\par}\fi
\usepackage{lastpage}
\usepackage{hyperref}
\hypersetup{hidelinks}

% ============================================================
% Palette « Pop » OnBuch+ — mêmes hex que la classe P (plus_ui.dart)
% ============================================================
\definecolor{poporange}{HTML}{F59321}
\definecolor{popdark}{HTML}{E07A0C}
\definecolor{popcream}{HTML}{FFF6E8}
\definecolor{popink}{HTML}{1C1714}
\definecolor{poppurple}{HTML}{7A5AE0}
\definecolor{popgreen}{HTML}{2E9E6A}
\definecolor{poppink}{HTML}{E0457B}
\definecolor{popblue}{HTML}{2F7DE1}
\definecolor{popgold}{HTML}{C98A12}
\definecolor{popmuted}{HTML}{8A7F76}
\definecolor{popline}{HTML}{EADFD2}
\definecolor{popgray}{HTML}{8A7F76}
\colorlet{popgrey}{popgray}
% Alias tolérants (le modèle nomme parfois les couleurs différemment)
\colorlet{popwhite}{white}
\colorlet{popblack}{popink}
\colorlet{popred}{poppink}
\colorlet{popyellow}{popgold}
% Teintes claires (fonds de tableaux/encadrés) — le modèle nomme parfois
% les couleurs avec un suffixe L (ex. popblueL) sans les avoir définies.
\colorlet{poporangeL}{poporange!20}
\colorlet{popdarkL}{popdark!20}
\colorlet{poppurpleL}{poppurple!20}
\colorlet{popgreenL}{popgreen!20}
\colorlet{poppinkL}{poppink!20}
\colorlet{popblueL}{popblue!20}
\colorlet{popgoldL}{popgold!20}
\colorlet{popredL}{popred!20}
\colorlet{popgrayL}{popgray!20}
% Teintes claires pour fonds de tableaux / zones
\colorlet{popblueL}{popblue!10!white}
\colorlet{poporangeL}{poporange!14!white}
\colorlet{popgreenL}{popgreen!12!white}
\colorlet{poppurpleL}{poppurple!12!white}
\colorlet{poppinkL}{poppink!12!white}
\colorlet{popgoldL}{popgold!14!white}

\pagecolor{popcream}

% ============================================================
% Typographie
% ============================================================
\defaultfontfeatures{Ligatures=TeX}
\setmainfont{PlusJakartaSans-Medium.ttf}[Path=../../../fonts/,BoldFont=PlusJakartaSans-Bold.ttf]
% Maths en sans-serif (Fira Math) avec chiffres et lettres droites de Plus
% Jakarta, pour que les formules se fondent dans le texte.
\usepackage[math-style=ISO,bold-style=ISO]{unicode-math}
\setmathfont{FiraMath-Regular.otf}[Path=../../../fonts/]
\setmathfont{PlusJakartaSans-Medium.ttf}[Path=../../../fonts/,range={up/num,up/{Latin,latin},\mathup}]
\usepackage{icomma} % virgule décimale sans espace en mode math
% Caractères mathématiques Unicode absents des polices de texte (Plus Jakarta,
% Archivo Black) : rendus par leur équivalent mathématique, en texte comme en
% formule — sinon ils s'affichent en carré vide (ex. « Arithmétique dans ℤ »).
\usepackage{newunicodechar}
\newunicodechar{ℝ}{\ensuremath{\mathbb{R}}}
\newunicodechar{ℤ}{\ensuremath{\mathbb{Z}}}
\newunicodechar{ℂ}{\ensuremath{\mathbb{C}}}
\newunicodechar{ℕ}{\ensuremath{\mathbb{N}}}
\newunicodechar{ℚ}{\ensuremath{\mathbb{Q}}}
\newunicodechar{𝔻}{\ensuremath{\mathbb{D}}}
\newunicodechar{α}{\ensuremath{\alpha}}
\newunicodechar{β}{\ensuremath{\beta}}
\newunicodechar{γ}{\ensuremath{\gamma}}
\newunicodechar{δ}{\ensuremath{\delta}}
\newunicodechar{ε}{\ensuremath{\varepsilon}}
\newunicodechar{θ}{\ensuremath{\theta}}
\newunicodechar{λ}{\ensuremath{\lambda}}
\newunicodechar{μ}{\ensuremath{\mu}}
\newunicodechar{σ}{\ensuremath{\sigma}}
\newunicodechar{τ}{\ensuremath{\tau}}
\newunicodechar{φ}{\ensuremath{\varphi}}
\newunicodechar{ψ}{\ensuremath{\psi}}
\newunicodechar{ω}{\ensuremath{\omega}}
\newunicodechar{Σ}{\ensuremath{\Sigma}}
% majuscules produites par \MakeUppercase dans les titres (α-aminés -> Α)
\newunicodechar{Α}{\ensuremath{\alpha}}
\newunicodechar{Β}{\ensuremath{\beta}}
\newunicodechar{Γ}{\ensuremath{\Gamma}}
\newunicodechar{Δ}{\ensuremath{\Delta}}
\newunicodechar{Ε}{\ensuremath{\varepsilon}}
\newunicodechar{Θ}{\ensuremath{\Theta}}
\newunicodechar{Λ}{\ensuremath{\Lambda}}
\newunicodechar{Μ}{\ensuremath{\mu}}
\newunicodechar{Τ}{\ensuremath{\tau}}
\newunicodechar{Φ}{\ensuremath{\Phi}}
\newunicodechar{Ψ}{\ensuremath{\Psi}}
\newunicodechar{Ω}{\ensuremath{\Omega}}
\newunicodechar{↦}{\ensuremath{\mapsto}}
\newunicodechar{∈}{\ensuremath{\in}}
\newunicodechar{∉}{\ensuremath{\notin}}
\newunicodechar{∧}{\ensuremath{\wedge}}
\newunicodechar{⇔}{\ensuremath{\Leftrightarrow}}
\newunicodechar{⟺}{\ensuremath{\Longleftrightarrow}}
\newunicodechar{⇒}{\ensuremath{\Rightarrow}}
\newunicodechar{⟹}{\ensuremath{\Longrightarrow}}
\newunicodechar{∘}{\ensuremath{\circ}}
\newunicodechar{∪}{\ensuremath{\cup}}
\newunicodechar{∩}{\ensuremath{\cap}}
\newunicodechar{≡}{\ensuremath{\equiv}}
\newunicodechar{⊂}{\ensuremath{\subset}}
\newunicodechar{⊥}{\ensuremath{\perp}}
\newunicodechar{∃}{\ensuremath{\exists}}
\newunicodechar{∀}{\ensuremath{\forall}}
\newunicodechar{∛}{\ensuremath{\sqrt[3]{\,}}}
\newunicodechar{∜}{\ensuremath{\sqrt[4]{\,}}}
\newunicodechar{⃗}{\ensuremath{{}^{\to}}}
\newunicodechar{ⁿ}{\ifmmode^{n}\else\textsuperscript{\ensuremath{n}}\fi}
\newunicodechar{ˣ}{\ifmmode^{x}\else\textsuperscript{\ensuremath{x}}\fi}
\newunicodechar{ᵇ}{\ifmmode^{b}\else\textsuperscript{\ensuremath{b}}\fi}
\newunicodechar{ʳ}{\ifmmode^{r}\else\textsuperscript{\ensuremath{r}}\fi}
\newunicodechar{ᵘ}{\ifmmode^{u}\else\textsuperscript{\ensuremath{u}}\fi}
\newunicodechar{ʸ}{\ifmmode^{y}\else\textsuperscript{\ensuremath{y}}\fi}
\newunicodechar{ᵃ}{\ifmmode^{a}\else\textsuperscript{\ensuremath{a}}\fi}
\newunicodechar{ᵉ}{\ifmmode^{e}\else\textsuperscript{\ensuremath{e}}\fi}
\newunicodechar{ᵏ}{\ifmmode^{k}\else\textsuperscript{\ensuremath{k}}\fi}
\newunicodechar{ᵖ}{\ifmmode^{p}\else\textsuperscript{\ensuremath{p}}\fi}
\newunicodechar{ᴺ}{\ifmmode^{N}\else\textsuperscript{\ensuremath{N}}\fi}
\newunicodechar{ᶜ}{\ifmmode^{c}\else\textsuperscript{\ensuremath{c}}\fi}
\newunicodechar{ʰ}{\ifmmode^{h}\else\textsuperscript{\ensuremath{h}}\fi}
\newunicodechar{ᵠ}{\ifmmode^{q}\else\textsuperscript{\ensuremath{q}}\fi}
\newunicodechar{ₙ}{\ifmmode_{n}\else\textsubscript{\ensuremath{n}}\fi}
\newunicodechar{ₐ}{\ifmmode_{a}\else\textsubscript{\ensuremath{a}}\fi}
\newunicodechar{ᵢ}{\ifmmode_{i}\else\textsubscript{\ensuremath{i}}\fi}
\newunicodechar{ᵣ}{\ifmmode_{r}\else\textsubscript{\ensuremath{r}}\fi}
\newunicodechar{ₖ}{\ifmmode_{k}\else\textsubscript{\ensuremath{k}}\fi}
\newunicodechar{ᵦ}{\ifmmode_{\beta}\else\textsubscript{\ensuremath{\beta}}\fi}
\newunicodechar{ₓ}{\ifmmode_{x}\else\textsubscript{\ensuremath{x}}\fi}
\newfontfamily\popdisplay{ArchivoBlack-Regular.ttf}[Path=../../../fonts/]
\newfontfamily\poplogo{SpaceGrotesk-Bold.ttf}[Path=../../../fonts/]
\newfontfamily\popsemi{PlusJakartaSans-SemiBold.ttf}[Path=../../../fonts/]
\newfontfamily\popxbold{PlusJakartaSans-ExtraBold.ttf}[Path=../../../fonts/]
\newfontfamily\popxboldls{PlusJakartaSans-ExtraBold.ttf}[Path=../../../fonts/,LetterSpace=3.0]

\setlist[enumerate]{leftmargin=1.7em,itemsep=5pt,topsep=4pt}
\setlist[itemize]{leftmargin=1.7em,itemsep=4pt,topsep=4pt,label={\color{poporange}\textbullet}}
\setlength{\parindent}{0pt}
\setlength{\parskip}{5pt}
\renewcommand{\arraystretch}{1.25}

% pgfplots : style Pop par défaut
\pgfplotsset{
  every axis/.append style={axis line style={popink,line width=1pt},tick style={popink},
    grid style={popline,line width=0.5pt},label style={font=\small},tick label style={font=\footnotesize}},
  popcurve/.style={poporange,line width=1.8pt,no markers,smooth},
}
% Même style disponible dans un \draw TikZ simple (hors pgfplots)
\tikzset{popcurve/.style={poporange,line width=1.8pt}}
\tikzset{popgrid/.style={popline,very thin}}
\colorlet{popcurve}{poporange} % le nom du style est parfois utilisé comme couleur

% ============================================================
% Pill / squiggle
% ============================================================
\newcommand{\poppill}[3]{%
  \tikz[baseline=-0.55ex]\node[draw=popink,line width=1.2pt,rounded corners=2.6pt,%
    fill=#2,inner xsep=6.5pt,inner ysep=3pt]{%
    \popxboldls\fontsize{7}{7}\selectfont\color{#3}\MakeUppercase{#1}};%
}
\newcommand{\popsquiggle}[1]{%
  \tikz[baseline=-0.4ex]{
    \draw[#1,line width=2.6pt,line cap=round]
      plot [smooth] coordinates {(0,0.09) (0.34,0.01) (0.68,0.21) (1.02,0.01) (1.36,0.10)};
  }%
}
% Mot-clé surligné (marqueur orange sous le texte)
\newcommand{\cle}[1]{{\popxbold\color{popdark}#1}}

% ============================================================
% En-tête / pied
% ============================================================
\pagestyle{fancy}
\fancyhf{}
\renewcommand{\headrulewidth}{0pt}
\renewcommand{\footrulewidth}{0pt}
\lhead{\raisebox{-0.15ex}{\poplogo\fontsize{13}{13}\selectfont\color{popink}OnBuch\color{poporange}+}}
\rhead{\poppill{\DOCMATIERE\ \textbullet\ \DOCNIVEAU\ \textbullet\ Leçon \DOCLECON}{white}{popink}}
\lfoot{\popxbold\fontsize{7.6}{7.6}\selectfont\color{popmuted}\MakeUppercase{Cours signé OnBuch+}\ \textbullet\ {\popsemi\DOCTITRE}}
\rfoot{\tikz[baseline=-0.6ex]\node[rounded corners=8.5pt,fill=popink,minimum height=17pt,minimum width=17pt,inner xsep=5pt,inner sep=0]{\popxbold\fontsize{8}{8}\selectfont\color{popcream}\thepage\,/\,\pageref*{LastPage}};}

% ============================================================
% Titres
% ============================================================
\newcommand{\chaptitre}[2]{%
  \begin{tcolorbox}[enhanced,colback=poporange,colframe=popink,boxrule=2.6pt,arc=11pt,
      drop shadow=popink,left=15pt,right=15pt,top=12pt,bottom=11pt,before skip=3pt,after skip=18pt]
    \poppill{#2}{popink}{popcream}\par\vspace{9pt}
    {\raggedright\hyphenpenalty=10000\exhyphenpenalty=10000\popdisplay\fontsize{22}{24}\selectfont\color{popcream}\MakeUppercase{#1}\par}
  \end{tcolorbox}
}
% Section numérotée : pastille orange avec le numéro + titre Archivo
\newcounter{popsec}
\newcommand{\coursec}[1]{%
  \stepcounter{popsec}\setcounter{popsub}{0}%
  \par\vspace{16pt}\noindent
  \begin{minipage}{\linewidth}
  \tikz[baseline=-0.7ex]\node[draw=popink,line width=1.6pt,fill=poporange,rounded corners=4pt,minimum size=22pt,inner sep=0]{\popdisplay\fontsize{12}{12}\selectfont\color{popcream}\thepopsec};%
  \hspace{8pt}{\popdisplay\fontsize{15}{17}\selectfont\color{popink}\MakeUppercase{#1}}\par
  \vspace{4pt}\noindent{\color{poporange}\rule{36pt}{3.6pt}}
  \end{minipage}\par\nopagebreak\vspace{8pt}
}
\newcounter{popsub}
\newcommand{\courssub}[1]{%
  \stepcounter{popsub}%
  \par\vspace{9pt}\noindent{\popxbold\fontsize{11.5}{13}\selectfont\color{popink}{\color{poporange}\thepopsec.\thepopsub}\hspace{6pt}#1}\par\nopagebreak\vspace{4pt}
}

% ============================================================
% Boîtes pédagogiques — même recette : fond blanc, bordure épaisse colorée,
% ombre dure encre, badge plein. Titre optionnel [#1].
% ============================================================
\tcbset{popbase/.style={breakable,enhanced,colback=white,boxrule=2.3pt,arc=9pt,
  shadow={3pt}{-3pt}{0pt}{popink},left=14pt,right=14pt,top=11pt,bottom=11pt,before skip=11pt,after skip=15pt,
  fontupper=\popsemi\fontsize{10.6}{14.5}\selectfont\color{popink},
  fontlower=\popsemi\fontsize{10.6}{14.5}\selectfont\color{popink}}}
% Coche dessinée (le glyphe ✓ n'existe pas dans les polices du document)
\newcommand{\popcheck}[1]{\tikz[baseline=-0.5ex]\draw[#1,line width=1.6pt,line cap=round,line join=round](-0.13,0.0)--(-0.04,-0.09)--(0.13,0.1);}
\providecommand{\coche}{\popcheck{popgreen}}
\providecommand{\resultat}[1]{\par\smallskip\noindent\fcolorbox{popgreen}{popgreenL}{\parbox{\dimexpr\linewidth-2\fboxsep-2\fboxrule\relax}{\textbf{Résultat.} #1}}\par\smallskip}
\newcommand{\popboxhead}[3]{\poppill{#1}{#2}{white}\ifx\relax#3\relax\else\hspace{6pt}{\popxbold\fontsize{10.6}{12}\selectfont\color{popink}#3}\fi\par\vspace{7pt}}

\newtcolorbox{aretenir}[1][]{popbase,colframe=popgreen,before upper={\popboxhead{À retenir}{popgreen}{#1}}}
\newtcolorbox{exemplebox}[1][]{popbase,colframe=poporange,before upper={\popboxhead{Exemple}{poporange}{#1}}}
\newtcolorbox{definition}[1][]{popbase,colframe=popblue,before upper={\popboxhead{Définition}{popblue}{#1}}}
\newtcolorbox{propriete}[1][]{popbase,colframe=poppurple,before upper={\popboxhead{Propriété}{poppurple}{#1}}}
\newtcolorbox{methode}[1][]{popbase,colframe=popgold,before upper={\popboxhead{Méthode}{popgold}{#1}}}
\newtcolorbox{attention}[1][]{popbase,colframe=poppink,before upper={\popboxhead{Attention}{poppink}{#1}}}
\newtcolorbox{experience}[1][]{popbase,colframe=popblue,colback=popblueL,before upper={\popboxhead{Expérience}{popblue}{#1}}}
% « Le savais-tu ? » : carte violette pleine (contexte, culture, Cameroun)
\newtcolorbox{savaistu}[1][]{popbase,colback=poppurple,colframe=popink,
  fontupper=\popsemi\fontsize{10.4}{14.2}\selectfont\color{white},
  before upper={\poppill{Le savais-tu\,?}{white}{poppurple}\ifx\relax#1\relax\else\hspace{6pt}{\popxbold\color{white}#1}\fi\par\vspace{7pt}}}
% Exercice résolu : énoncé en haut, \tcblower puis la solution sur fond crème
\newcounter{popexo}\newcounter{popexr}
\newtcolorbox{exoresolu}[1][]{popbase,colframe=poporange,colbacklower=popcream,
  segmentation style={popink,line width=1.2pt,dashed},
  before upper={\stepcounter{popexr}\popboxhead{Exercice résolu \thepopexr}{poporange}{#1}},
  before lower={\poppill{Solution}{popink}{popcream}\par\vspace{6pt}}}
% Exercice (non résolu en place) — la correction est dans la partie Corrigés
\newtcolorbox{exercice}[1][]{popbase,colframe=popink,
  before upper={\stepcounter{popexo}\popboxhead{Exercice \thepopexo}{popink}{#1}}}
% Corrigé
\newtcolorbox{corrige}[1][]{popbase,colframe=popgreen,colback=popgreenL,
  before upper={\popboxhead{Corrigé}{popgreen}{#1}}}
% Objectifs / prérequis / bilan
\newtcolorbox{objectifs}{popbase,colframe=popink,colback=poporangeL,
  before upper={\popboxhead{Objectifs de la leçon}{popink}{}}}
\newtcolorbox{prerequis}{popbase,colframe=popink,colback=popblueL,
  before upper={\popboxhead{Prérequis}{popblue}{}}}
\newtcolorbox{bilan}{popbase,colframe=popink,colback=white,boxrule=2.8pt,
  before upper={\popboxhead{Fiche bilan}{poporange}{L'essentiel à connaître pour l'examen}}}
% Figure encadrée avec légende
\newtcolorbox{popfigure}[1][]{enhanced,breakable=false,colback=white,colframe=popline,boxrule=1.4pt,arc=8pt,
  left=10pt,right=10pt,top=10pt,bottom=8pt,before skip=10pt,after skip=12pt,halign=center,
  lower separated=false,halign lower=center,
  fontlower=\popsemi\fontsize{9}{11.5}\selectfont\color{popmuted},
  #1}
\newcommand{\legende}[1]{\par\vspace{6pt}{\popsemi\fontsize{9}{11.5}\selectfont\color{popmuted}#1}}

% ============================================================
% Signature OnBuch+
% ============================================================
\newcommand{\poplogoinline}[1]{{\poplogo\fontsize{#1}{#1}\selectfont\color{popink}OnBuch\color{poporange}+}}
\newcommand{\signatureonbuch}{%
  \begin{tcolorbox}[enhanced,colback=popink,colframe=popink,boxrule=2.2pt,arc=10pt,
      drop shadow=poporange,left=14pt,right=14pt,top=10pt,bottom=10pt,before skip=0pt,after skip=0pt]
    \tikz[baseline=-0.7ex]\node[circle,fill=popgreen,draw=popcream,line width=1.3pt,minimum size=22pt,inner sep=0]{\popcheck{white}};%
    \hspace{9pt}\begin{minipage}[c]{0.62\linewidth}
      {\popxbold\fontsize{12}{13}\selectfont\color{popcream}Signé\ \textbullet\ {\poplogo OnBuch\color{poporange}+}}\par\vspace{2pt}
      {\raggedright\popsemi\fontsize{8}{10.5}\selectfont\color{popline}\MakeUppercase{Équipe pédagogique OnBuch+}\\\MakeUppercase{Conforme au programme officiel MINESEC}\par}
    \end{minipage}\hfill
    {\poplogo\fontsize{22}{22}\selectfont\color{popcream}OnBuch\color{poporange}+}
  \end{tcolorbox}%
}

% ============================================================
% Page de garde
% #1 titre · #2 matière · #3 niveau · #4 module · #5 n° leçon · #6 durée ·
% #7 description / objectifs (texte ou itemize)
% ============================================================
\newcommand{\pagegarde}[7]{%
  \thispagestyle{empty}
  \newgeometry{a4paper,margin=1.7cm,top=1.6cm,bottom=1.4cm}%
  \begin{tikzpicture}[remember picture,overlay]
    \fill[poporange!18] ([xshift=-3.2cm,yshift=-3.6cm]current page.north east) circle (4.6cm);
    \fill[poppurple!14] ([xshift=2.4cm,yshift=7.6cm]current page.south west) circle (3.8cm);
  \end{tikzpicture}%
  {\poplogo\fontsize{30}{30}\selectfont\color{popink}OnBuch\color{poporange}+}\hfill
  \raisebox{0.6ex}{\poppill{Cours signé \textbullet\ Premium}{popink}{popcream}}\par
  \vspace{1pt}\popsquiggle{poppurple}\par
  \vspace{8pt}
  \begin{tcolorbox}[enhanced,colback=poporange,colframe=popink,boxrule=2.8pt,arc=13pt,
      drop shadow=popink,left=18pt,right=18pt,top=12pt,bottom=13pt,after skip=10pt]
    \poppill{#2 \textbullet\ #3}{popink}{popcream}\hspace{5pt}\poppill{#4}{white}{popink}\par\vspace{8pt}
    {\popdisplay\fontsize{14}{14}\selectfont\color{popink}LEÇON #5}\par\vspace{4pt}
    {\raggedright\hyphenpenalty=10000\exhyphenpenalty=10000\popdisplay\fontsize{25}{27}\selectfont\color{popcream}\MakeUppercase{#1}\par}
    \vspace{8pt}
    {\popxbold\fontsize{9.5}{9.5}\selectfont\color{popink}\MakeUppercase{Durée conseillée : #6}}
  \end{tcolorbox}
  \begin{tcolorbox}[enhanced,colback=white,colframe=popink,boxrule=2.2pt,arc=10pt,
      drop shadow=popink,left=16pt,right=16pt,top=10pt,bottom=10pt,after skip=8pt]
    \poppill{Dans cette leçon}{poppurple}{white}\par\vspace{5pt}
    \popsemi\fontsize{9.3}{12.6}\selectfont\color{popink}#7
  \end{tcolorbox}
  \noindent\begin{minipage}[t]{0.487\linewidth}
    \begin{tcolorbox}[enhanced,colback=popgreen,colframe=popink,boxrule=2.2pt,arc=10pt,
        drop shadow=popink,left=11pt,right=11pt,top=10pt,bottom=10pt,height=3.1cm,valign=center]
      \tcbox[colback=white,colframe=popink,boxrule=1.3pt,arc=5pt,boxsep=0pt,nobeforeafter,
        left=3pt,right=3pt,top=3pt,bottom=3pt,valign=center]{\qrcode[height=1.9cm]{https://play.google.com/store/apps/details?id=com.luvvix.onbuch&hl=fr}}%
      \hspace{8pt}\begin{minipage}[c]{0.5\linewidth}
        {\popdisplay\fontsize{11}{12.5}\selectfont\color{white}\MakeUppercase{Télécharge OnBuch}}\par\vspace{4pt}
        {\raggedright\popsemi\fontsize{8.4}{11}\selectfont\color{white}Gratuit \textbullet\ Google Play\\Tuteur IA, annales, concours, résultats\par}
      \end{minipage}
    \end{tcolorbox}
  \end{minipage}\hfill
  \begin{minipage}[t]{0.487\linewidth}
    \begin{tcolorbox}[enhanced,colback=poppurple,colframe=popink,boxrule=2.2pt,arc=10pt,
        drop shadow=popink,left=11pt,right=11pt,top=10pt,bottom=10pt,height=3.1cm,valign=center]
      \tikz[baseline=-0.7ex]\node[circle,fill=white,draw=popink,line width=1.3pt,minimum size=1.6cm,inner sep=0]{\popdisplay\fontsize{24}{24}\selectfont\color{poppurple}+};%
      \hspace{8pt}\begin{minipage}[c]{0.6\linewidth}
        {\popdisplay\fontsize{11}{12.5}\selectfont\color{white}\MakeUppercase{Passe à OnBuch+}}\par\vspace{4pt}
        {\raggedright\popsemi\fontsize{8.4}{11}\selectfont\color{white}Tous les cours signés, Tuteur IA illimité, fiches PDF — onglet OnBuch+ de l'app\par}
      \end{minipage}
    \end{tcolorbox}
  \end{minipage}
  \vspace{8pt}
  \popcontenu
  \vfill
  \signatureonbuch
  \restoregeometry
  \newpage
}

% Rangée « Ce cours contient » (page de garde)
\newcommand{\poptile}[3]{% #1 couleur #2 gros texte #3 légende
  \begin{tcolorbox}[enhanced,colback=#1,colframe=popink,boxrule=2pt,arc=9pt,shadow={2.5pt}{-2.5pt}{0pt}{popink},
      left=6pt,right=6pt,top=9pt,bottom=9pt,halign=center,valign=center,height=1.75cm,width=\linewidth,nobeforeafter]
    {\popdisplay\fontsize{12.5}{13}\selectfont\color{white}\MakeUppercase{#2}}\par\vspace{3pt}
    {\centering\popsemi\fontsize{7.8}{9.5}\selectfont\color{white}#3\par}
  \end{tcolorbox}}
\newcommand{\popcontenu}{%
  \par\noindent{\popxboldls\fontsize{7.5}{7.5}\selectfont\color{popmuted}CE COURS SIGNÉ CONTIENT}\par\vspace{7pt}
  \noindent\begin{minipage}[t]{0.235\linewidth}\poptile{popblue}{Cours}{complet, illustré, pas à pas}\end{minipage}\hfill
  \begin{minipage}[t]{0.235\linewidth}\poptile{poporange}{Méthodes}{et exercices résolus}\end{minipage}\hfill
  \begin{minipage}[t]{0.235\linewidth}\poptile{poppink}{Exercices}{type examen + corrigés}\end{minipage}\hfill
  \begin{minipage}[t]{0.235\linewidth}\poptile{popgreen}{Fiche bilan}{l'essentiel à retenir}\end{minipage}\par}

% Sommaire
\newtcolorbox{sommaire}{enhanced,colback=white,colframe=popink,boxrule=2.2pt,arc=10pt,
  drop shadow=popink,left=18pt,right=18pt,top=13pt,bottom=13pt,before skip=6pt,after skip=20pt,
  before upper={\poppill{Sommaire}{poppurple}{white}\par\vspace{9pt}\popsemi\fontsize{10.6}{14.5}\selectfont\color{popink}}}

% Signature de fin de cours — poussée tout en bas de la dernière page
\newcommand{\findecours}{%
  \par\vfill\nopagebreak
  \begin{center}\popsquiggle{poporange}\end{center}\vspace{4pt}
  \signatureonbuch
}
\AtBeginDocument{\def\llbracket{\mathopen{[\mkern-3.5mu[}}\def\rrbracket{\mathclose{]\mkern-3.5mu]}}} % intervalles d'entiers (glyphe ⟦ absent de la police)
% Glyphes absents de Fira Math : redéfinis à partir de glyphes disponibles
\AtBeginDocument{%
  \def\setminus{\mathbin{\backslash}}\let\smallsetminus\setminus
  \def\vdots{\mathord{\vbox{\baselineskip3.2pt\lineskiplimit0pt\kern4pt\hbox{.}\hbox{.}\hbox{.}}}}%
  \def\ddots{\mathinner{\mkern1mu\raise6.5pt\hbox{.}\mkern2mu\raise3.6pt\hbox{.}\mkern2mu\raise0.7pt\hbox{.}\mkern1mu}}%
  \def\triangle{\mathord{\scalebox{0.9}{$\symup{\Delta}$}}}%
  \def\nabla{\mathord{\raisebox{\depth}{\scalebox{1}[-1]{$\symup{\Delta}$}}}}%
  \def\checkmark{\popcheck{popdark}}%
  \def\star{\mathbin{\ast}}\def\bigstar{\mathord{\ast}}%
  \def\diamond{\mathbin{\scalebox{0.6}{$\Diamond$}}}%
  \def\bigcup{\mathop{\mathchoice{\raisebox{-0.3em}{\scalebox{1.7}{$\cup$}}}{\scalebox{1.25}{$\cup$}}{\cup}{\cup}}\displaylimits}%
  \def\bigcap{\mathop{\mathchoice{\raisebox{-0.3em}{\scalebox{1.7}{$\cap$}}}{\scalebox{1.25}{$\cap$}}{\cap}{\cap}}\displaylimits}%
  \def\lesseqqgtr{\mathrel{\lessgtr}}\def\bigoplus{\mathop{\oplus}\displaylimits}%
}
\providecommand{\ssi}{si et seulement si} % abréviation fréquente
\AtBeginDocument{\providecommand{\wideparen}{\overparen}} % arc de cercle

%% FIGFIX-BEGIN v5 — correctifs globaux des figures (docs/onbuch-generation/tools/figfix)
\usepackage{adjustbox}\usepackage{etoolbox}\tcbuselibrary{raster}
% toute figure TikZ/pgfplots plus large que la ligne est réduite à la largeur disponible
\BeforeBeginEnvironment{tikzpicture}{\begin{adjustbox}{max width=\linewidth}}
\AfterEndEnvironment{tikzpicture}{\end{adjustbox}}
% légendes pgfplots lisibles par-dessus les courbes
\pgfplotsset{every axis/.append style={legend style={fill=white,fill opacity=0.92,text opacity=1,draw=black!15,inner sep=3pt}}}
% étiquettes d'arêtes (syntaxe edge["..."]) avec fond blanc
\tikzset{every edge quotes/.append style={fill=white,inner sep=1.2pt}}
%% FIGFIX-END
\begin{document}

\pagegarde{Réduction des risques liés aux catastrophes naturelles}{SVT}{Tle D}{Module III : L'Éducation à l'Environnement et au Développement Durable — Les mouvements de la lithosphère et leurs conséquences}{N15}{10 h}{Cette leçon décrit la morphologie des fonds océaniques et le fonctionnement des rifts, puis explique les mouvements des plaques lithosphériques (divergence, convergence, coulissage) et les catastrophes naturelles qui en résultent. Elle se termine par les stratégies de gestion des risques : prévision, prédiction et prévention, appliquées au contexte camerounais.
\vspace{4pt}
\begin{multicols}{2}\small
\begin{itemize}
\item À la fin de cette leçon, je sais analyser des données topographiques (profils bathymétriques) des fonds océaniques et y reconnaître dorsales, rifts, fosses et plaines abyssales
\item À la fin de cette leçon, je sais reconstituer les étapes de la formation d'un rift et du plancher océanique (expansion océanique)
\item À la fin de cette leçon, je sais réaliser une maquette simple du rift et du plancher océanique
\item À la fin de cette leçon, je sais exploiter les données expérimentales de fusion des péridotites pour expliquer la naissance du magma en zone d'accrétion et en zone de subduction
\item À la fin de cette leçon, je sais décrire les mouvements de convergence (subduction, collision) et de coulissage des plaques et identifier leurs facteurs clés
\item À la fin de cette leçon, je sais associer chaque type de catastrophe naturelle (séisme, tsunami, volcanisme, glissement de terrain, inondation) à un contexte géodynamique
\end{itemize}\end{multicols}}

\begin{sommaire}
\begin{multicols}{2}
\begin{enumerate}
\item Morphologie des fonds océaniques
\item Le rift et l'expansion du plancher océanique
\item Naissance des magmas : données expérimentales de fusion des péridotites
\item Convergence et coulissage des plaques
\item Les catastrophes naturelles liées aux mouvements des plaques
\item Gestion des catastrophes naturelles : prévision, prédiction, prévention
\item Activité d'intégration
\item Méthodes et exercices résolus
\item Exercices
\item Corrigés des exercices
\item Fiche bilan
\end{enumerate}
\end{multicols}
\end{sommaire}

\begin{objectifs}
\begin{itemize}
\item À la fin de cette leçon, je sais analyser des données topographiques (profils bathymétriques) des fonds océaniques et y reconnaître dorsales, rifts, fosses et plaines abyssales
\item À la fin de cette leçon, je sais reconstituer les étapes de la formation d'un rift et du plancher océanique (expansion océanique)
\item À la fin de cette leçon, je sais réaliser une maquette simple du rift et du plancher océanique
\item À la fin de cette leçon, je sais exploiter les données expérimentales de fusion des péridotites pour expliquer la naissance du magma en zone d'accrétion et en zone de subduction
\item À la fin de cette leçon, je sais décrire les mouvements de convergence (subduction, collision) et de coulissage des plaques et identifier leurs facteurs clés
\item À la fin de cette leçon, je sais associer chaque type de catastrophe naturelle (séisme, tsunami, volcanisme, glissement de terrain, inondation) à un contexte géodynamique
\item À la fin de cette leçon, je sais distinguer prévision, prédiction et prévention dans la gestion des catastrophes naturelles
\item À la fin de cette leçon, je sais élaborer un outil de sensibilisation (affiche, dépliant, sketch) sur la gestion d'une catastrophe naturelle
\end{itemize}
\end{objectifs}

\begin{prerequis}
\begin{itemize}
\item Structure interne du globe terrestre (croûte, manteau, noyau) vue en classe de Première
\item Notion de lithosphère et d'asthénosphère (rigidité mécanique des enveloppes terrestres)
\item Notions de pression et de température augmentant avec la profondeur (géotherme)
\item Lecture et interprétation de graphiques et de courbes expérimentales
\item Repérage des grandes masses continentales et océaniques sur une carte du monde
\end{itemize}
\end{prerequis}

\newpage

\coursec{Morphologie des fonds océaniques}

\textbf{Situation d'accroche.} Le 22 mai 2019, un glissement de terrain à Bafoussam, dans la région de l'Ouest, a enseveli plusieurs habitations et fait de nombreuses victimes. Chaque année, les inondations de Yaoundé et de Douala emportent des vies et détruisent des quartiers entiers. Le mont Cameroun, volcan actif qui domine la ville de Buéa, est entré en éruption en 1999 puis en 2000, couvrant les plantations de laves. Ailleurs dans le monde, le séisme suivi d'un tsunami qui a frappé le Japon en 2011 a fait près de \num{20000} morts. Pourquoi certaines régions du globe — et du Cameroun — sont-elles plus exposées que d'autres aux séismes, tsunamis et éruptions ? Que se passe-t-il sous les océans et sous les continents pour produire de telles catastrophes ? Et surtout, peut-on les \cle{prévoir}, les \cle{prédire} ou s'en \cle{prémunir} ? Pour répondre à ces questions, il faut d'abord connaître le relief caché sous les océans : c'est là, comme nous le verrons, que se joue l'essentiel de la dynamique de la Terre. Cette leçon part donc de l'exploration des fonds océaniques pour aboutir, à la fin, à la gestion des risques naturels.

\courssub{Des océans immenses et longtemps inconnus}

Les océans couvrent environ \SI{71}{\percent} de la surface du globe, soit près de \SI{361}{millions de km^2}, avec une profondeur moyenne voisine de \SI{3800}{m}. Pourtant, jusqu'au milieu du \textsc{xx}\textsuperscript{e} siècle, l'homme connaissait mieux la surface de la Lune que le fond des mers : la lumière ne pénètre pas au-delà de quelques centaines de mètres, et la pression écrasante interdit toute exploration directe facile. Comment alors cartographier ce monde invisible ? La réponse est venue du son : c'est la \cle{bathymétrie}.

\begin{definition}[Bathymétrie]
La \cle{bathymétrie} (du grec \emph{bathus}, profond, et \emph{metron}, mesure) est la science qui mesure les profondeurs des océans et permet d'établir la topographie (le relief) des fonds marins. Elle repose aujourd'hui sur les \cle{sondeurs acoustiques} : un navire émet une onde sonore vers le fond et mesure le temps que met l'écho à revenir. Les \cle{écho-sondeurs multifaisceaux} modernes balayent une large bande sous le navire et fournissent des cartes en trois dimensions d'une grande précision.
\end{definition}

\courssub{Le principe de la mesure bathymétrique}

L'idée est simple et élégante : le son se propage dans l'eau de mer à une vitesse connue, environ $v = \SI{1500}{m.s^{-1}}$. Si l'on connaît le temps $t$ mis par l'onde pour faire l'aller-retour entre le navire et le fond, on en déduit la profondeur $d$. Attention : le temps mesuré correspond à \emph{deux} trajets (descente et remontée), d'où la division par 2 :

\[ d = \frac{v \times t}{2} \]

\begin{experience}[Mesurer une profondeur à l'écho-sondeur]
\textbf{Protocole.} Un navire océanographique équipé d'un sondeur émet verticalement une impulsion ultrasonore. Un récepteur enregistre l'écho réfléchi par le fond. On note le temps d'aller-retour $t$.\par
\textbf{Observations.} Au large du golfe de Guinée, l'écho revient au bout de $t = \SI{6,0}{s}$ ; au-dessus de la dorsale médio-atlantique, il revient au bout de $t = \SI{3,3}{s}$ seulement.\par
\textbf{Interprétation.} Dans le premier cas : $d = \dfrac{1500 \times 6{,}0}{2} = \SI{4500}{m}$ ; dans le second : $d = \dfrac{1500 \times 3{,}3}{2} = \SI{2475}{m}$. Le fond est donc beaucoup plus \emph{élevé} au niveau de la dorsale : c'est un relief sous-marin.\par
\textbf{Conclusion.} En enchaînant les mesures le long d'une traversée, on obtient un \cle{profil bathymétrique} : une coupe du relief océanique.
\end{experience}

\begin{exoresolu}[Calcul de profondeur]
Un écho-sondeur émet une onde acoustique dont l'écho est reçu au bout de $t = \SI{8,0}{s}$. Sachant que la vitesse du son dans l'eau de mer est $v = \SI{1500}{m.s^{-1}}$, calculer la profondeur du fond à cet endroit. À quelle structure du profil océanique cette valeur peut-elle correspondre ?
\tcblower
\textbf{Solution.} L'onde parcourt deux fois la profondeur $d$ (aller puis retour), donc :
\begin{align*}
2d &= v \times t \\
d &= \frac{v \times t}{2} = \frac{1500 \times 8{,}0}{2} = \SI{4000}{m}
\end{align*}
Une profondeur de \SI{4000}{m} correspond typiquement à une \cle{plaine abyssale}. (Piège classique : oublier de diviser par 2 donnerait \SI{8000}{m}, valeur absurde ici.)
\end{exoresolu}

\courssub{Analyse d'un profil topographique transocéanique}

Grâce à des milliers de profils bathymétriques, les océanographes ont dressé la carte des fonds. Prenons l'exemple d'une traversée de l'océan Atlantique, de l'Afrique vers l'Amérique du Sud. Loin d'être une cuvette plate et monotone, le fond présente un relief \emph{organisé}.

\begin{methode}[Analyser un profil bathymétrique]
\begin{enumerate}
\item Repérer la zone la plus \textbf{élevée} du profil (profondeur la plus faible) : c'est l'axe de la \cle{dorsale}.
\item Chercher au sommet de cette dorsale une \textbf{vallée axiale} étroite et encaissée : c'est le \cle{rift}.
\item Observer de part et d'autre la \textbf{symétrie} des reliefs : collines abyssales puis plaines abyssales, à peu près à la même profondeur à égale distance de l'axe.
\item Repérer les zones les plus \textbf{profondes} du profil (fossés étroits dépassant parfois \SI{8000}{m}) : ce sont les \cle{fosses océaniques}, situées en bordure des océans.
\item Vérifier la présence de décalages de l'axe de la dorsale : ils traduisent les \cle{failles transformantes}.
\end{enumerate}
\end{methode}

\begin{popfigure}
\begin{center}
\begin{tikzpicture}
\begin{axis}[
width=\linewidth, height=7.2cm,
xlabel={Distance le long du profil (km)},
ylabel={Profondeur (m)},
xmin=0, xmax=6000, ymin=-8000, ymax=0,
xtick={0,1000,2000,3000,4000,5000,6000},
ytick={0,-2000,-4000,-6000,-8000},
grid=both, grid style={popline!40},
axis line style={popdark},
label style={popdark},
]
\addplot[popcurve, domain=0:6000, samples=300, smooth]
coordinates {
(0,-500) (200,-1500) (500,-3500) (800,-4800) (1200,-5200) (1600,-5000)
(2000,-4600) (2300,-4200) (2500,-3500) (2700,-2800) (2900,-2500)
(3000,-3400) (3100,-2500) (3300,-2800) (3500,-3500) (3800,-4300)
(4200,-4900) (4600,-5200) (5000,-5100) (5400,-6800) (5600,-7400)
(5800,-5000) (6000,-800)
};
\addplot[popblue!15, domain=0:6000] coordinates {(0,-8000) (6000,-8000)} \closedcycle;
\node[popdark, font=\small\bfseries] at (axis cs:3000,-1400) {Dorsale médio-atlantique};
\node[poppink, font=\small\bfseries] at (axis cs:3000,-4200) {Rift axial};
\draw[->, poppink, thick] (axis cs:3000,-3900) -- (axis cs:3000,-3350);
\node[popblue, font=\small] at (axis cs:1200,-5800) {Plaine abyssale};
\node[popblue, font=\small] at (axis cs:4600,-5800) {Plaine abyssale};
\node[popgreen, font=\small] at (axis cs:2200,-3300) {Collines abyssales};
\node[popgreen, font=\small] at (axis cs:3900,-3300) {Collines abyssales};
\node[poporange, font=\small\bfseries] at (axis cs:5550,-5600) {Fosse};
\draw[->, poporange, thick] (axis cs:5560,-5900) -- (axis cs:5600,-7100);
\end{axis}
\end{tikzpicture}
\end{center}
\legende{Profil bathymétrique transatlantique schématique (Afrique à gauche, Amérique du Sud à droite). La dorsale médio-atlantique s'élève vers \SI{-2500}{m} et est creusée en son axe par le rift ; de part et d'autre s'alignent les collines abyssales puis les plaines abyssales (vers \SI{-5000}{m}), de façon remarquablement symétrique. Une fosse profonde marque la bordure de l'océan.}
\end{popfigure}

\courssub{Les grandes structures des fonds océaniques}

\textbf{La dorsale médio-océanique et le rift.} Au milieu de l'Atlantique court un immense relief sous-marin, large de \SI{1000}{km} à \SI{2000}{km} et long de plusieurs dizaines de milliers de kilomètres : la \cle{dorsale médio-atlantique}. Son sommet se situe vers \SI{2500}{m} de profondeur, soit environ \SI{2500}{m} \emph{au-dessus} des plaines voisines. Fait capital : son axe est creusé par une vallée étroite (quelques dizaines de kilomètres de large) et profonde de \SI{1000}{m} à \SI{2000}{m} : le \cle{rift}.

\begin{definition}[Dorsale médio-océanique et rift]
Une \cle{dorsale médio-océanique} est un relief sous-marin élevé, allongé et continu (l'ensemble des dorsales forme un réseau d'environ \SI{60000}{km} à travers tous les océans). Son axe est creusé par une vallée d'effondrement appelée \cle{rift}, siège d'une activité volcanique et sismique intense. C'est au niveau des dorsales que se forme le nouveau plancher océanique (zone d'\cle{accrétion}).
\end{definition}

\textbf{Les plaines et les collines abyssales.} De part et d'autre de la dorsale, le fond s'enfonce progressivement : on rencontre d'abord les \cle{collines abyssales}, relief vallonné de quelques centaines de mètres, puis les \cle{plaines abyssales}, surfaces très planes situées entre \SI{4000}{m} et \SI{6000}{m} de profondeur. Leur planéité s'explique par l'accumulation de sédiments fins qui comblent les irrégularités du basalte sous-jacent.

\begin{definition}[Plaine abyssale]
Une \cle{plaine abyssale} est une vaste étendue plane des grands fonds océaniques, située entre \SI{4000}{m} et \SI{6000}{m} de profondeur, recouverte de sédiments. Elle constitue la partie la plus ancienne et la plus profonde du plancher océanique, loin de la dorsale.
\end{definition}

\textbf{Les fosses océaniques.} En bordure de certains océans (surtout le Pacifique), le fond plonge brutalement dans des fossés étroits et très profonds : les \cle{fosses océaniques}. La plus profonde, la fosse des Mariannes, atteint \SI{-11034}{m} : on pourrait y enfoncer l'Everest (\SI{8849}{m}) et il resterait encore plus de \SI{2000}{m} d'eau au-dessus ! Les fosses sont bordées soit d'\cle{arcs insulaires} volcaniques (Japon, Aléoutiennes), soit de \cle{chaînes de montagnes continentales} (les Andes, le long de la fosse du Pérou-Chili).

\begin{definition}[Fosse océanique]
Une \cle{fosse océanique} est une dépression étroite et très profonde (de \SI{6000}{m} à plus de \SI{11000}{m}) située en bordure d'océan. Elle marque une zone d'\cle{enfoncement} (subduction) où le plancher océanique plonge dans le manteau ; elle s'accompagne d'une forte activité sismique et volcanique.
\end{definition}

\textbf{Les failles transformantes.} L'axe des dorsales n'est pas continu : il est découpé en segments décalés les uns par rapport aux autres par de grandes fractures verticales, les \cle{failles transformantes}. Le long de ces failles, deux portions de plancher coulissent l'une contre l'autre, ce qui provoque des séismes. En Atlantique équatorial, la faille Romanche décale la dorsale de plus de \SI{600}{km}.

\begin{popfigure}
\begin{center}
\begin{tikzpicture}[scale=0.92, every node/.style={font=\small}]
% Coupe de la dorsale
\begin{scope}
\fill[popblue!12] (-5.2,-2.6) rectangle (5.2,0);
\draw[popdark, very thick] (-5.2,0) -- (-2.4,-0.5) -- (-1.2,-0.9) -- (-0.5,-0.6) -- (-0.35,-1.7) -- (0.35,-1.7) -- (0.5,-0.6) -- (1.2,-0.9) -- (2.4,-0.5) -- (5.2,0);
\fill[poporange!25] (-0.35,-1.7) -- (0.35,-1.7) -- (0.5,-2.6) -- (-0.5,-2.6) -- cycle;
\draw[->, poppink, very thick] (-0.9,-2.2) -- (-2.2,-2.2);
\draw[->, poppink, very thick] (0.9,-2.2) -- (2.2,-2.2);
\node[popdark, font=\small\bfseries] at (-3.3,0.35) {Dorsale};
\node[poppink, font=\small\bfseries] at (0,-2.05) {Rift};
\node[popblue, font=\small] at (4.0,-0.75) {Collines};
\node[popblue, font=\small] at (-4.2,-0.75) {Plaine};
\node[popmuted, font=\footnotesize] at (0,-2.9) {Coupe transversale d'une dorsale (le rift est la vallée axiale)};
\end{scope}
% Vue en plan : failles transformantes
\begin{scope}[yshift=-5.6cm]
\foreach \y/\xa/\xb in {0.9/-4.6/-1.4, 0.3/-1.0/2.2, -0.3/2.6/4.6}{
\draw[popdark, very thick] (\xa,\y) -- (\xb,\y);
\fill[poppink] (\xa,\y) circle (1.6pt);
\fill[poppink] (\xb,\y) circle (1.6pt);
}
\draw[poporange, very thick, dashed] (-1.4,0.9) -- (-1.0,0.3);
\draw[poporange, very thick, dashed] (2.2,0.3) -- (2.6,-0.3);
\draw[->, popgreen, thick] (-1.2,1.15) -- (-2.4,1.15);
\draw[->, popgreen, thick] (-1.2,0.65) -- (0.0,0.65);
\draw[->, popgreen, thick] (2.4,0.55) -- (1.2,0.55);
\draw[->, popgreen, thick] (2.4,0.05) -- (3.6,0.05);
\node[poporange, font=\small\bfseries] at (-3.4,-0.55) {Failles transformantes};
\draw[->, poporange] (-2.2,-0.45) -- (-1.25,0.55);
\node[popdark, font=\small] at (3.6,1.15) {Segments de dorsale};
\node[popmuted, font=\footnotesize] at (0,-0.95) {Vue en plan : l'axe de la dorsale est décalé par les failles transformantes};
\end{scope}
\end{tikzpicture}
\end{center}
\legende{La dorsale et ses structures. \textbf{Haut} : coupe transversale montrant le relief élevé de la dorsale, la vallée axiale du rift (en orange) et l'écartement du plancher de part et d'autre (flèches roses). \textbf{Bas} : vue en plan de l'axe de la dorsale (traits noirs) découpé en segments décalés par les failles transformantes (pointillés orange) ; les flèches vertes indiquent le sens de coulissage des deux côtés de la faille.}
\end{popfigure}

\courssub{Bilan : une topographie organisée et symétrique}

\begin{aretenir}[L'essentiel sur la morphologie des fonds océaniques]
\begin{itemize}
\item Les fonds océaniques sont cartographiés par \cle{bathymétrie} acoustique : $d = \dfrac{v \times t}{2}$.
\item La \cle{dorsale médio-océanique} est un relief élevé (vers \SI{-2500}{m}) dont l'axe est creusé par le \cle{rift}.
\item De part et d'autre : \cle{collines abyssales} puis \cle{plaines abyssales} (\SI{4000}{m} à \SI{6000}{m}), de façon \textbf{symétrique} par rapport à l'axe de la dorsale.
\item Les \cle{fosses océaniques} (jusqu'à \SI{-11034}{m} dans la fosse des Mariannes) bordent certains océans, associées à des arcs insulaires ou à des chaînes de montagnes.
\item Les \cle{failles transformantes} décalent l'axe des dorsales.
\item Cette organisation symétrique n'est pas un hasard : elle témoigne de la formation du plancher océanique au niveau des dorsales, puis de son écartement progressif — ce sera l'objet de la section suivante.
\end{itemize}
\end{aretenir}

\begin{exemplebox}[Quelques chiffres de référence]
La dorsale médio-atlantique culmine vers \SI{-2500}{m} alors que les plaines abyssales voisines s'étendent vers \SI{-5000}{m} : la dorsale constitue donc un relief de près de \SI{2500}{m} de hauteur, comparable à une grande chaîne de montagnes terrestre. À l'opposé, la fosse des Mariannes plonge à \SI{-11034}{m}. Entre ces deux extrêmes, tout le relief océanique s'organise autour de l'axe des dorsales.
\end{exemplebox}

\begin{attention}[Ne pas confondre dorsale et fosse]
Erreur fréquente au Baccalauréat : confondre \cle{dorsale} et \cle{fosse}. Ce sont deux structures \textbf{opposées} : la dorsale est un relief \emph{élevé}, zone d'\emph{expansion} où naît le plancher océanique (accrétion) ; la fosse est une dépression \emph{profonde}, zone d'\emph{enfoncement} où le plancher disparaît (subduction). Sur un profil, la dorsale se reconnaît à sa faible profondeur et à son rift axial ; la fosse, à sa grande profondeur et à sa position en bordure d'océan.
\end{attention}

\begin{savaistu}[Le mont Cameroun et la Ligne volcanique du Cameroun]
Avec ses \SI{4100}{m}, le mont Cameroun est le point culminant d'Afrique de l'Ouest. Il appartient à la \cle{Ligne volcanique du Cameroun}, un alignement de volcans de plus de \SI{1500}{km} qui traverse le pays du Sud-Ouest vers l'Adamaoua et se prolonge dans le golfe de Guinée par les îles de Bioko, São Tomé et Príncipe. Cette ligne volcanique rappelle que le Cameroun, lui aussi, est concerné par la dynamique interne de la Terre : les éruptions de 1999 et 2000, et surtout la catastrophe du lac Nyos en 1986 (dégazage brutal de \ce{CO2} d'origine volcanique, plus de \num{1700} morts), montrent l'importance de comprendre ces phénomènes pour mieux s'en protéger.
\end{savaistu}

\begin{exercice}[Lecture d'un profil bathymétrique]
Un profil bathymétrique réalisé dans l'océan Indien donne les profondeurs suivantes d'ouest en est : \SI{-5200}{m}, \SI{-4600}{m}, \SI{-2800}{m}, \SI{-3900}{m}, \SI{-2700}{m}, \SI{-4700}{m}, \SI{-5300}{m}.
\begin{enumerate}
\item Repérer la position de la dorsale et celle du rift. Justifier.
\item Calculer la hauteur approximative du relief de la dorsale par rapport aux plaines abyssales.
\item Un écho revient au sondeur au bout de \SI{7,0}{s} au-dessus d'une plaine abyssale. Vérifier que la profondeur mesurée est cohérente avec le profil ($v = \SI{1500}{m.s^{-1}}$).
\end{enumerate}
\end{exercice}

\begin{corrige}[Exercice : lecture d'un profil bathymétrique]
\begin{enumerate}
\item La dorsale correspond aux profondeurs les plus faibles, autour de \SI{-2800}{m} et \SI{-2700}{m} ; entre ces deux points, l'enfoncement à \SI{-3900}{m} matérialise la vallée axiale : le \cle{rift}. Les valeurs symétriques de part et d'autre (\SI{-4600}{m}/\SI{-4700}{m} puis \SI{-5200}{m}/\SI{-5300}{m}) confirment l'organisation symétrique du plancher.
\item Hauteur du relief : $5300 - 2700 = \SI{2600}{m}$ environ au-dessus des plaines abyssales.
\item $d = \dfrac{v \times t}{2} = \dfrac{1500 \times 7{,}0}{2} = \SI{5250}{m}$, valeur tout à fait cohérente avec les profondeurs des plaines abyssales du profil (\SI{-5200}{m} à \SI{-5300}{m}).
\end{enumerate}
\end{corrige}

La topographie des fonds océaniques est donc tout sauf chaotique : dorsales axiales creusées d'un rift, plaines abyssales symétriques, fosses marginales, décalages par failles transformantes. Reste à comprendre \emph{pourquoi} ce relief est organisé ainsi : c'est ce que révèle l'étude du fonctionnement du rift et de l'expansion du plancher océanique, objet de la section suivante.

\coursec{Le rift et l'expansion du plancher océanique}

Dans la section précédente, nous avons décrit la morphologie des fonds océaniques : les plaines abyssales, les fosses et surtout les \cle{dorsales océaniques}, ces chaînes montagneuses sous-marines qui parcourent tous les océans du globe sur plus de \SI{60000}{km}. Nous avons remarqué que l'axe de chaque dorsale est creusé d'une vallée étroite, le \cle{rift}. Pourquoi une montagne est-elle fendue en son sommet ? Pourquoi les séismes et les volcans sous-marins s'y concentrent-ils ? C'est à ces questions que répond la présente section : nous allons montrer que le rift est la cicatrice d'un phénomène immense et permanent, la \cle{naissance de la croûte océanique}, et que les océans s'élargissent en permanence.

\courssub{Observation : le rift, une vallée axiale bordée de failles}

\begin{definition}[Le rift océanique]
Le \cle{rift} est une vallée axiale étroite (largeur de 10 à \SI{40}{km}, profondeur de 500 à \SI{1500}{m}) qui occupe le sommet des dorsales océaniques. Il est bordé de part et d'autre par des \cle{failles normales} (failles dont le compartiment central s'est affaissé par rapport aux compartiments latéraux). Le rift est le siège d'un \cle{volcanisme basaltique} (laves fluides, pauvres en silice) et de \cle{séismes superficiels} (foyers à moins de \SI{10}{km} de profondeur), de faible magnitude mais fréquents.
\end{definition}

\begin{experience}[Observation du rift de la dorsale médio-atlantique]
\textbf{Document.} Les campagnes bathymétriques et les plongées de submersibles (par exemple le \emph{Nautile}) sur la dorsale médio-atlantique ont fourni les données suivantes :
\begin{itemize}
\item la vallée axiale est encaissée entre deux épaufrures bordées de failles normales décalant le relief en gradins ;
\item on observe des coulées de lave récentes, en forme de coussins (\emph{pillow-lavas}), et des fissures d'où s'échappent des fluides chauds ;
\item les séismes enregistrés sont superficiels et alignés sur l'axe ;
\item le flux de chaleur géothermique est maximal à l'axe et décroît quand on s'en éloigne.
\end{itemize}
\textbf{Interprétation.} Les failles normales traduisent un \emph{étirement} de la lithosphère ; la lave basaltique, dont la composition chimique est proche de celle du manteau, indique une remontée de matériaux profonds ; la chaleur élevée confirme la présence de matériaux chauds sous l'axe. Tout se passe comme si le plancher océanique était \emph{écartelé} en deux au niveau du rift.
\end{experience}

\begin{attention}[Faille normale ou faille inverse ?]
Dans une \cle{faille normale}, le compartiment situé au-dessus du plan de faille (le \emph{toit}) descend par rapport au \emph{mur} : c'est la signature d'une \textbf{distension} (étirement). Dans une faille inverse, le toit remonte : c'est la signature d'une compression. Au niveau du rift, les failles sont \emph{normales} : le rift est donc une zone d'étirement, non de compression. Ne jamais inverser cette conclusion dans une interprétation de coupe géologique.
\end{attention}

\courssub{Interprétation : le rift, une zone de divergence et d'accrétion}

\begin{definition}[Accrétion océanique]
L'\cle{accrétion océanique} est la formation de nouvelle croûte océanique au niveau de l'axe des dorsales, par solidification du magma basaltique issu de la fusion partielle du manteau. La dorsale est donc appelée \cle{zone d'accrétion} ou \cle{zone de divergence} : deux plaques lithosphériques s'y écartent l'une de l'autre.
\end{definition}

\begin{definition}[Expansion océanique]
L'\cle{expansion océanique} (ou expansion du plancher océanique) est l'écartement continu de deux plaques lithosphériques de part et d'autre de l'axe d'une dorsale, accompagné de la création permanente de croûte océanique neuve au niveau du rift. L'océan s'élargit ainsi symétriquement de part et d'autre de la dorsale.
\end{definition}

Le mécanisme peut se résumer ainsi : sous la dorsale, des matériaux \cle{mantelliques chauds} (péridotites de l'asthénosphère) remontent par convection. En remontant, ils subissent une \emph{décompression} qui provoque leur fusion partielle (nous démontrerons ce point expérimentalement dans la section suivante). Le liquide basaltique ainsi produit s'injecte dans la fissure axiale du rift : une partie refroidit lentement en profondeur (gabbros), une partie sort en surface et fige au contact de l'eau de mer (basaltes en coussins). Cette croûte neuve est ensuite poussée latéralement, de part et d'autre, par les injections suivantes, comme deux tapis roulants s'éloignant de l'axe.

\begin{definition}[La lithosphère]
La \cle{lithosphère} est l'enveloppe rigide externe de la Terre, d'une épaisseur de 70 à \SI{100}{km} sous les océans (jusqu'à \SI{250}{km} sous les continents). Elle comprend la \emph{croûte} (océanique ou continentale) et la partie supérieure rigide du \emph{manteau} (manteau lithosphérique). Elle repose sur l'\cle{asthénosphère}, zone du manteau plus ductile (déformable) qui permet le mouvement des plaques.
\end{definition}

\begin{popfigure}
\begin{center}
\begin{tikzpicture}[scale=0.9, every node/.style={font=\small}]
% Asthénosphère
\fill[poporange!45] (-6.5,-3.2) rectangle (6.5,-1.2);
\node at (0,-2.9) {\textbf{Asthénosphère (manteau ductile, péridotites chaudes)}};
% Remontée centrale
\fill[poporange!75] (-1.1,-3.2) -- (-0.5,-0.4) -- (0.5,-0.4) -- (1.1,-3.2) -- cycle;
\node[font=\footnotesize, align=center] at (0,-1.6) {remontée\\mantellique};
% Flèches de convection
\draw[-{Stealth[length=3mm]}, popdark, line width=1.1pt] (-3.6,-2.3) -- (-2.2,-2.3);
\draw[-{Stealth[length=3mm]}, popdark, line width=1.1pt] (3.6,-2.3) -- (2.2,-2.3);
% Lithosphère gauche et droite
\fill[popblue!25] (-6.5,0.2) -- (-0.5,0.2) -- (-0.5,-0.4) -- (-1.1,-1.2) -- (-6.5,-1.2) -- cycle;
\fill[popblue!25] (6.5,0.2) -- (0.5,0.2) -- (0.5,-0.4) -- (1.1,-1.2) -- (6.5,-1.2) -- cycle;
% Mer
\fill[popblue!12] (-6.5,0.2) rectangle (6.5,1.6);
\node[popblue, font=\footnotesize] at (5.3,1.35) {eau de mer};
% Vallée axiale (rift)
\fill[popblue!12] (-0.5,0.2) -- (-0.5,1.0) -- (0.5,1.0) -- (0.5,0.2) -- cycle;
\draw[popdark, line width=1pt] (-0.5,0.2) -- (-0.5,1.0);
\draw[popdark, line width=1pt] (0.5,0.2) -- (0.5,1.0);
\node[font=\footnotesize, align=center] at (0,1.35) {\textbf{rift}};
% Failles normales
\draw[popdark, line width=1pt] (-2.6,0.2) -- (-2.2,-1.2);
\draw[popdark, line width=1pt] (-4.2,0.2) -- (-3.8,-1.2);
\draw[popdark, line width=1pt] (2.6,0.2) -- (2.2,-1.2);
\draw[popdark, line width=1pt] (4.2,0.2) -- (3.8,-1.2);
\node[font=\footnotesize, align=center] at (-3.4,0.75) {failles\\normales};
\node[font=\footnotesize, align=center] at (3.4,0.75) {failles\\normales};
% Injections basaltiques
\fill[popgreen!60] (-0.18,-0.4) rectangle (0.18,0.95);
\node[font=\footnotesize, align=center] at (0,-0.75) {injection de\\basalte};
% Pillow lavas
\foreach \x in {-0.75,-0.35,0.35,0.75} {\draw[popgreen!80, fill=popgreen!40] (\x,0.2) circle (0.16);}
\node[font=\footnotesize, align=center, anchor=west] at (1.0,-0.55) {pillow-lavas};
\draw[popgreen!80, -{Stealth[length=2mm]}] (1.0,-0.45) -- (0.55,0.15);
% Flèches d'écartement
\draw[-{Stealth[length=3.5mm]}, poppink, line width=1.4pt] (-1.6,1.15) -- (-3.0,1.15);
\draw[-{Stealth[length=3.5mm]}, poppink, line width=1.4pt] (1.6,1.15) -- (3.0,1.15);
\node[poppink, font=\footnotesize] at (-2.3,1.45) {plaque A};
\node[poppink, font=\footnotesize] at (2.3,1.45) {plaque B};
% Âges symétriques
\node[font=\footnotesize] at (-5.6,0.6) {âge croissant $\leftarrow$};
\node[font=\footnotesize] at (5.6,0.6) {$\rightarrow$ âge croissant};
\node[font=\footnotesize] at (0,-3.55) {};
\end{tikzpicture}
\end{center}
\legende{Coupe schématique d'une dorsale océanique et de son rift. Sous l'axe, la remontée de péridotites chaudes de l'asthénosphère (orange) alimente des injections de basalte (vert) qui créent la croûte neuve. Les failles normales (traits noirs inclinés) encadrent la vallée axiale. Les plaques A et B s'écartent symétriquement (flèches roses) : l'âge des basaltes croît quand on s'éloigne de l'axe.}
\end{popfigure}

\courssub{Les preuves de l'expansion océanique}

\textbf{Première preuve : l'âge des basaltes croît de part et d'autre de l'axe.}

Les forages profonds réalisés dans les océans (programmes DSDP puis ODP) ont permis de prélever les basaltes du plancher océanique et de les dater par les méthodes de radiochronologie (K/Ar). Le résultat est frappant : \emph{aucun basalte océanique n'est plus vieux que 180 Ma environ}, alors que les continents conservent des roches âgées de plusieurs milliards d'années. De plus, plus un basalte est éloigné de l'axe de la dorsale, plus il est âgé, et la répartition des âges est \emph{symétrique} par rapport à l'axe : à égale distance de part et d'autre, on trouve des basaltes de même âge. C'est exactement ce que prédit le modèle de l'expansion : la croûte se forme à l'axe, puis vieillit en s'éloignant.

\begin{popfigure}
\begin{center}
\begin{tikzpicture}
\begin{axis}[
width=0.85\linewidth, height=6.5cm,
xlabel={Distance à l'axe de la dorsale (km)},
ylabel={Âge du basalte (Ma)},
xmin=0, xmax=550, ymin=0, ymax=55,
grid=major, grid style={popline!40},
legend style={at={(0.03,0.97)}, anchor=north west, font=\footnotesize},
]
\addplot[popcurve, domain=0:500, samples=100] {0.1*x};
\addlegendentry{Droite d'étalonnage : pente $= 10\ \mathrm{km/Ma}$}
\addplot[only marks, mark=*, mark size=2pt, poporange] coordinates {(100,10) (200,20) (300,30) (400,40) (500,50)};
\addlegendentry{Datations de basaltes forés}
\addplot[only marks, mark=square*, mark size=2pt, poppurple] coordinates {(150,16) (350,33) (450,47)};
\addlegendentry{Autres mesures (dispersion expérimentale)}
\end{axis}
\end{tikzpicture}
\end{center}
\legende{Âge du plancher océanique en fonction de la distance à l'axe de la dorsale. Les points de mesure s'alignent sur une droite passant par l'origine : la vitesse d'écartement est constante. La pente de la droite (ici $10\ \mathrm{km/Ma}$, soit \SI{1}{cm/an}) donne le demi-taux d'expansion.}
\end{popfigure}

\textbf{Deuxième preuve : les anomalies magnétiques symétriques.}

Les basaltes contiennent des minéraux ferrimagnétiques (magnétite). En refroidissant sous la \cle{température de Curie} (environ \SI{580}{\celsius} pour la magnétite), ces minéraux s'aimantent dans le sens du champ magnétique terrestre du moment, comme de minuscules boussoles figées. Or, on sait que le champ magnétique terrestre s'est \emph{inversé} de nombreuses fois au cours des temps géologiques (le pôle Nord magnétique devenant pôle Sud et inversement). Les mesures magnétométriques effectuées au-dessus des dorsales révèlent des bandes alternées d'\cle{anomalies magnétiques positives} (aimantation dans le sens du champ actuel) et \cle{négatives} (aimantation inversée), \emph{parfaitement symétriques par rapport à l'axe de la dorsale}. C'est la preuve décisive : chaque bande s'est formée à l'axe pendant une période de polarité donnée, puis a été séparée en deux moitiés qui ont dérivé de part et d'autre. Cette découverte (Vine et Matthews, 1963) a définitivement validé l'hypothèse de l'expansion océanique proposée par Hess en 1962.

\begin{popfigure}
\begin{center}
\begin{tikzpicture}[scale=1.0, every node/.style={font=\footnotesize}]
% Axe
\fill[poporange!70] (-0.4,0) rectangle (0.4,3.2);
\node[align=center] at (0,3.55) {\textbf{axe de la dorsale}\\(rift)};
% Bandes symétriques
% Gauche
\fill[popblue!55] (-1.2,0) rectangle (-0.4,3.2);
\fill[popcream] (-2.4,0) rectangle (-1.2,3.2);
\fill[popblue!55] (-3.0,0) rectangle (-2.4,3.2);
\fill[popcream] (-4.0,0) rectangle (-3.0,3.2);
\fill[popblue!55] (-5.4,0) rectangle (-4.0,3.2);
% Droite
\fill[popblue!55] (0.4,0) rectangle (1.2,3.2);
\fill[popcream] (1.2,0) rectangle (2.4,3.2);
\fill[popblue!55] (2.4,0) rectangle (3.0,3.2);
\fill[popcream] (3.0,0) rectangle (4.0,3.2);
\fill[popblue!55] (4.0,0) rectangle (5.4,3.2);
% Contours
\draw[popdark] (-5.4,0) rectangle (5.4,3.2);
% Flèches d'écartement
\draw[-{Stealth[length=3mm]}, poppink, line width=1.3pt] (-0.6,2.6) -- (-2.0,2.6);
\draw[-{Stealth[length=3mm]}, poppink, line width=1.3pt] (0.6,2.6) -- (2.0,2.6);
% Annotations
\node at (-2.9,-0.45) {anomalie $+$};
\node at (2.9,-0.45) {anomalie $+$};
\node at (-1.8,-0.45) {anomalie $-$};
\node at (1.8,-0.45) {anomalie $-$};
\node[align=center] at (0,-1.05) {âge croissant $\longleftarrow$ \hspace{2.2cm} $\longrightarrow$ âge croissant};
% Légende couleurs
\fill[popblue!55] (-5.4,4.0) rectangle (-5.0,4.4); \node[anchor=west] at (-4.9,4.2) {polarité normale (comme aujourd'hui)};
\fill[popcream, draw=popdark] (0.6,4.0) rectangle (1.0,4.4); \node[anchor=west] at (1.1,4.2) {polarité inversée};
\end{tikzpicture}
\end{center}
\legende{Anomalies magnétiques du plancher océanique : bandes alternées de polarité normale (bleu) et inversée (crème), symétriques par rapport à l'axe de la dorsale (orange). Chaque bande s'est aimantée à l'axe lors de sa formation, puis a dérivé latéralement : le plancher océanique est un « enregistreur » des inversions du champ magnétique terrestre.}
\end{popfigure}

\courssub{Le taux d'expansion : mesure et calcul}

\begin{methode}[Calculer un taux d'expansion océanique]
\begin{enumerate}
\item Repérer sur la carte (ou le profil) la \textbf{distance} $d$ qui sépare un basalte daté de l'\textbf{axe} de la dorsale (et non la distance entre deux basaltes situés de part et d'autre !).
\item Noter l'\textbf{âge} $a$ de ce basalte, donné par la datation radiochronologique.
\item Appliquer la relation : \[ v = \dfrac{d}{a} \]
\item Convertir les unités : $1\ \mathrm{km/Ma} = \dfrac{10^{5}\ \mathrm{cm}}{10^{6}\ \mathrm{ans}} = \SI{0,1}{cm/an}$, donc $10\ \mathrm{km/Ma} = \SI{1}{cm/an}$.
\item Si l'énoncé donne la \textbf{largeur totale} d'une bande ou de l'océan entre deux continents, penser à \textbf{diviser par 2} pour obtenir le demi-taux d'expansion (vitesse d'une seule plaque par rapport à l'axe).
\end{enumerate}
\end{methode}

\begin{exemplebox}[Un calcul de référence]
Un basalte prélevé à \SI{100}{km} de l'axe de la dorsale médio-atlantique est âgé de 10 Ma. Le demi-taux d'expansion est :
\[ v = \dfrac{d}{a} = \dfrac{\SI{100}{km}}{\SI{10}{Ma}} = 10\ \mathrm{km/Ma} = \SI{1}{cm/an} \]
L'Atlantique s'élargit donc d'environ \SI{2}{cm/an} au total (\SI{1}{cm/an} de chaque côté de l'axe), soit la vitesse de croissance de nos ongles. Selon les dorsales, les taux varient de 1 à \SI{5}{cm/an} (dorsales lentes comme l'Atlantique ; dorsales rapides comme le Pacifique Est, jusqu'à \SI{15}{cm/an} localement).
\end{exemplebox}

\begin{exoresolu}[Taux d'expansion du Pacifique Est]
Sur la dorsale du Pacifique Est, une anomalie magnétique datée de 4 Ma se retrouve sous forme de deux bandes symétriques dont les bords externes sont distants de \SI{320}{km}. Calculer le demi-taux d'expansion de cette dorsale en cm/an.
\tcblower
\textbf{Solution.} La largeur totale de \SI{320}{km} correspond à \emph{deux} moitiés symétriques ayant dérivé de part et d'autre de l'axe. La distance parcourue par chaque moitié est donc :
\[ d = \dfrac{320}{2} = \SI{160}{km} \]
Le demi-taux d'expansion est :
\[ v = \dfrac{d}{a} = \dfrac{\SI{160}{km}}{\SI{4}{Ma}} = 40\ \mathrm{km/Ma} \]
Conversion : $40\ \mathrm{km/Ma} = 40 \times \SI{0,1}{cm/an} = \SI{4}{cm/an}$.
\textbf{Conclusion.} Le demi-taux d'expansion de la dorsale du Pacifique Est est de \SI{4}{cm/an} : c'est une dorsale \emph{rapide}, environ quatre fois plus rapide que la dorsale médio-atlantique.
\end{exoresolu}

\begin{attention}[Les deux pièges classiques du calcul]
\begin{itemize}
\item \textbf{Oublier de diviser par 2} quand la distance donnée est la largeur totale entre deux bandes symétriques ou entre deux continents : on obtient alors un taux double de la valeur attendue.
\item \textbf{Se tromper dans la conversion} : retenir la correspondance $10\ \mathrm{km/Ma} = \SI{1}{cm/an}$ et vérifier systématiquement l'ordre de grandeur (les taux d'expansion sont toujours de l'ordre du cm/an, jamais du m/an).
\end{itemize}
\end{attention}

\courssub{Structure de la lithosphère océanique (savoir admis)}

Les dragages, les forages et l'étude des \cle{ophiolites} (fragments de croûte océanique charriés sur les continents, par exemple dans les chaînes de montagnes) ont permis d'établir le modèle de la croûte océanique. De haut en bas, on distingue :
\begin{itemize}
\item une couche de \cle{sédiments} (quelques centaines de mètres, d'épaisseur croissante avec l'âge du plancher) ;
\item les \cle{basaltes en coussins} ou \emph{pillow-lavas} (environ 1 à \SI{2}{km}) : laves figées au contact de l'eau de mer ;
\item les \cle{gabbros} (environ 4 à \SI{5}{km}) : roches de \emph{même composition chimique} que les basaltes mais à gros cristaux, formées par refroidissement lent du magma en profondeur ;
\item les \cle{péridotites} du manteau lithosphérique, sous la \cle{discontinuité de Mohorovičić} (Moho).
\end{itemize}
L'ensemble croûte + manteau supérieur rigide constitue la lithosphère océanique, dont l'épaisseur augmente en s'éloignant de la dorsale (elle refroidit et s'épaissit avec l'âge).

\begin{savaistu}[Les pillow-lavas, témoins d'un refroidissement brutal]
Quand la lave basaltique à environ \SI{1200}{\celsius} s'écoule sous l'eau de mer, sa surface fige quasi instantanément en une pellicule de verre, tandis que l'intérieur encore fluide gonfle puis perce l'enveloppe, formant un nouveau « coussin ». Il en résulte des empilements de formes arrondies, les \cle{pillow-lavas}, que les plongeurs des submersibles observent au niveau du rift. Retrouver des pillow-lavas dans une chaîne de montagnes (ophiolites) prouve qu'un océan existait autrefois à cet endroit. Au Cameroun, la ligne volcanique du Mont Cameroun produit des basaltes comparables, mais émis à l'air libre : ils forment des coulées cordées ou en blocs, et non des coussins.
\end{savaistu}

\begin{attention}[Le magma ne sort pas d'un « océan de magma »]
Erreur fréquente : imaginer que la lave du rift provient d'une mer de magma permanente sous la croûte. En réalité, le manteau sous la dorsale est \emph{solide} (péridotites) ; le magma basaltique est produit \emph{localement et ponctuellement} par \cle{fusion partielle} (quelques pour cent du manteau fond) due à la \textbf{décompression} lors de la remontée asthénosphérique. Nous le démontrerons expérimentalement dans la section suivante grâce aux courbes de fusion des péridotites.
\end{attention}

\courssub{TP : réalisation d'une maquette du rift et du plancher océanique}

\begin{experience}[Maquette de l'expansion océanique avec des matériaux locaux]
\textbf{Matériel.} Une boîte en carton dont on découpe le fond en deux moitiés ; deux feuilles de papier kraft (ou deux bandes de tissu) ; des feutres ou de la peinture ; de la colle ; éventuellement de la pâte à modeler ou de la boue séchée pour figurer le relief.

\textbf{Protocole.}
\begin{enumerate}
\item Pratiquer une fente centrale dans le fond de la boîte : elle représente le \emph{rift}.
\item Passer sous la boîte deux bandes de papier accolées, dont les extrémités ressortent par la fente et sont tirées vers l'extérieur, de part et d'autre.
\item Colorier sur les bandes des \emph{bandes alternées} (bleu / blanc) symétriques par rapport à la fente : ce sont les anomalies magnétiques.
\item Tirer lentement et simultanément les deux bandes vers l'extérieur tout en faisant « émerger » du papier neuf (colorié en rouge-orangé) par la fente : c'est la croûte neuve injectée au rift.
\item Figurer les failles normales par des gradins en carton de part et d'autre de la fente, et les pillow-lavas par de petites boules de pâte à modeler alignées sur l'axe.
\end{enumerate}

\textbf{Observations.} Les bandes colorées les plus anciennes sont les plus éloignées de la fente ; la symétrie axiale est respectée ; plus on tire vite, plus la « dorsale » s'élargit rapidement.

\textbf{Interprétation.} La maquette reproduit les trois faits essentiels : création de croûte neuve à l'axe, écartement symétrique des plaques, vieillissement du plancher avec la distance à l'axe. Elle permet de visualiser pourquoi l'âge des basaltes et les anomalies magnétiques sont des enregistrements chronologiques de l'expansion.
\end{experience}

\begin{aretenir}[L'essentiel de la section]
\begin{itemize}
\item Le \cle{rift} est une vallée axiale bordée de failles normales, siège d'un volcanisme basaltique et de séismes superficiels : c'est une zone de \textbf{divergence}.
\item L'\cle{accrétion océanique} crée en permanence de la croûte neuve à l'axe des dorsales, par fusion partielle du manteau remontant (décompression).
\item L'\cle{expansion océanique} est prouvée par la croissance symétrique de l'âge des basaltes et par la symétrie des anomalies magnétiques de part et d'autre de l'axe.
\item Le taux d'expansion se calcule par $v = d/a$ (1 à \SI{5}{cm/an} en général ; diviser par 2 si la distance est bilatérale).
\item La croûte océanique comprend, de haut en bas : sédiments, pillow-lavas, gabbros ; elle repose sur les péridotites du manteau lithosphérique.
\end{itemize}
\end{aretenir}

\begin{exercice}[Pour t'entraîner]
Sur la dorsale médio-atlantique, un basalte foré à \SI{250}{km} à l'ouest de l'axe est âgé de 25 Ma ; un basalte foré à \SI{250}{km} à l'est est âgé de 25 Ma également.
\begin{enumerate}
\item Que traduit l'égalité de ces deux âges ?
\item Calculer le demi-taux d'expansion de cette dorsale en cm/an.
\item Quelle serait la largeur de l'océan créé en 100 Ma à ce rythme ?
\end{enumerate}
\end{exercice}

\begin{corrige}[Exercice]
\begin{enumerate}
\item L'égalité des âges à égale distance de l'axe traduit la \textbf{symétrie} de l'expansion océanique : la croûte se forme à l'axe puis est repoussée à la même vitesse de part et d'autre.
\item $v = \dfrac{d}{a} = \dfrac{\SI{250}{km}}{\SI{25}{Ma}} = 10\ \mathrm{km/Ma} = \SI{1}{cm/an}$.
\item En 100 Ma, chaque plaque parcourt $100\ \mathrm{Ma} \times 10\ \mathrm{km/Ma} = \SI{1000}{km}$. La largeur totale créée est donc $2 \times 1000 = \SI{2000}{km}$.
\end{enumerate}
\end{corrige}

Ainsi, le rift n'est pas une simple cassure : c'est la \emph{fabrique} du plancher océanique. Mais d'où vient exactement le magma basaltique injecté à l'axe ? Pour le comprendre, il faut examiner les données expérimentales de la fusion des péridotites : c'est l'objet de la section suivante.

\coursec{Naissance des magmas : données expérimentales de fusion des péridotites}

Nous venons de voir que le plancher océanique naît au niveau du rift, où du \cle{magma} basaltique s'épanche en permanence et se solidifie pour former la croûte océanique. Nous verrons aussi que les zones de subduction sont le siège d'un volcanisme explosif (andésites). Une question fondamentale se pose alors : \emph{d'où vient ce magma ?} Le manteau terrestre est-il un océan de roche liquide ? Pourquoi fond-il précisément sous les dorsales et au-dessus des plaques plongeantes ? C'est à ces questions que répond la \cle{pétrologie expérimentale}, en reproduisant en laboratoire les conditions de pression et de température qui règnent dans les profondeurs de la Terre.

\courssub{Le problème : d'où vient le magma ?}

\begin{experience}[Une observation qui intrigue]
\textbf{Observation.} Les sondages profonds, l'étude des ondes sismiques et les enclaves remontées par les volcans montrent que le manteau terrestre est constitué d'une roche \textbf{solide}, la péridotite. Pourtant, les volcans des dorsales (Islande, dorsale médio-atlantique) et des zones de subduction (Mont Cameroun, chaîne des Cascades, Andes) émettent des laves \textbf{liquides} à plus de \SI{1000}{\celsius}.

\textbf{Problème.} Comment une roche solide du manteau peut-elle produire un magma liquide, et pourquoi la fusion se produit-elle préférentiellement sous les dorsales et les zones de subduction ?

\textbf{Hypothèses.} (H1) La température augmente localement jusqu'à dépasser le point de fusion de la roche. (H2) D'autres facteurs que la température (pression, présence d'eau) modifient les conditions de fusion.

\textbf{Vérification.} Il faut déterminer expérimentalement à quelle température la péridotite commence à fondre en fonction de la pression, puis comparer ces conditions aux conditions réelles du manteau.
\end{experience}

\courssub{Les outils du raisonnement : péridotite, solidus, géotherme}

\begin{definition}[Péridotite]
La \cle{péridotite} est une roche \textbf{magmatique grenue} (à cristaux visibles), de couleur verte, constituée essentiellement d'\textbf{olivine} et de \textbf{pyroxènes}. C'est la roche constitutive du \textbf{manteau supérieur} de la Terre. On la trouve à la surface sous forme d'enclaves arrachées par les magmas ou dans les massifs d'ophiolites (fragments de lithosphère océanique charriés sur les continents).
\end{definition}

\begin{definition}[Solidus]
Le \cle{solidus} d'une roche est la courbe qui, dans un diagramme pression--température, sépare le domaine où la roche est \textbf{totalement solide} du domaine où elle commence à \textbf{fondre}. En dessous du solidus : roche solide ; au-dessus : présence de liquide (magma) en proportion croissante avec la température. Une roche étant un mélange de minéraux, elle ne fond pas à une température unique : le solidus marque le \emph{début} de la fusion.
\end{definition}

\begin{definition}[Géotherme]
Le \cle{géotherme} est la courbe qui représente la variation de la \textbf{température en fonction de la profondeur} (ou de la pression) à l'intérieur de la Terre, dans un contexte géodynamique donné. Le géotherme n'est pas le même partout : il est \textbf{élevé} sous les dorsales (manteau chaud qui remonte) et \textbf{abaissé} dans les zones de subduction (plaque froide qui s'enfonce).
\end{definition}

\begin{attention}[Le manteau n'est pas un océan de magma]
Erreur très fréquente au Baccalauréat : écrire que « le manteau est fondu » ou que « le magma provient d'une couche liquide permanente ». C'est \textbf{faux} : les ondes sismiques S (qui ne se propagent pas dans les liquides) traversent le manteau, ce qui prouve qu'il est \textbf{solide}. La fusion est \textbf{partielle} (quelques pourcents de la roche fondent) et \textbf{locale} (uniquement là où les conditions pression--température franchissent le solidus). Le magma est un liquide qui s'extrait d'une roche solide, comme l'eau qu'on essore d'une éponge mouillée.
\end{attention}

\courssub{La pétrologie expérimentale : fondre une péridotite en laboratoire}

\begin{experience}[Détermination expérimentale du solidus de la péridotite]
\textbf{Matériel.} Échantillons de péridotite naturelle (quelques milligrammes), capsule métallique scellée (or, platine), presse à haute pression (enclumes de diamant ou presse à piston-cylindre) capable d'atteindre plusieurs gigapascals (GPa), four permettant de dépasser \SI{1500}{\celsius}, thermocouples. Deux séries d'expériences : péridotite \textbf{sèche} (anhydre) et péridotite \textbf{hydratée} (eau ajoutée dans la capsule).

\textbf{Protocole.}
\begin{enumerate}
\item Placer l'échantillon dans la capsule et le porter à une pression $P$ donnée (simulant une profondeur : 1 GPa $\approx$ 33 km).
\item Chauffer progressivement et noter la température à laquelle apparaissent les \textbf{premières gouttes de liquide} (observées après trempe, au microscope : présence de verre).
\item Recommencer pour différentes pressions, avec et sans eau.
\item Reporter les points expérimentaux dans un diagramme ($T$, $P$) et tracer les courbes de début de fusion.
\end{enumerate}

\textbf{Résultats.}
\begin{itemize}
\item Péridotite \textbf{anhydre} (sèche) : le début de fusion se situe vers \SI{1050}{\celsius} à la surface (0 GPa) et monte jusqu'à environ \SI{1450}{\celsius} à 3 GPa ($\approx$ 100 km) : le solidus anhydre est une courbe \textbf{croissante} avec la pression.
\item Péridotite \textbf{hydratée} : le début de fusion est abaissé de \textbf{plusieurs centaines de degrés} (vers \SI{900}{\celsius}--\SI{1100}{\celsius} sous 3 GPa) : le solidus hydraté se situe nettement \textbf{en dessous} du solidus anhydre.
\end{itemize}

\textbf{Interprétation.} La température de début de fusion de la péridotite \textbf{augmente avec la pression} et \textbf{diminue fortement en présence d'eau}. Ces deux résultats expérimentaux sont la clé de la genèse des magmas.
\end{experience}

\begin{propriete}[Propriétés admises, établies expérimentalement]
\begin{enumerate}
\item La température de début de fusion de la péridotite \textbf{augmente lorsque la pression augmente} : plus une roche est profonde, plus il faut la chauffer pour la fondre. Réciproquement, une \textbf{diminution de pression} (décompression) peut faire passer une roche chaude dans le domaine de la fusion \textbf{sans apport de chaleur}.
\item La présence d'\textbf{eau abaisse la température de début de fusion} des péridotites de plusieurs centaines de degrés Celsius (solidus hydraté inférieur au solidus anhydre). Cette propriété est \textbf{admise} au programme : elle résulte des expériences de pétrologie et n'est pas à démontrer.
\end{enumerate}
\end{propriete}

\courssub{Exploitation du diagramme pression--température}

\begin{popfigure}
\begin{center}
\begin{tikzpicture}
\begin{axis}[
  width=0.95\linewidth, height=9cm,
  xlabel={Température (\si{\celsius})},
  ylabel={Pression (GPa)},
  xmin=600, xmax=1700, ymin=0, ymax=3.5,
  xtick={600,800,1000,1200,1400,1600},
  ytick={0,1,2,3},
  y dir=reverse,
  grid=both, grid style={popline!40},
  legend style={at={(0.02,0.98)}, anchor=north west, font=\small, draw=popline, fill=popcream},
  clip=false
]
% Solidus anhydre (dry) - Katz et al. 2003 approx
\addplot[popcurve, popink, very thick, mark=none] coordinates {
  (1050,0) (1150,1) (1300,2) (1450,3) (1550,3.5)
};
\addlegendentry{Solidus péridotite anhydre (sèche)}

% Solidus hydraté (wet) - approx
\addplot[popcurve, popblue, very thick, mark=none] coordinates {
  (800,0) (900,1) (1000,2) (1100,3) (1150,3.5)
};
\addlegendentry{Solidus péridotite hydratée}

% Géotherme dorsale (chaud, adiabatique ~0.3-0.5 K/km, Tpot ~1350)
\addplot[popcurve, poporange, very thick, mark=none] coordinates {
  (0,0) (1300,1) (1400,2) (1450,3) (1480,3.5)
};
\addlegendentry{Géotherme de dorsale (chaud)}

% Géotherme subduction (froid)
\addplot[popcurve, popgreen, very thick, mark=none] coordinates {
  (0,0) (600,1) (800,2) (900,3) (950,3.5)
};
\addlegendentry{Géotherme de subduction (froid)}

% Point A : Intersection Dorsale / Solidus Anhydre ~ (1320, 1.6) -> ~53 km? Wait. Ridge melting starts ~60-80km (2-2.5 GPa).
% Let's adjust Ridge Geotherm to cross dry solidus at ~2.2 GPa.
% Dry solidus at 2.2 GPa ~ 1330 C. Ridge Geotherm at 2.2 GPa ~ 1330 C.
% Subduction Geotherm crosses Wet Solidus at ~2.8 GPa. Wet solidus at 2.8 GPa ~ 1080 C. Subduction Geotherm at 2.8 GPa ~ 1080 C.

% Point A
\node[circle, fill=popink, inner sep=2.5pt, label={[popink, font=\small]above right:{A : fusion par décompression}}] at (axis cs:1330,2.2) {};
% Point B
\node[circle, fill=popblue, inner sep=2.5pt, label={[popblue, font=\small]below right:{B : fusion par hydratation}}] at (axis cs:1080,2.8) {};

% Flèche Décompression (verticale vers le haut = baisse P)
\draw[-{Stealth[length=3mm]}, popink, very thick] (axis cs:1330,2.8) -- (axis cs:1330,2.3);
\node[popink, font=\small, rotate=90, anchor=south] at (axis cs:1350,2.55) {Décompression};

% Flèche Hydratation (horizontale vers la gauche = baisse T solidus)
\draw[-{Stealth[length=3mm]}, popblue, very thick] (axis cs:1350,2.8) -- (axis cs:1120,2.8);
\node[popblue, font=\small, anchor=south] at (axis cs:1230,2.85) {Apport d'eau};

\end{axis}
\end{tikzpicture}
\end{center}
\legende{Diagramme pression--température (pression croissante vers le bas, comme la profondeur). Les courbes de début de fusion (solidus) de la péridotite sèche (rose) et hydratée (bleue) sont comparées aux géothermes d'une dorsale (orange) et d'une zone de subduction (verte). Au point A, le géotherme de dorsale croise le solidus anhydre vers 2,2 GPa ($\approx$ 73 km) : la péridotite qui remonte entre en fusion \textbf{par décompression}. Au point B, le géotherme froid de subduction, qui ne croiserait jamais le solidus anhydre, croise le solidus hydraté vers 2,8 GPa ($\approx$ 92 km) : la fusion y est déclenchée \textbf{par l'apport d'eau}.}
\end{popfigure}

\begin{methode}[Exploiter un diagramme pression--température (méthode exigible au Bac)]
Face à un document donnant un solidus et un ou plusieurs géothermes, procède toujours ainsi :
\begin{enumerate}
\item \textbf{Repérer les axes} : en général la température en abscisse, la pression (ou la profondeur) en ordonnée, croissante vers le bas. Vérifier les unités (\si{\celsius}, GPa ou km).
\item \textbf{Identifier le solidus} : c'est la courbe de début de fusion. Tout point situé \emph{à droite} (côté chaud) du solidus correspond à une roche partiellement fondue ; tout point \emph{à gauche} correspond à une roche solide.
\item \textbf{Tracer ou repérer le géotherme} du contexte étudié (dorsale : géotherme chaud ; subduction : géotherme froid).
\item \textbf{Chercher l'intersection} entre le géotherme et le solidus : c'est le point où la roche atteint les conditions de début de fusion. Noter la pression (donc la profondeur) correspondante.
\item \textbf{Conclure sur le mécanisme} : si le géotherme croise le solidus anhydre par remontée de la roche, c'est une fusion par \textbf{décompression} ; si le géotherme ne croise que le solidus hydraté, c'est une fusion par \textbf{hydratation} (apport d'eau).
\end{enumerate}
\end{methode}

\courssub{La fusion par décompression : le magma des dorsales}

Sous une dorsale océanique, la divergence des plaques crée un vide qui provoque la \textbf{remontée du manteau} (péridotites chaudes de l'asthénosphère). Or cette remontée est quasi \textbf{adiabatique} : la roche monte vite et n'a pas le temps de perdre sa chaleur. Sa température reste donc presque constante tandis que la \textbf{pression diminue fortement}. Sur le diagramme, le point représentatif de la péridotite se déplace \emph{vers le haut} (décompression) à température presque constante, et finit par \textbf{croiser le solidus anhydre} (point A de la figure) : la roche commence à fondre \textbf{sans aucun apport de chaleur}. C'est la \cle{fusion partielle par décompression}.

\begin{exemplebox}[Des ordres de grandeur à connaître pour le Bac]
Sous une dorsale, la péridotite asthénosphérique à environ \SI{1300}{\celsius} commence à fondre lorsqu'elle atteint une profondeur de l'ordre de \textbf{60 à 80 km} (pression de 2 à 2,5 GPa). Le \textbf{taux de fusion partielle} est de l'ordre de \textbf{10 à 20 \%} : seule cette fraction de la roche fond. Le liquide obtenu, moins dense que la roche solide restante, s'extrait, s'accumule et remonte : c'est un \textbf{magma basaltique}, qui alimente le volcanisme des dorsales et forme les gabbros et les basaltes de la croûte océanique.
\end{exemplebox}

\courssub{La fusion par hydratation : le magma des zones de subduction}

Dans une zone de subduction, la plaque océanique \textbf{froide} plonge dans le manteau : le géotherme y est \textbf{anormalement bas}. Un simple raisonnement sur le diagramme montre que ce géotherme froid \textbf{ne croise jamais le solidus anhydre} : sans autre facteur, le manteau au-dessus de la plaque plongeante ne devrait \emph{pas} fondre. Pourtant, les zones de subduction sont bordées d'arcs volcaniques actifs (Andes, Japon, et chez nous la ligne volcanique du Cameroun s'inscrit dans un contexte de dynamique lithosphérique comparable, bien que souvent rattachée à un volcanisme intraplaque). Il faut donc un facteur supplémentaire : \textbf{l'eau}.

La plaque plongeante transporte des sédiments hydratés et des minéraux hydratés (amphiboles, serpentines, loiwsite). En s'enfonçant, elle subit une augmentation de pression et de température qui \textbf{déshydrate} ces minéraux (métamorphisme de la plaque) : l'eau libérée s'échappe et \textbf{monte} dans le manteau sus-jacent (le « coin mantellique »), situé au-dessus de la plaque. Or, d'après la propriété expérimentale admise, l'eau \textbf{abaisse le solidus} de la péridotite de plusieurs centaines de degrés. Le géotherme froid de subduction croise alors le \textbf{solidus hydraté} (point B de la figure) : la péridotite du coin mantellique fond \textbf{par hydratation}, à une température bien inférieure à celle qui serait nécessaire à sec. C'est la \cle{fusion partielle par hydratation}.

Le magma ainsi produit, initialement basaltique, s'enrichit en silice au cours de sa remontée à travers la lithosphère continentale épaisse (fusion de matériaux de la croûte, cristallisation fractionnée, mélange magmatique), pour donner des laves de type \textbf{andésitique}, visqueuses, responsables d'un volcanisme \textbf{explosif} caractéristique des zones de subduction.

\begin{definition}[Fusion partielle]
La \cle{fusion partielle} est la fusion d'une \textbf{fraction seulement} d'une roche (quelques pourcents à environ 20 \%). Parce que les minéraux ne fondent pas tous à la même température, les premiers liquides formés ont une composition \textbf{différente} de la roche mère : la fusion partielle d'une péridotite produit un magma \textbf{basaltique} (plus riche en silice et en éléments incompatibles que la péridotite). Le taux de fusion partielle est admis : il est de \textbf{quelques pourcents} à une vingtaine de pourcents selon les contextes.
\end{definition}

\begin{popfigure}
\begin{center}
\begin{tikzpicture}[font=\small, scale=0.95, transform shape]
% ===== Panneau gauche : dorsale =====
\begin{scope}
% Manteau (asthénosphère)
\fill[poporangeL] (0,-0.5) rectangle (6.5,-5);
% Lithosphère océanique de part et d'autre (fine)
\fill[popmuted!70] (0,-0.5) rectangle (2.5,-1.3);
\fill[popmuted!70] (4.0,-0.5) rectangle (6.5,-1.3);
% Rift / axe (zone d'épanchement)
\fill[popink!80] (3.25,-5) -- (3.55,-5) -- (3.45,-0.5) -- (3.35,-0.5) -- cycle;
% Flèches de divergence des plaques
\draw[-{Stealth[length=3mm]}, very thick, popdark] (2.3,-0.9) -- (1.0,-0.9) node[left, popdark, font=\footnotesize] {Divergence};
\draw[-{Stealth[length=3mm]}, very thick, popdark] (4.2,-0.9) -- (5.5,-0.9);
% Remontée péridotite (flèches courbes)
\draw[-{Stealth[length=3mm]}, very thick, popink] (2.6,-4.5) .. controls (2.8,-3.0) .. (3.25,-1.8);
\draw[-{Stealth[length=3mm]}, very thick, popink] (3.9,-4.5) .. controls (3.7,-3.0) .. (3.35,-1.8);
% Zone de fusion partielle (zone hachurée/colorée)
\fill[popink!20, draw=popink, thick] (2.5,-3.0) ellipse (1.0 and 0.6);
\node[popink, font=\footnotesize, align=center] at (1.0,-3.0) {Zone de fusion\\partielle par\\décompression};
% Légendes
\node[font=\bfseries, popdark] at (3.25,0.4) {Zone d'accrétion (Dorsale)};
\node[font=\footnotesize, popdark] at (5.7,-0.2) {Croûte océanique};
\node[font=\footnotesize, popink] at (3.25,-0.2) {Magma basaltique $\uparrow$};
\node[font=\footnotesize] at (3.25,-4.8) {Manteau (péridotite) chaud qui remonte};
\end{scope}

% ===== Panneau droit : subduction =====
\begin{scope}[xshift=8.5cm]
% Manteau (coin mantellique)
\fill[popgreenL] (0,-0.5) rectangle (6.5,-5);
% Plaque plongeante (océanique, froide)
\fill[popmuted!70, rotate around={-30:(0.5,-0.5)}] (0.5,-0.5) rectangle (6.0,-1.4);
% Croûte continentale (sus-jacente)
\fill[poppurpleL] (3.0,-0.5) rectangle (6.5,0.5);
% Volcan (cône)
\fill[popink!70] (4.5,0.5) -- (5.5,0.5) -- (5.0,1.6) -- cycle;
\draw[popink, thick] (5.0,1.6) -- (5.0,-1.5) node[below, popink, font=\footnotesize] {Conduit magmatique};
% Flèches d'eau libérée par la plaque (montent vers le coin mantellique)
\draw[-{Stealth[length=3mm]}, very thick, popblue] (2.5,-3.0) -- (3.5,-2.0);
\draw[-{Stealth[length=3mm]}, very thick, popblue] (3.0,-3.5) -- (4.0,-2.5);
\node[popblue, font=\footnotesize, align=center] at (1.8,-2.2) {Eau libérée\\par déshydratation};
% Zone de fusion partielle dans le coin mantellique
\fill[popink!20, draw=popink, thick] (4.2,-2.0) ellipse (0.9 and 0.5);
\node[popink, font=\footnotesize, align=center] at (5.8,-2.5) {Zone de fusion\\partielle par\\hydratation};
% Flèche subduction
\draw[-{Stealth[length=3mm]}, very thick, popdark, rotate around={-30:(5.5,-1.0)}] (5.5,-1.0) -- (6.5,-1.0) node[right, popdark, font=\footnotesize] {Subduction};
% Légendes
\node[font=\bfseries, popdark] at (3.25,1.5) {Zone de subduction};
\node[font=\footnotesize, popdark, rotate=-30] at (2.0,-1.5) {Plaque plongeante (froide)};
\node[font=\footnotesize] at (4.8,-0.2) {Croûte continentale};
\node[font=\footnotesize, popink] at (5.0,1.9) {Magma andésitique};
\end{scope}
\end{tikzpicture}
\end{center}
\legende{Comparaison des deux mécanismes de production magmatique. \textbf{À gauche} : sous une dorsale, la péridotite chaude remonte (flèches roses) et fond par décompression ; le magma basaltique s'épanche au rift. \textbf{À droite} : la plaque plongeante se déshydrate ; l'eau (flèches bleues) monte dans le manteau sus-jacent (coin mantellique) et y déclenche une fusion partielle par hydratation ; le magma andésitique alimente un volcan explosif.}
\end{popfigure}

\courssub{Bilan et exercice résolu}

\begin{aretenir}[Deux mécanismes, deux magmas]
\begin{itemize}
\item \textbf{Zone d'accrétion (dorsale)} : remontée du manteau $\Rightarrow$ \textbf{décompression} $\Rightarrow$ le géotherme chaud croise le \textbf{solidus anhydre} vers 60--80 km de profondeur (2--2,5 GPa) $\Rightarrow$ fusion partielle (10--20 \%) $\Rightarrow$ \textbf{magma basaltique} (volcanisme effusif, formation de la croûte océanique).
\item \textbf{Zone de subduction} : déshydratation de la plaque plongeante $\Rightarrow$ \textbf{apport d'eau} qui abaisse le solidus $\Rightarrow$ le géotherme froid croise le \textbf{solidus hydraté} $\Rightarrow$ fusion partielle $\Rightarrow$ magma basaltique primaire $\Rightarrow$ évolution vers \textbf{magma andésitique} (volcanisme explosif, arcs insulaires et cordillères).
\item Dans les deux cas, le manteau reste \textbf{solide} : seule une faible fraction de la péridotite fond, localement. La composition du magma diffère de celle de la roche mère parce que la fusion est \textbf{partielle}.
\end{itemize}
\end{aretenir}

\begin{exoresolu}[Exploiter un diagramme pression--température]
\textbf{Énoncé.} Un document donne, dans un diagramme (température en abscisse, pression en ordonnée), le solidus de la péridotite anhydre, le solidus de la péridotite hydratée et deux géothermes G1 et G2. G1 croise le solidus anhydre à 2,2 GPa ; G2 ne croise pas le solidus anhydre mais croise le solidus hydraté à 2,8 GPa. 1) Attribuer chaque géotherme à un contexte géodynamique en justifiant. 2) Préciser pour chaque contexte le mécanisme de fusion et la nature du magma produit. 3) Calculer la profondeur approximative du début de fusion dans le contexte de G1 (1 GPa $\approx$ 33 km).
\tcblower
\textbf{Solution.} 1) G1 croise le solidus anhydre : les conditions de fusion sont atteintes \emph{sans eau}, donc par simple décompression d'un manteau chaud qui remonte : G1 est le géotherme d'une \textbf{dorsale} (géotherme chaud). G2 ne croise pas le solidus anhydre (manteau trop froid) mais croise le solidus hydraté : la fusion n'est possible qu'en présence d'eau, ce qui caractérise le manteau situé au-dessus d'une \textbf{plaque en subduction} (géotherme froid). 2) Dorsale : fusion partielle \textbf{par décompression} $\Rightarrow$ magma \textbf{basaltique}. Subduction : fusion partielle \textbf{par hydratation} $\Rightarrow$ magma \textbf{andésitique}. 3) Profondeur : $2{,}2 \times 33 \approx 73$ km, valeur cohérente avec les 60--80 km admis sous les dorsales.
\end{exoresolu}

\begin{attention}[Pièges de rédaction au Baccalauréat]
\begin{itemize}
\item Ne jamais écrire que « la chaleur fait fondre le manteau sous la dorsale » : la température de la péridotite \textbf{n'augmente pas} lors de sa remontée adiabatique ; c'est la \textbf{baisse de pression} qui déclenche la fusion.
\item Ne pas confondre les deux solidus : dans une subduction, c'est l'\textbf{eau} qui permet la fusion, pas une élévation de température.
\item Ne pas dire que la plaque plongeante fond : c'est le \textbf{manteau au-dessus de la plaque} (coin mantellique) qui fond grâce à l'eau libérée par la plaque.
\item Préciser toujours « fusion \textbf{partielle} » et donner un ordre de grandeur du taux (quelques pourcents à 20 \%).
\item Ne pas oublier l'unité de pression (GPa) ou de profondeur (km) lors de la lecture du diagramme.
\end{itemize}
\end{attention}

\begin{savaistu}[Le volcanisme camerounais et la prévention des risques]
La \textbf{ligne volcanique du Cameroun}, qui s'étend du golfe de Guinée (île de Bioko) au mont Cameroun (\SI{4095}{m}, plus haut sommet d'Afrique de l'Ouest) puis vers l'Adamaoua (lac Nyos, lac Monoun), illustre que la production de magma n'est pas limitée aux dorsales et subductions : il existe aussi un volcanisme \textbf{intraplaque} lié à des remontées localisées de manteau chaud (points chauds, zones de distension lithosphérique). Le mont Cameroun, volcan actif dont les éruptions récentes (1999, 2000) ont menacé les plantations de la CDC et les villages de la côte (Buea, Limbe), rappelle que comprendre la naissance des magmas est une étape indispensable de la \textbf{prévention des risques volcaniques}, que nous étudierons dans la dernière partie de cette leçon. La catastrophe limnique du lac Nyos (1986), due à une éruption gazeuse (\ce{CO2}) d'origine magmatique profonde, souligne aussi l'importance de la surveillance géochimique.
\end{savaistu}

En résumé, la pétrologie expérimentale fournit la clé du magmatisme terrestre : la péridotite du manteau, bien que solide, fond partiellement et localement chaque fois que les conditions de pression, de température et d'hydratation franchissent son solidus. Deux mécanismes dominent — la \textbf{décompression} sous les dorsales et l'\textbf{hydratation} au-dessus des zones de subduction — qui expliquent respectivement le magmatisme basaltique des océans et le magmatisme andésitique des arcs volcaniques. Nous pouvons maintenant aborder les autres frontières de plaques : les zones de convergence et de coulissage.

\coursec{Convergence et coulissage des plaques}

Dans la section précédente, nous avons vu comment les données expérimentales de fusion des péridotites expliquent la naissance des magmas : fusion par décompression sous les dorsales (zones d'accrétion) et fusion hydratée au-dessus des zones de subduction. Mais une question demeure : \emph{pourquoi} la lithosphère océanique plonge-t-elle dans l'asthénosphère, et que deviennent les plaques qui se rapprochent ou qui glissent l'une contre l'autre ? C'est l'objet de cette section : décrire les \cle{mouvements de convergence} et les \cle{mouvements de coulissage} des plaques, et identifier leurs facteurs clés.

\courssub{Rappel : une lithosphère découpée en plaques mobiles}

\begin{aretenir}[Le modèle de la tectonique des plaques (savoir admis)]
La \cle{lithosphère} (croûte + manteau supérieur rigide, épaisseur d'environ \SI{100}{km} sous les océans) est découpée en une douzaine de \cle{plaques} rigides qui se déplacent, à des vitesses de quelques centimètres par an (\SI{1}{cm/an} à \SI{10}{cm/an}), sur l'\cle{asthénosphère}, zone du manteau supérieur solide mais \emph{ductile} (déformable) présentant un faible degré de fusion partielle. Trois types de limites séparent les plaques :
\begin{itemize}
\item les limites \textbf{divergentes} (dorsales, rifts) : les plaques s'écartent, il y a création de croûte océanique ;
\item les limites \textbf{convergentes} : les plaques se rapprochent (subduction ou collision) ;
\item les limites \textbf{coulissantes} (failles transformantes) : les plaques glissent horizontalement l'une le long de l'autre.
\end{itemize}
\end{aretenir}

\begin{popfigure}
\begin{center}
\begin{tikzpicture}[scale=0.92, every node/.style={font=\small}]
% Afrique
\draw[fill=popgreenL, draw=popdark, thick] (0.5,2.2) .. controls (1.2,3.4) and (2.2,3.8) .. (2.8,3.2)
 .. controls (3.4,2.7) and (3.1,1.6) .. (2.5,0.9) .. controls (2.0,0.3) and (1.2,0.6) .. (0.9,1.3)
 .. controls (0.6,1.8) and (0.3,1.9) .. (0.5,2.2) -- cycle;
\node at (1.8,2.1) {\textbf{Africaine}};
% Amérique du Sud
\draw[fill=poporangeL, draw=popdark, thick] (-4.6,1.6) .. controls (-4.0,2.6) and (-3.2,2.7) .. (-2.9,2.0)
 .. controls (-2.6,1.3) and (-3.0,0.4) .. (-3.5,-0.2) .. controls (-3.9,-0.7) and (-4.5,-0.2) .. (-4.7,0.6)
 .. controls (-4.8,1.1) and (-4.8,1.3) .. (-4.6,1.6) -- cycle;
\node at (-3.8,1.2) {\textbf{Amérique}};
% Eurasie
\draw[fill=popblueL, draw=popdark, thick] (0.6,3.9) .. controls (2.0,4.6) and (4.6,4.6) .. (5.6,3.9)
 .. controls (5.2,3.4) and (3.6,3.5) .. (2.9,3.5) .. controls (1.8,3.5) and (0.9,3.5) .. (0.6,3.9) -- cycle;
\node at (3.2,4.05) {\textbf{Eurasiatique}};
% Pacifique
\draw[fill=poppurpleL, draw=popdark, thick] (-6.3,3.6) .. controls (-5.4,4.3) and (-4.6,4.2) .. (-4.4,3.5)
 .. controls (-4.9,3.0) and (-5.9,3.0) .. (-6.3,3.6) -- cycle;
\node at (-5.4,3.6) {\textbf{Pacifique}};
% Indo-australienne
\draw[fill=popgoldL, draw=popdark, thick] (4.3,0.4) .. controls (5.0,1.1) and (5.9,1.0) .. (6.1,0.3)
 .. controls (5.6,-0.3) and (4.7,-0.3) .. (4.3,0.4) -- cycle;
\node at (5.2,0.4) {\textbf{Indo-austr.}};
% Dorsale médio-atlantique
\draw[very thick, poporange, decorate, decoration={zigzag, segment length=4, amplitude=1.2}] (-1.6,4.2) -- (-1.9,2.6) -- (-1.7,1.0) -- (-2.0,-0.4);
\node[poporange, font=\footnotesize, rotate=80] at (-2.35,2.0) {dorsale};
% Flèches de divergence
\draw[->, very thick, poporange] (-1.9,2.4) -- (-2.9,2.4);
\draw[->, very thick, poporange] (-1.6,2.4) -- (-0.6,2.4);
% Convergence Himalaya
\draw[->, very thick, poppink] (4.9,1.5) -- (4.3,2.6);
\draw[->, very thick, poppink] (3.4,3.3) -- (4.0,2.7);
\node[poppink, font=\footnotesize] at (4.6,2.2) {collision};
% Subduction Andes
\draw[->, very thick, popblue] (-5.6,0.6) -- (-4.9,0.9);
\node[popblue, font=\footnotesize] at (-5.6,0.2) {subduction};
% Coulissage San Andreas
\draw[very thick, poppurple] (-5.9,2.6) -- (-4.7,3.1);
\draw[->, thick, poppurple] (-5.8,2.85) -- (-5.3,3.05);
\draw[->, thick, poppurple] (-5.1,2.55) -- (-5.6,2.4);
\node[poppurple, font=\footnotesize] at (-6.1,2.2) {coulissage};
\end{tikzpicture}
\end{center}
\legende{Carte très simplifiée de quelques grandes plaques lithosphériques et de leurs mouvements relatifs. Zigzag orange : dorsale médio-atlantique (divergence) ; flèches roses face à face : convergence de collision (Himalaya) ; flèche bleue : subduction (côte ouest de l'Amérique du Sud) ; doubles flèches violettes opposées : coulissage (faille de San Andreas).}
\end{popfigure}

\courssub{Les mouvements de convergence : deux plaques se rapprochent}

Lorsque deux plaques se rapprochent, la surface de lithosphère située entre elles doit \emph{disparaître} ou se \emph{déformer}. L'issue dépend de la \cle{nature des lithosphères en présence}, et plus précisément de leur \cle{densité} :
\begin{itemize}
\item la lithosphère \textbf{océanique} est mince (environ \SI{7}{km} de croûte basaltique sur un manteau lithosphérique) et \textbf{dense} ($d \approx \num{3,0}$ à $\num{3,3}$) ;
\item la lithosphère \textbf{continentale} est épaisse (croûte granitique de \SI{30}{km} à \SI{40}{km}) et \textbf{peu dense} ($d \approx \num{2,7}$ pour la croûte) : elle ne peut pas plonger profondément dans l'asthénosphère.
\end{itemize}
Deux cas se présentent donc : la \textbf{subduction} et la \textbf{collision}.

\textbf{La subduction : la plongée de la lithosphère océanique}

\begin{definition}[Subduction]
La \cle{subduction} est l'enfoncement d'une lithosphère océanique, dense, sous une autre lithosphère (océanique ou continentale), dans l'asthénosphère. La zone de plongée est marquée en surface par une \cle{fosse océanique} (dépression pouvant atteindre \SI{11000}{m} de profondeur, comme la fosse des Mariannes).
\end{definition}

\begin{experience}[Observation : la distribution des foyers sismiques près des Andes]
\textbf{Document.} On reporte sur une coupe perpendiculaire à la cordillère des Andes la profondeur des foyers des séismes enregistrés pendant plusieurs années, en fonction de leur distance à la fosse du Pérou-Chili.

\textbf{Observations.} Les foyers ne sont pas répartis au hasard : ils s'alignent selon un \emph{plan incliné} d'environ 30 à 45°, partant de la fosse (séismes superficiels, profondeur $< \SI{70}{km}$) et s'enfonçant sous le continent (séismes intermédiaires puis profonds, jusqu'à environ \SI{700}{km}).

\textbf{Interprétation.} Les séismes traduisent des ruptures brutales de roches \emph{rigides}. Leur alignement incliné matérialise donc une plaque rigide — la lithosphère océanique — qui s'enfonce sous la lithosphère continentale. C'est la preuve directe de la subduction.
\end{experience}

\begin{definition}[Plan de Bénioff]
Le \cle{plan de Bénioff} (ou plan de Wadati-Bénioff) est la surface inclinée le long de laquelle s'alignent les foyers des séismes, de plus en plus profonds à mesure qu'on s'éloigne de la fosse sous la plaque chevauchante. Il matérialise la position de la plaque plongeante.
\end{definition}

\begin{popfigure}
\begin{center}
\begin{tikzpicture}[scale=1.0, every node/.style={font=\small}]
% Asthénosphère
\fill[poporangeL] (-6.5,-3.4) rectangle (6.5,-0.9);
\node[popdark] at (0,-3.1) {\textbf{Asthénosphère} (ductile, partiellement fondue)};
% Lithosphère océanique
\fill[popblueL] (-6.5,-0.9) rectangle (-0.2,0);
\draw[popdark, thick] (-6.5,0) -- (-0.2,0);
\node at (-3.4,0.35) {\textbf{Lithosphère océanique} (dense)};
\node[font=\footnotesize] at (-3.4,-0.5) {croûte basaltique};
% Plaque plongeante
\fill[popblueL] (-0.2,-0.9) -- (4.6,-3.4) -- (5.4,-3.4) -- (0.6,-0.9) -- cycle;
% Lithosphère continentale
\fill[popgreenL] (0.6,-0.9) rectangle (6.5,0.3);
\draw[popdark, thick] (0.6,0.3) -- (6.5,0.3);
\node at (3.9,0.65) {\textbf{Lithosphère continentale}};
% Fosse
\draw[popdark, very thick] (-0.2,0) -- (0.2,-0.25) -- (0.6,0.3);
\node[font=\footnotesize, popblue] at (-0.9,-0.55) {fosse};
% Volcan
\draw[fill=popmuted!60, draw=popdark] (2.6,0.3) -- (3.2,1.3) -- (3.8,0.3) -- cycle;
\draw[fill=poppink!70] (3.1,1.3) -- (3.2,1.5) -- (3.3,1.3) -- cycle;
\node[font=\footnotesize] at (4.6,1.25) {volcan explosif (andésitique)};
% Cheminée magmatique
\draw[poppink, very thick, ->] (3.0,-1.7) -- (3.2,1.2);
\node[poppink, font=\footnotesize, rotate=72] at (2.55,-0.4) {magma};
% Plan de Bénioff : foyers
\foreach \x/\y in {0.1/-0.5, 0.7/-0.9, 1.3/-1.2, 1.9/-1.6, 2.5/-2.0, 3.1/-2.4, 3.7/-2.7, 4.3/-3.0}{
  \node[popink, font=\footnotesize] at (\x,\y) {$\times$};}
\draw[popink, dashed, thick] (-0.1,-0.35) -- (4.7,-3.2);
\node[popink, font=\footnotesize, rotate=-32] at (2.0,-2.35) {plan de Bénioff};
% Flèches de mouvement
\draw[->, very thick, popblue] (-4.5,-1.5) -- (-2.5,-1.5);
\draw[->, very thick, popgreen!60!black] (5.5,-1.5) -- (4.3,-1.5);
% Échelle profondeur
\draw[<->, popmuted] (-6.2,0) -- (-6.2,-3.4) node[midway, left, font=\footnotesize] {profondeur};
\end{tikzpicture}
\end{center}
\legende{Coupe schématique d'une zone de subduction (type cordillère des Andes). La lithosphère océanique dense plonge sous la lithosphère continentale ; les croix noires matérialisent les foyers sismiques alignés selon le plan de Bénioff ; la fusion hydratée des péridotites au-dessus de la plaque plongeante alimente un volcanisme explosif andésitique.}
\end{popfigure}

\begin{popfigure}
\begin{center}
\begin{tikzpicture}
\begin{axis}[
  width=0.85\linewidth, height=6.2cm,
  xlabel={Distance à la fosse (km)},
  ylabel={Profondeur du foyer (km)},
  xmin=0, xmax=700, ymin=-700, ymax=0,
  y dir=reverse,
  grid=major, grid style={popline!50},
  xtick={0,100,200,300,400,500,600,700},
  ytick={0,100,200,300,400,500,600,700},
  legend style={at={(0.03,0.97)}, anchor=north west, font=\footnotesize},
]
\addplot[only marks, mark=*, mark size=1.6pt, popink] coordinates {
 (20,15) (60,30) (110,55) (150,80) (200,110) (240,140) (290,175)
 (330,210) (380,255) (420,300) (470,350) (510,400) (560,455) (610,520) (660,590)
};
\addlegendentry{Foyers sismiques enregistrés}
\addplot[popcurve, poporange, domain=0:700, samples=50] {0.9*x};
\addlegendentry{Alignement du plan de Bénioff}
\end{axis}
\end{tikzpicture}
\end{center}
\legende{Profondeur des foyers sismiques en fonction de la distance à la fosse (données fictives représentatives d'une subduction de type Andes). L'alignement croissant des foyers matérialise le plan de Bénioff, c'est-à-dire la plaque plongeante.}
\end{popfigure}

\textbf{Facteurs clés de la subduction : densité et pesanteur}

Pourquoi la plaque plonge-t-elle ? Deux facteurs conjuguent leurs effets :
\begin{enumerate}
\item La \cle{densité de la lithosphère océanique} : en s'éloignant de la dorsale, la lithosphère \emph{refroidit}, s'épaissit et se \emph{densifie}. Une lithosphère jeune et chaude (près de la dorsale) est peu dense et « flotte » haut sur l'asthénosphère ; une lithosphère âgée (plus de 50 à 100 Ma), froide, devient \emph{plus dense que l'asthénosphère sous-jacente}.
\item La \cle{pesanteur} : dès que la densité de la plaque dépasse celle de l'asthénosphère, son propre poids l'entraîne vers le bas : on parle d'\emph{aspiration} de la plaque (\emph{slab pull} des auteurs anglo-saxons). C'est le moteur principal du mouvement des plaques.
\end{enumerate}

\begin{propriete}[Condition de l'enfoncement d'une plaque]
C'est la \cle{différence de densité} entre la lithosphère océanique et l'asthénosphère qui permet l'enfoncement de la plaque plongeante : la subduction n'est possible que si
\[ d_{\text{lithosphère océanique}} > d_{\text{asthénosphère}}. \]
Comme la densité de la lithosphère augmente avec son âge (refroidissement progressif), ce sont les lithosphères océaniques \emph{anciennes et froides} qui subductent. À l'inverse, la lithosphère continentale, trop peu dense, ne peut pas s'enfoncer durablement.
\end{propriete}

\begin{exemplebox}[La subduction et ses marqueurs : la ceinture de feu du Pacifique]
Autour de l'océan Pacifique, la plaque Pacifique et la plaque de Nazca plongent sous les plaques voisines. On y observe systématiquement l'association de trois marqueurs : une \textbf{fosse océanique} profonde, des \textbf{séismes} dont les foyers s'approfondissent sous la plaque chevauchante (plan de Bénioff), et un \textbf{volcanisme explosif} de type andésitique (Andes, Japon, Indonésie), alimenté par la fusion hydratée des péridotites du manteau au-dessus de la plaque plongeante (voir section précédente).
\end{exemplebox}

\begin{attention}[Erreur fréquente : « la plaque continentale flotte sur la plaque océanique »]
Il est \textbf{faux} de dire que la plaque continentale « flotte » sur la plaque océanique ou que celle-ci « passe par-dessous » comme sur un tapis roulant. Les deux plaques reposent \emph{côte à côte} sur l'asthénosphère ; c'est la \textbf{lithosphère océanique qui plonge dans l'asthénosphère} parce qu'elle est devenue plus dense qu'elle. De même, évite d'écrire que « la croûte continentale plonge » : c'est la \emph{lithosphère} (croûte + manteau lithosphérique) qui est en mouvement.
\end{attention}

\textbf{La collision : quand deux continents se rencontrent}

Que se passe-t-il lorsque la convergence rapproche deux lithosphères \emph{continentales} ? La subduction d'une plaque océanique peut d'abord consumer tout l'océan situé entre deux continents ; les continents finissent alors par se rencontrer. Or la lithosphère continentale, peu dense ($d \approx \num{2,7}$), \emph{ne peut pas plonger} profondément dans l'asthénosphère.

\begin{definition}[Collision]
La \cle{collision} est la convergence de deux lithosphères continentales, trop peu denses pour subducter. La compression raccourcit et épaissit les croûtes : les roches se \textbf{plissent} et se \textbf{faillent} (chevauchements), formant une \cle{chaîne de montagnes} de collision.
\end{definition}

\begin{exemplebox}[L'Himalaya, chaîne de collision en activité]
Il y a environ 50 Ma, l'océan Téthys séparait l'Inde de l'Asie. La subduction de la lithosphère océanique de la plaque indienne sous l'Asie a fermé cet océan ; puis le continent indien est entré en collision avec le continent eurasiatique. La compression, qui se poursuit encore aujourd'hui (l'Inde avance d'environ \SI{5}{cm/an} vers le nord), a produit l'\textbf{Himalaya}, dont les sommets s'élèvent encore d'environ \SI{1}{cm/an} (Everest : \SI{8849}{m}). On n'y observe \emph{ni fosse océanique, ni volcanisme andésitique}, mais des séismes superficiels à intermédiaires et d'épaissis empilements de chevauchements.
\end{exemplebox}

\begin{popfigure}
\begin{center}
\begin{tikzpicture}[scale=1.0, every node/.style={font=\small}]
% Asthénosphère
\fill[poporangeL] (-6.5,-3.2) rectangle (6.5,-1.3);
\node[popdark] at (0,-2.9) {\textbf{Asthénosphère}};
% Continent gauche
\fill[popgreenL] (-6.5,-1.3) rectangle (-0.3,0.2);
\draw[popdark, thick] (-6.5,0.2) -- (-0.3,0.2);
\node at (-3.4,0.55) {\textbf{Lithosphère continentale 1}};
% Continent droit
\fill[popblueL] (0.3,-1.3) rectangle (6.5,0.2);
\draw[popdark, thick] (0.3,0.2) -- (6.5,0.2);
\node at (3.6,0.55) {\textbf{Lithosphère continentale 2}};
% Chaîne de montagnes : plis
\draw[popdark, very thick, fill=popgoldL]
 (-0.3,0.2) .. controls (-0.9,1.6) and (-0.3,1.7) .. (0,0.9)
 .. controls (0.3,1.9) and (0.9,1.8) .. (0.6,0.7)
 .. controls (0.5,0.4) and (0.4,0.3) .. (0.3,0.2) -- cycle;
\node[font=\footnotesize] at (0,2.15) {chaîne de montagnes (plis, chevauchements)};
% Racine crustale
\fill[popgreen!50!popblue!50!white] (-0.9,-1.3) .. controls (-0.5,-2.4) and (0.5,-2.4) .. (0.9,-1.3) -- cycle;
\node[font=\footnotesize, popdark] at (2.6,-2.0) {racine crustale (croûte épaissie)};
\draw[->, popdark, thin] (1.7,-1.95) -- (0.5,-1.8);
% Flèches de convergence
\draw[->, very thick, poppink] (-4.6,-0.6) -- (-2.6,-0.6);
\draw[->, very thick, poppink] (4.6,-0.6) -- (2.6,-0.6);
% Séismes superficiels
\foreach \x/\y in {-0.5/-0.5, 0.2/-0.7, 0.6/-0.4}{
  \node[popink, font=\footnotesize] at (\x,\y) {$\times$};}
\node[popink, font=\footnotesize] at (-2.2,-0.95) {séismes superficiels};
\end{tikzpicture}
\end{center}
\legende{Coupe schématique d'une zone de collision (type Himalaya). Aucune plaque ne plonge : la compression raccourcit la lithosphère, provoque plissements et chevauchements, élève une chaîne de montagnes et crée une racine crustale profonde. Il n'y a ni fosse, ni volcanisme andésitique ; les séismes restent superficiels à intermédiaires.}
\end{popfigure}

\courssub{Les mouvements de coulissage : les failles transformantes}

\begin{definition}[Faille transformante]
Une \cle{faille transformante} est une faille verticale le long de laquelle deux plaques \textbf{coulissent horizontalement} l'une par rapport à l'autre. Il n'y a \emph{ni création ni destruction de croûte} : la surface totale de lithosphère est conservée. Les failles transformantes relient souvent deux segments de dorsale décalés.
\end{definition}

\begin{experience}[Observation : le décalage des structures le long de la faille de San Andreas]
\textbf{Document.} En Californie, la faille de San Andreas sépare la plaque Pacifique (au sud-ouest) de la plaque Amérique du Nord (au nord-est). On dispose de photographies aériennes et de mesures répétées depuis plus d'un siècle.

\textbf{Observations.} Des structures linéaires continues autrefois — ruisseaux, clôtures, routes, lignes de chemin de fer — sont aujourd'hui \emph{décalées horizontalement} de plusieurs mètres, parfois de plusieurs centaines de mètres pour les plus anciennes. La région est secouée par des séismes fréquents, dont les foyers sont toujours \emph{superficiels} (moins de \SI{20}{km} de profondeur). Aucune fosse, aucun volcan n'y est associé.

\textbf{Interprétation.} Le décalage horizontal progressif des repères prouve un \emph{coulissage} : les deux plaques glissent l'une le long de l'autre (environ \SI{5}{cm/an}). Les contraintes s'accumulent dans les roches rigides superficielles, puis se relâchent brutalement lors de ruptures : ce sont les séismes superficiels (ex. séisme de San Francisco, 1906).
\end{experience}

\begin{popfigure}
\begin{center}
\begin{tikzpicture}[scale=1.0, every node/.style={font=\small}]
% Bloc gauche
\fill[popblueL] (-5.5,-2.2) rectangle (0,2.2);
% Bloc droit
\fill[popgreenL] (0,-2.2) rectangle (5.5,2.2);
% Faille
\draw[popink, very thick] (0,-2.2) -- (0,2.2);
\node[popink, rotate=90] at (-0.28,0) {\footnotesize faille transformante};
% Ruisseau décalé
\draw[popblue, very thick] (-5.5,1.2) -- (0,1.2);
\draw[popblue, very thick] (0,-0.6) -- (5.5,-0.6);
\draw[popblue, ->, thick] (-4.6,1.2) -- (-3.8,1.2);
\draw[popblue, ->, thick] (3.8,-0.6) -- (4.6,-0.6);
\node[popblue, font=\footnotesize] at (-3.2,1.55) {ruisseau décalé};
\node[popblue, font=\footnotesize] at (3.2,-0.95) {ruisseau décalé};
% Accolade du décalage
\draw[<->, poppink, very thick] (0.35,1.2) -- (0.35,-0.6) node[midway, right, font=\footnotesize] {décalage};
% Flèches de coulissage
\draw[->, very thick, poppurple] (-3.5,1.9) -- (-1.5,1.9);
\draw[->, very thick, poppurple] (1.5,-1.9) -- (3.5,-1.9);
\node at (-3.2,-1.5) {\textbf{Plaque A}};
\node at (3.2,1.5) {\textbf{Plaque B}};
% Séismes
\foreach \p in {(0.15,0.6), (-0.15,-1.3), (0.12,1.7)}{
  \node[popink, font=\footnotesize] at \p {$\times$};}
\node[popink, font=\footnotesize, align=center] at (-3.0,0.2) {séismes\\superficiels};
\end{tikzpicture}
\end{center}
\legende{Vue en plan d'une faille transformante (type San Andreas). Les deux plaques coulissent horizontalement en sens inverse ; les structures linéaires (ici un ruisseau) sont décalées ; les séismes, fréquents, ont des foyers superficiels. Il n'y a ni création ni destruction de croûte.}
\end{popfigure}

\begin{attention}[Ne confonds pas transformante et subduction]
Le long d'une faille transformante, les séismes sont \textbf{superficiels} et il n'y a \textbf{ni volcanisme, ni fosse}. Si un document montre des foyers de plus en plus profonds en s'éloignant d'une fosse, ce n'est \emph{pas} un coulissage mais une \emph{subduction}. Retiens aussi que le coulissage conserve la surface de lithosphère : il n'y a ni accrétion ni disparition de croûte.
\end{attention}

\courssub{Bilan : les trois types de limites de plaques}

\begin{aretenir}[Tableau récapitulatif des limites de plaques]
\begin{center}
\begin{tabularx}{\linewidth}{|l|X|X|X|}
\hline
\rowcolor{popblueL} \textbf{Critère} & \textbf{Divergente} & \textbf{Convergente} & \textbf{Coulissante} \\
\hline
Mouvement relatif & écartement & rapprochement & glissement horizontal \\
\hline
Structure associée & dorsale, rift & fosse (subduction) ou chaîne de montagnes (collision) & faille transformante \\
\hline
Bilan de croûte & création (accrétion) & destruction (subduction) ou raccourcissement (collision) & ni création ni destruction \\
\hline
Séismes & superficiels & superficiels à profonds (plan de Bénioff) ; superficiels en collision & superficiels, fréquents \\
\hline
Volcanisme & basaltique, effusif & explosif, andésitique (subduction) ; absent (collision) & absent \\
\hline
Exemple & dorsale médio-atlantique & Andes (subduction), Himalaya (collision) & faille de San Andreas \\
\hline
\end{tabularx}
\end{center}
\end{aretenir}

\begin{methode}[Identifier un type de limite de plaque à partir de documents]
Face à un ensemble de documents (carte des séismes, coupe topographique, carte du volcanisme), procède par étapes :
\begin{enumerate}
\item \textbf{Cherche une fosse océanique} sur les profils topographiques et la répartition en \emph{profondeur} des foyers sismiques.
\item \textbf{Conclus selon les indices} :
\begin{itemize}
\item fosse + séismes de plus en plus profonds (plan de Bénioff) + volcanisme explosif $\Rightarrow$ \textbf{subduction} ;
\item chaîne de montagnes \emph{sans fosse}, séismes superficiels à intermédiaires, pas de volcanisme $\Rightarrow$ \textbf{collision} ;
\item décalage horizontal de structures + séismes superficiels uniquement, sans volcanisme $\Rightarrow$ \textbf{coulissage} ;
\item relief volcanique axial + séismes superficiels + roches basaltiques $\Rightarrow$ \textbf{divergence} (dorsale).
\end{itemize}
\item \textbf{Justifie toujours} ta réponse en citant précisément les indices du document (profondeurs en km, présence/absence de fosse, nature du volcanisme).
\end{enumerate}
\end{methode}

\begin{exoresolu}[Identifier la limite de plaque au large du Pérou]
\textbf{Énoncé.} Au large du Pérou, on observe : (1) une fosse océanique atteignant \SI{6000}{m} de profondeur ; (2) des séismes dont les foyers s'enfoncent jusqu'à \SI{600}{km} sous la cordillère des Andes en s'éloignant de la fosse ; (3) des volcans explosifs à laves andésitiques. Identifie le type de limite de plaques et explique le mécanisme à l'œuvre.
\tcblower
\textbf{Solution détaillée.}
\begin{itemize}
\item \emph{Analyse des indices.} La fosse (indice 1) traduit l'enfoncement d'une plaque ; l'approfondissement progressif des foyers sismiques sous le continent (indice 2) dessine un plan de Bénioff, qui matérialise une plaque plongeante ; le volcanisme andésitique explosif (indice 3) résulte de la fusion hydratée des péridotites du manteau au-dessus de cette plaque.
\item \emph{Conclusion.} Il s'agit d'une limite \textbf{convergente de type subduction} : la lithosphère océanique de la plaque de Nazca, âgée, froide et donc plus dense que l'asthénosphère, plonge sous la lithosphère continentale sud-américaine, entraînée par la pesanteur.
\end{itemize}
\end{exoresolu}

\begin{savaistu}[La Ligne du Cameroun, témoin des mouvements lithosphériques]
Le Cameroun possède un alignement volcanique remarquable, la \cle{Ligne du Cameroun}, qui s'étend sur plus de \SI{1000}{km}, depuis les îles du golfe de Guinée (Bioko, São Tomé) jusqu'aux monts Oku et Bamboutos, en passant par le \textbf{Mont Cameroun} (\SI{4095}{m}), volcan encore actif (dernières éruptions en 1999 et 2000). Cet alignement, à la fois océanique \emph{et} continental, est interprété comme la trace de mouvements lithosphériques anciens liés à l'ouverture de l'Atlantique Sud et à la remontée de matériel mantellique chaud sous la plaque africaine. C'est un excellent exemple à mobiliser au Baccalauréat pour montrer que le volcanisme n'est pas limité aux limites de plaques actuelles : il existe aussi un volcanisme \emph{intraplaque}, lié à des points chauds ou à d'anciennes fractures lithosphériques.
\end{savaistu}

\begin{exercice}[À toi de jouer]
En Californie, les mesures GPS montrent que deux points situés de part et d'autre de la faille de San Andreas s'éloignent l'un de l'autre \emph{parallèlement à la faille} à raison de \SI{5}{cm/an}. Les séismes enregistrés ont des foyers à moins de \SI{15}{km} de profondeur et aucun volcan n'est présent.
\begin{enumerate}
\item Identifie le type de limite de plaques en justifiant ta réponse.
\item Calcule le décalage horizontal accumulé en 200 ans.
\item Explique pourquoi les séismes y sont superficiels.
\end{enumerate}
\end{exercice}

\begin{corrige}[Exercice : la faille de San Andreas]
\begin{enumerate}
\item Limite \textbf{coulissante} (faille transformante) : déplacement horizontal parallèle à la faille, séismes superficiels, absence de volcanisme et de fosse.
\item Décalage : $d = \SI{5}{cm/an} \times 200\ \text{ans} = \SI{1000}{cm} = \SI{10}{m}$.
\item Le coulissage concerne la lithosphère rigide : les contraintes s'accumulent dans les roches cassantes proches de la surface ; en profondeur, l'asthénosphère ductile se déforme sans rupture brutale, donc sans séisme.
\end{enumerate}
\end{corrige}

\coursec{Les catastrophes naturelles liées aux mouvements des plaques}

\courssub{De la dynamique des plaques aux aléas géologiques : comprendre pour mieux prévenir}

Dans les sections précédentes, nous avons analysé la morphologie des fonds océaniques, le fonctionnement des rifts (expansion du plancher océanique) et les mécanismes de convergence et de coulissage des plaques lithosphériques. Nous avons vu que ces mouvements, lents à l'échelle humaine (quelques centimètres par an), engendrent d'immenses contraintes mécaniques au sein des roches. Lorsque ces contraintes dépassent la résistance des matériaux, l'énergie accumulée est libérée brutalement : c'est l'origine des \cle{séismes} et, par conséquent, des \cle{tsunamis} et de l'essentiel du \cle{volcanisme}. Les mouvements verticaux (soulèvement, subsidence) et le relief jeune créé par la tectonique (chaînes de montagnes, failles actives) préparent également le terrain pour les \cle{glissements de terrain} et les \cle{inondations}, souvent aggravés par l'action humaine. Cette section établit le lien direct entre la géodynamique interne et les risques majeurs auxquels sont confrontées les populations, y compris au Cameroun.

\begin{definition}[Aléa, vulnérabilité, risque et catastrophe naturelle]
Un \textbf{aléa naturel} est un phénomène géologique (séisme, éruption, glissement de terrain) ou climatique (cyclone, crue) potentiellement dangereux, caractérisé par sa localisation, son intensité, sa fréquence et sa probabilité d'occurrence.
La \textbf{vulnérabilité} désigne le degré de perte (humaine, matérielle, économique, environnementale) qu'une population, des infrastructures ou un écosystème subiraient s'ils étaient touchés par un aléa de donnée intensité. Elle dépend de la densité de population, de la qualité de l'urbanisation (normes parasismiques), de la pauvreté, de la déforestation, etc.
Le \textbf{risque} est la convolution de l'aléa et de la vulnérabilité : $Risque = Aléa \times Vulnérabilité$.
Une \textbf{catastrophe naturelle} survient lorsqu'un aléa frappe une zone vulnérable, causant des dommages que la communauté ne peut surmonter sans aide extérieure.
\end{definition}

\begin{attention}[Piège conceptuel : Aléa $\neq$ Risque $\neq$ Catastrophe]
Un séisme de magnitude 7 dans un désert inhabité est un \emph{aléa fort} mais un \emph{risque nul} et \emph{pas une catastrophe}. Le même séisme dans une mégapole mal construite est une \emph{catastrophe majeure}. Au Cameroun, le risque volcanique du Mont Cameroun est élevé non seulement par l'aléa (fréquence éruptive) mais surtout par la forte vulnérabilité (populations denses sur les flancs, ville de Limbé au pied). Réduire le risque, c'est agir sur la vulnérabilité (aménagement, construction, éducation) car on ne peut supprimer l'aléa tectonique.
\end{attention}

\courssub{Genèse des magmas : apports de la fusion expérimentale des péridotites}

La compréhension du volcanisme (et donc du risque volcanique) repose sur la connaissance de l'origine des magmas. Les expériences de fusion de la péridotite (roche du manteau supérieur) sous pression et température contrôlées, menées notamment par \textbf{Green et Ringwood (1967)} et précisées par \textbf{Jaques et Green (1980)}, fournissent les contraintes physiques essentielles.

\begin{experience}[Fusion de la péridotite : rôle de l'eau et de la pression]
\textbf{Dispositif :} Capsule en or ou platine contenant de la péridotite naturelle (ou synthétique) $\pm$ eau ($\ce{H2O}$), soumise à haute pression ($P$, via presse à pistons ou enclumes de diamant) et haute température ($T$), puis quenchée (refroidissement brutal). Analyse des phases présentes (microscope électronique, microsonde).
\textbf{Résultats clés (données chiffrées) :}
\begin{itemize}
    \item \textbf{Péridotite sèche (anhydre) :} Le \textbf{solidus} (début de fusion) suit une courbe $T(^\circ\mathrm{C}) \approx 1100 + 130 \times P(\mathrm{GPa})$ pour $P < 3$ GPa. Le \textbf{liquidus} (fusion totale) est $\sim 150^\circ\mathrm{C}$ plus haut. \textbf{La fusion partielle ($\sim 10-25\%$) produit un magma basaltique (MORB)}.
    \item \textbf{Péridotite hydratée (avec $\ce{H2O}$) :} L'eau abaisse drastiquement le solidus : $\Delta T \approx -200$ à $-300^\circ\mathrm{C}$ à $2-3$ GPa. Le solidus "humide" coupe le géotherme de subduction. \textbf{La fusion partielle ($\sim 5-15\%$) produit un magma andésitique/basaltique riche en eau}.
    \item \textbf{Effet de la décompression adiabatique :} Le manteau remonte (dorsale, point chaud) $\rightarrow$ $P \downarrow$ à $T$ quasi constante (adabatique) $\rightarrow$ croisement du solidus sec $\rightarrow$ \textbf{fusion par décompression} (zones d'accrétion, points chauds).
    \item \textbf{Effet du flux hydrique (subduction) :} La plaque subductée se déshydrate (minéraux hydratés $\rightarrow$ amphibole, lawsonite, chlorite $\rightarrow$ $\ce{H2O}$ libérée) $\rightarrow$ $\ce{H2O}$ remonte dans le coin de manteau sus-jacent $\rightarrow$ abaisse le solidus $\rightarrow$ \textbf{fusion par flux hydrique} (arcs volcaniques).
\end{itemize}
\textbf{Conclusion :} La nature du magma (basalte vs andésite/rhyolite) et son volume sont dictés par le \textbf{chemin $P-T$ suivi par le manteau} (décompression vs apport d'eau), lui-même imposé par la \textbf{géodynamique des plaques}.
\end{experience}

\begin{popfigure}
\begin{tikzpicture}[scale=0.9, >=Stealth]
% Axes
\draw[->, thick] (0,0) -- (0,6) node[left] {$P$ (GPa)};
\draw[->, thick] (0,0) -- (7,0) node[below] {$T$ ($^\circ\mathrm{C}$)};
% Échelles
\foreach \y/\yl in {1/1, 2/2, 3/3, 4/4, 5/5} \draw (0,\y) -- (-0.1,\y) node[left] {\yl};
\foreach \x/\xl in {1/1000, 2/1100, 3/1200, 4/1300, 5/1400, 6/1500} \draw (\x,0) -- (\x,-0.1) node[below] {\xl};
% Solidus sec
\draw[popred, very thick, domain=0:3, samples=50] plot (\x, {1.1 + 0.13*\x + 0.02*\x*\x}); % approx courbe
\node[popred, right] at (3, 1.8) {\textbf{Solidus sec}};
% Liquidus sec
\draw[popred, thick, dashed, domain=0:3, samples=50] plot (\x, {1.3 + 0.13*\x + 0.02*\x*\x});
\node[popred, right] at (3, 2.1) {Liquidus sec};
% Solidus humide
\draw[popblue, very thick, domain=0.5:3, samples=50] plot (\x, {0.8 + 0.1*\x});
\node[popblue, right] at (3, 1.1) {\textbf{Solidus humide ($\ce{H2O}$ sat.)}};
% Géotherme manteau normal (adabatique)
\draw[popgreen, thick, domain=0:3, samples=50] plot (\x, {1.3 + 0.05*\x});
\node[popgreen, above] at (3, 1.45) {Adiabatique manteau};
% Géotherme subduction (froid)
\draw[poppurple, thick, domain=0:3, samples=50] plot (\x, {0.6 + 0.15*\x});
\node[poppurple, below] at (3, 1.05) {Géotherme subduction};
% Zones de fusion
\fill[popred!20] (0,1.1) -- plot[domain=0:3] (\x, {1.1 + 0.13*\x + 0.02*\x*\x}) -- plot[domain=3:0] (\x, {1.3 + 0.05*\x}) -- cycle;
\node[popred!80!popdark] at (1.5, 1.4) {\textbf{Fusion dorsale}};
\fill[popblue!20] (0.5,0.85) -- plot[domain=0.5:3] (\x, {0.8 + 0.1*\x}) -- plot[domain=3:0.5] (\x, {0.6 + 0.15*\x}) -- cycle;
\node[popblue!80!popdark] at (2, 0.9) {\textbf{Fusion subduction}};
% Flèches processus
\draw[->, thick, popdark] (4, 1.5) -- (5.5, 1.5) node[right, align=left] {Remontée adiabatique\\(Dorsale, Point chaud)};
\draw[->, thick, popdark] (4, 0.9) -- (5.5, 0.9) node[right, align=left] {Injection $\ce{H2O}$\\(Subduction)};
\end{tikzpicture}
\legende{Diagramme $P-T$ schématique montrant les solidus de la péridotite (sec et humide) et les géothermes. \textbf{Zone rose :} Fusion par décompression adiabatique (dorsales, points chauds) $\rightarrow$ Magmas basaltiques (MORB, OIB). \textbf{Zone bleue :} Fusion par flux hydrique (subduction) $\rightarrow$ Magmas andésitiques à rhyolitiques, riches en volatils. Données issues des expériences à haute pression (Green \& Ringwood, 1967 ; Jaques \& Green, 1980).}
\end{popfigure}

\courssub{Les séismes : libération brutale de l'énergie élastique}

\begin{experience}[Modélisation de la rupture fragile et du cycle sismique (d'après H.F. Reid, 1910)]
\textbf{Observation :} En déformant lentement une règle en plastique ou une plaquette de chocolat entre les mains, on observe qu'elle résiste d'abord élastiquement (elle reprend sa forme si on relâche), puis casse net avec un bruit sec (rupture fragile) quand la contrainte dépasse la résistance.
\textbf{Interprétation :} Les roches de la lithosphère se comportent de même aux limites de plaques. Le mouvement relatif des plaques impose une déformation lente (accumulation d'énergie élastique) le long des failles verrouillées. Lorsque la contrainte de cisaillement dépasse la résistance au frottement de la faille, une rupture brutale se propage : c'est le séisme. L'énergie élastique stockée est convertie en ondes sismiques, chaleur et travail de fracture.
\textbf{Conclusion :} Le séisme est la manifestation superficielle d'une rupture profonde le long d'une faille préexistante ou nouvelle. La \textbf{théorie du rebond élastique} (Elastic Rebound Theory) explique la cyclicité : accumulation lente $\rightarrow$ rupture brutale $\rightarrow$ relâchement $\rightarrow$ nouvelle accumulation.
\end{experience}

\begin{popfigure}
\begin{tikzpicture}[scale=1.1, >=Stealth]
% Couches géologiques
\fill[popcream] (-4,-2) rectangle (4,0);
\fill[popblueL] (-4,-4) rectangle (4,-2);
\fill[poporangeL] (-4,-6) rectangle (4,-4);
% Faille
\draw[thick, popdark, decorate, decoration={zigzag, amplitude=2pt, segment length=4pt}] (-1,-6) -- (1,0);
% Foyer (Hypocentre)
\fill[popred] (0,-3) circle (4pt);
\node[popred, right=2pt] at (0,-3) {\textbf{Foyer (Hypocentre)}};
% Épicentre
\fill[popblue] (0,0) circle (4pt);
\node[popblue, above=2pt] at (0,0) {\textbf{Épicentre}};
\draw[dashed, popmuted] (0,-3) -- (0,0);
% Ondes sismiques (cercles concentriques)
\foreach \r/\c in {1.5/popred, 2.5/poporange, 3.5/popyellow} {
    \draw[\c, thick] (0,-3) circle (\r);
    \draw[\c, thick, dashed] (0,-3) circle (\r+0.2);
}
% Flèches de mouvement des blocs
\draw[<->, thick, popgreen] (-3.5,-1) -- (-2.5,-1) node[midway, above] {Mvt plaque};
\draw[<->, thick, popgreen] (2.5,-5) -- (3.5,-5) node[midway, above] {Mvt plaque};
% Légende ondes
\node[align=left, popdark] at (4.5, -1) {Ondes P (primaires)\\\& Ondes S (secondaires)};
\node[align=left, popdark] at (4.5, -4) {Ondes de surface\\ (Rayleigh, Love)};
% Profondeur
\draw[<->, thin] (-4.5,0) -- (-4.5,-3) node[midway, left] {Profondeur focale};
\end{tikzpicture}
\legende{Schéma d'un séisme : le foyer (hypocentre) est le point de rupture initial dans la lithosphère ; l'épicentre est sa projection verticale à la surface. Les ondes sismiques (P, S, de surface) se propagent sphériquement depuis le foyer. La profondeur focale classe les séismes : superficiels ($< 70$ km), intermédiaires ($70-300$ km), profonds ($> 300$ km, uniquement en zones de subduction).}
\end{popfigure}

\begin{definition}[Foyer, épicentre, magnitude, intensité]
Le \textbf{foyer} (ou \textbf{hypocentre}) est le point précis à l'intérieur de la Terre où la rupture démarre (coordonnées géographiques + profondeur). L'\textbf{épicentre} est le point à la surface terrestre (sol ou fond marin) situé à la verticale du foyer.
La \textbf{magnitude} ($M$) quantifie l'\textbf{énergie totale libérée} au foyer. L'échelle de Richter (logarithmique, base 10) est historique ; on utilise aujourd'hui la magnitude de moment $M_w$ (basée sur le moment sismique $M_0 = \mu \times S \times D$, avec $\mu$ module de rigidité, $S$ surface de rupture, $D$ glissement moyen). \textbf{Un séisme n'a qu'une seule magnitude.}
L'\textbf{intensité} (échelle MSK ou EMS-98, degrés I à XII) qualifie les \textbf{effets ressentis} (ressenti humain, dégâts aux bâtiments, modifications du sol) \textbf{en un lieu donné}. Elle décroît avec la distance à l'épicentre et dépend de la géologie locale (amplification sur sédiments meubles). \textbf{Un séisme a autant d'intensités que de localités touchées.}
\end{definition}

\begin{attention}[Erreur fréquente au Bac : Magnitude vs Intensité]
Ne confondez jamais : \textbf{Magnitude = Énergie au foyer (valeur unique, échelle logarithmique : $M+1 \approx 32\times$ plus d'énergie)}. \textbf{Intensité = Dégâts/ressenti en surface (valeur variable selon les lieux, échelle linéaire I-XII)}. Exemple : un séisme $M_w=7,0$ à $10$ km de profondeur sous une ville dense $\rightarrow$ Intensité IX-X (désastre). Le même séisme à $600$ km de profondeur $\rightarrow$ Intensité II-III (peu ressenti).
\end{attention}

\begin{exemplebox}[Le séisme de Sumatra-Andaman, 26 décembre 2004]
\textbf{Contexte géodynamique :} Subduction de la plaque indienne sous la plaque birmane (plaque eurasienne) le long de la fosse de Sunda.
\textbf{Données :} Magnitude de moment $M_w = 9,1$ (3ème plus grand séisme instrumental). Rupture sur $\sim 1300$ km de faille, glissement moyen $D \approx 10-15$ m, durée $\approx 500$ s. Profondeur focale $\sim 30$ km.
\textbf{Conséquence majeure :} Déplacement vertical brutal du plancher océanique ($\sim 5$ m de soulèvement) $\rightarrow$ génération d'un \textbf{tsunami} trans-océanique.
\textbf{Bilan :} $> 220\,000$ morts dans 14 pays riverains de l'océan Indien (Indonésie, Sri Lanka, Inde, Thaïlande, Somalie...). Absence de système d'alerte tsunami dans l'océan Indien à l'époque = vulnérabilité extrême.
\end{exemplebox}

\courssub{Répartition mondiale : la ceinture de feu du Pacifique}

\begin{popfigure}
\begin{tikzpicture}[scale=0.45]
% Fond océan
\fill[popblueL] (-180,-90) rectangle (180,90);
% Continents simplifiés (silhouettes)
\foreach \c in {popgreen!50!popdark, popgreen!60!popdark} {
    \fill[\c] (-180,90) -- (-180,30) -- (-150,30) -- (-120,50) -- (-90,60) -- (-60,70) -- (-30,90) -- cycle; % Am Nord
    \fill[\c] (-90,-10) -- (-90,-60) -- (-30,-60) -- (-30,-10) -- cycle; % Am Sud
    \fill[\c] (10,30) -- (10,70) -- (60,70) -- (60,30) -- cycle; % Europe/Asie O
    \fill[\c] (60,-10) -- (60,-40) -- (150,-40) -- (150,10) -- (120,10) -- (100,30) -- (60,30) -- cycle; % Asie E / Aus
    \fill[\c] (-20,-35) -- (-20,35) -- (50,35) -- (50,-35) -- cycle; % Afrique
}
% Ceinture de feu : points rouges (volcans) et cercles (séismes)
\foreach \x/\y in {
    -150/60, -140/55, -130/50, -120/45, -110/40, -100/35, -90/30, -80/20, -70/10, -60/0, -70/-10, -80/-20, -90/-30, -80/-40, -70/-50, % Amériques
    120/10, 130/20, 140/30, 150/40, 150/50, 140/55, 130/60, % Kamtchatka/Aléoutiennes/Japon
    120/-10, 125/-20, 130/-30, 120/-40, 110/-35, 100/-30, 95/-10, 90/0, 95/10, % Indonésie/Philippines
    170/-20, 175/-30, 180/-40, -180/-40, -170/-30, -160/-20 % Pacifique SW (Nouvelle-Zélande, Tonga)
} {
    \fill[popred!80!popdark] (\x,\y) circle (1.5);
    \draw[popred, thick] (\x,\y) circle (3);
}
% Légende
\node[anchor=south west, fill=white, fill opacity=0.8, text opacity=1, draw=popdark, rounded corners, align=center] at (-175,-85){
    \textbf{Ceinture de Feu du Pacifique} : $\approx 90\%$ des séismes mondiaux, $\approx 75\%$ des volcans actifs.\\
    \textcolor{popred}{\Large $\bullet$} Volcans actifs \quad \textcolor{popred}{\Large $\circ$} Séismes (épi/hypocentres)
};
% Dorsales médio-océaniques (séismes superficiels)
\draw[poporange, thick, decorate, decoration={snake, amplitude=2pt, segment length=10pt}] (-180,-10) -- (-100,-10) -- (-50,0) -- (0,-10) -- (50,-20) -- (180,-30);
\draw[poporange, thick, decorate, decoration={snake, amplitude=2pt, segment length=10pt}] (-180,10) -- (-100,10) -- (-50,20) -- (0,10) -- (50,20) -- (180,30);
\node[poporange, rotate=15] at (-150,15) {Dorsales : séismes superficiels, magnitude modérée};
\end{tikzpicture}
\legende{Carte mondiale simplifiée de la répartition des séismes et volcans. La \textbf{Ceinture de Feu du Pacifique} (subduction) concentre les séismes les plus puissants (profonds possibles) et le volcanisme explosif. Les \textbf{dorsales médio-océaniques} (accrétion) et les \textbf{failles transformantes} (coulissage, ex. San Andreas) génèrent des séismes superficiels ($< 30$ km) de magnitude généralement $< 8$. Les zones de collision continentale (Himalaya, Alpes) produisent des séismes superficiels destructeurs.}
\end{popfigure}

\courssub{Les tsunamis : quand l'océan s'emballe}

\begin{experience}[Mécanisme de génération d'un tsunami par séisme de subduction]
\textbf{Dispositif (maquette) :} Un bac rectangulaire rempli d'eau. Une plaque articulée au fond simule le plancher océanique (plaque chevauchante). Un piston la soulève brutalement de quelques centimètres.
\textbf{Observation :} Le soulèvement brutal déplace la colonne d'eau au-dessus. Une vague longue (longueur d'onde $\lambda \sim 100-200$ km) se propage radialement. En pleine mer, l'amplitude (hauteur) est faible ($\sim 0,5-1$ m), la période longue ($10-60$ min), la vitesse élevée ($v = \sqrt{g h} \approx 200$ m/s $\approx 720$ km/h pour $h=4000$ m). À l'approche de la côte ($h$ diminue), $v$ diminue, $\lambda$ diminue, l'énergie se concentre $\rightarrow$ \textbf{amplification} (hauteur $10-30$ m, voire plus), déferlement destructeur.
\textbf{Interprétation :} Le tsunami n'est \textbf{pas} une vague de vent. C'est une \textbf{onde de gravité à grande longueur d'onde} impliquant tout le volume d'eau. Le recul de la mer (retrait de l'eau) précède souvent l'arrivée de la crête (premier signe d'alerte naturel).
\textbf{Autres causes :} Glissement de terrain sous-marin (ex. 1998 Papouasie-Nouvelle-Guinée), éruption volcanique explosive sous-marine (ex. Hunga Tonga 2022), chute d'astéroïde.
\end{experience}

\begin{popfigure}
\begin{tikzpicture}[scale=0.8, >=Stealth]
% Profil subduction
\fill[popblueL] (-6,0) -- (6,0) -- (6,-1) -- (-6,-1) -- cycle; % Eau
\fill[poporangeL] (-6,-1) -- (6,-1) -- (4,-4) -- (-6,-4) -- cycle; % Plaque chevauchante
\fill[popgreenL] (-6,-4) -- (4,-4) -- (6,-6) -- (-6,-6) -- cycle; % Plaque subductrice
% Faille
\draw[thick, popred, decorate, decoration={zigzag, amplitude=3pt}] (-6,-1) -- (4,-4);
% Flèches mouvement
\draw[->, thick, popdark] (-5,-5) -- (-3,-5) node[midway, below] {Subduction};
\draw[->, thick, popdark] (5,-2) -- (3,-2) node[midway, above] {Verrouillage};
% Coseismique : soulèvement
\draw[<->, thick, popblue, double distance=2pt] (0,-1) -- (0,1) node[midway, right, align=center]{Soulèvement \\ coseismique $\Delta h$};
% Vague tsunami
\draw[popblue, very thick, domain=-6:6, samples=100] plot (\x, {0.5*sin(2*\x r) + 0.2*sin(4*\x r) + 1});
\draw[popblue, thick, domain=-6:6, samples=100] plot (\x, {0.3*sin(2*\x r - 30) + 0.1*sin(4*\x r - 30) + 1});
\draw[popblue, thick, domain=-6:6, samples=100] plot (\x, {0.2*sin(2*\x r - 60) + 0.05*sin(4*\x r - 60) + 1});
% Propagation
\draw[->, thick, popblue] (-6,1.5) -- (-7,1.5) node[left] {Propagation};
\draw[->, thick, popblue] (6,1.5) -- (7,1.5) node[right] {Propagation};
% Amplification côtière (zoom)
\begin{scope}[xshift=10cm, yshift=-1cm, scale=0.6]
\fill[popblueL] (-2,0) -- (2,0) -- (2,-3) -- (0,-3) -- (-2,0); % Côte pentue
\draw[popblue, very thick, domain=-2:2, samples=50] plot (\x, {1.5*exp(-(\x)^2/2)});
\draw[->, thick, popdark] (0,1.5) -- (0,2.5) node[above, align=center]{Amplification \\ (Shoaling)};
\draw[<->] (2,0) -- (2,-3) node[midway, right] {Profondeur $\downarrow$};
\end{scope}
% Légendes
\node[align=left, popdark] at (-7,-3) {1. Séisme de subduction\\ (Rupture faille)};
\node[align=left, popdark] at (-7,-4.5) {2. Déplacement vertical\\ brutal du fond marin};
\node[align=left, popdark] at (-7,-6) {3. Génération vague\\ longue période};
\node[align=left, popdark] at (13,-3) {4. Propagation rapide\\ en haute mer};
\node[align=left, popdark] at (13,-5) {5. Amplification\\ mortelle à la côte};
\end{tikzpicture}
\legende{Mécanisme de formation d'un tsunami par séisme de subduction. 1) Accumulation contraintes. 2) Rupture coseismique : soulèvement du plancher océanique (quelques mètres sur centaines de km). 3) Déplacement de la colonne d'eau $\rightarrow$ vague initiale. 4) Propagation transocéanique ($v \approx \sqrt{gh}$). 5) Effet de \textbf{shoaling} (refoulement) : ralentissement, raccourcissement longueur d'onde, croissance amplitude $\rightarrow$ submersion côtière.}
\end{popfigure}

\courssub{Le volcanisme : diversité des éruptions et des dangers}

La nature du magma (composition, température, teneur en gaz, viscosité) dicte le style éruptif, lui-même contrôlé par le contexte géodynamique (voir fusion péridotite ci-dessus).

\begin{propriete}[Typologie des éruptions et contextes géodynamiques]
\begin{center}
\begin{tabularx}{\linewidth}{|l|X|X|X|}\hline
\rowcolor{popblueL}
\textbf{Paramètre} & \textbf{Effusif (Hawaïen, Islandais)} & \textbf{Explosif (Vulcanien, Péléen, Plinien)} & \textbf{Phréatomagmatique} \\ \hline
\textbf{Contexte} & Dorsales (MORB), Points chauds (OIB), Rifts continentaux & Zones de \textbf{subduction} (Arcs insulaires, Cordillères) & Interaction magma/eau (lac, mer, nappe) \\ \hline
\textbf{Magma} & \textbf{Basaltique} : riche Mg,Fe ; pauvre SiO$_2$ ($< 52\%$) ; chaud ($1100-1200^\circ$C) ; \textbf{fluide} (faible viscosité) & \textbf{Andésitique / Rhyolitique} : riche SiO$_2$ ($> 55-70\%$) ; plus froid ($800-950^\circ$C) ; \textbf{très visqueux} (polymérisation Si-O) & Variable (souvent basaltique) \\ \hline
\textbf{Gaz} & Faible teneur ; dégazage passif facile (bulles montent) & Forte teneur ($H_2O, CO_2, SO_2$) ; bulles piégées $\rightarrow$ surpression & Vapeur d'eau massive (explosion thermique) \\ \hline
\textbf{Produits} & Coulées de lave fluides (pahoehoe, aa), fontaines de lave, scories & \textbf{Nuées ardentes} (pyroclastiques), cendres, ponces, blocs, bombes, lahars & Cendres fines, surges (nuées de base), maar \\ \hline
\textbf{Dangers directs} & Coulées lentes (destruction biens), gaz (SO$_2$) & \textbf{Nuées ardentes} (vitesse $> 100$ km/h, $T > 300^\circ$C), retombées cendres (toits, moteurs, santé), lahars (coulées de boue froides/chaudes) & Explosions imprévisibles, surges latérales \\ \hline
\textbf{Exemples} & Kilauea (Hawaï), Erta Ale (Éthiopie), Piton de la Fournaise, \textbf{Mont Cameroun}, \textbf{Nyiragongo} & Vésuve (79 ap. J.-C.), Pinatubo (1991), Montagne Pelée (1902) & Surtsey (Islande), Lac Nyos (dégazage limnique, voir encadré) \\ \hline
\end{tabularx}
\end{center}
\end{propriete}

\begin{savaistu}[La tragédie du Lac Nyos (Cameroun, 21 août 1986) : un aléa volcanique "silencieux"]
Le lac Nyos (région du Nord-Ouest, chaîne volcanique du Cameroun) est un lac de cratère (maar) profond ($\sim 200$ m), stratifié thermiquement. Le magmatisme profond (point chaud) injecte du \ce{CO2} qui se dissout dans les eaux profondes froides sous haute pression (stabilité type "bouteille de champagne").
\textbf{Déclencheur probable :} Glissement de terrain interne ou injection magmatique soudaine $\rightarrow$ renversement de la stratification (overturn) $\rightarrow$ dégazage brutal (ébullition explosive).
\textbf{Phénomène :} Libération de $\sim 1,6 \times 10^6$ tonnes de \ce{CO2} en quelques minutes. Nuage dense (plus lourd que l'air) s'écoulant dans les vallées à $\sim 50-100$ km/h sur $\sim 25$ km.
\textbf{Bilan :} $\approx 1\,746$ morts (asphyxie par anoxie/hypercapnie), 3\,500 têtes de bétail, végétation tuée.
\textbf{Prévention (Réduction du risque) :} Installation depuis 2001 de \textbf{tuyaux de dégazage} (auto-amorçage par bulles) pour prélever l'eau profonde riche en \ce{CO2} et la dégazer en surface de façon contrôlée. Surveillance continue (température, chimie, sismicité). C'est un exemple mondial de gestion de risque volcanique non-éruptif réussie.
\end{savaistu}

\courssub{Glissements de terrain et inondations : l'aggravation tectonique et anthropique}

Si les séismes, tsunamis et volcans sont des \textbf{aléas primaires} directs des mouvements de plaques, les glissements de terrain et les inondations sont souvent des \textbf{aléas secondaires} (déclenchés par un séisme, une éruption, une pluie intense sur relief tectonique) ou des aléas primaires climatiques dont la dangerosité est \textbf{exacerbée par le relief tectonique et l'aménagement humain}.

\begin{definition}[Glissement de terrain (Mouvement de versants)]
Déplacement vers le bas d'une masse de roches, débris ou terres sous l'effet de la gravité, le long d'une surface de rupture (glissement rotatif, translationnel) ou par écoulement (coulée de boue, lahar).
\textbf{Facteurs prédisposants (structurels) :} Relief fort (créé par tectonique/orogenèse), roches fracturées/friables (schistes, altérites), pente forte, dip des couches vers le vide.
\textbf{Facteurs déclenchants :} \textbf{Séismes} (secousses $\rightarrow$ perte cohésion, pression interstitielle), \textbf{Pluies intenses/prolongées} (saturation eau $\rightarrow$ hausse pression interstitielle, perte frottement), \textbf{Éruptions} (fonte neige/glace $\rightarrow$ lahars, surcharge cendres).
\textbf{Facteurs aggravants anthropiques :} \textbf{Déforestation} (perte ancrage racinaire, érosion), \textbf{Constructions anarchiques} (entailles au pied de pente, surcharge en tête, déviation écoulements), \textbf{Urbanisation non régulée} sur zones à risque connu.
\end{definition}

\begin{exemplebox}[Glissements de terrain au Cameroun : Bafoussam (2019) et Gouache (2023)]
\textbf{Bafoussam, 28 octobre 2019 :} Après plusieurs jours de pluies diluviennes (facteur déclenchant climatique), un glissement de terrain rotatif dans les altérites de base (saprolite) sur pente forte a emporté des habitations dans le quartier Nkeng. Bilan : $\approx 42$ morts. \textbf{Aggravants :} Urbanisation dense sur versants instables sans étude géotechnique, déforestation, absence de drainage.
\textbf{Gouache (Arrondissement de Galim, Région Ouest), 6 août 2023 :} Coulée de boue/débris (débris flow) déclenchée par orage violent sur sols volcaniques altérés (pentes du Massif du Bamboutos). Bilan : $\approx 20$ morts, maisons détruites. \textbf{Contexte :} Relief volcanique jeune, pentes raides, sols peu cohérents, agriculture sur brûlis/déforestation.
\textbf{Leçon :} Au Cameroun, le risque "mouvements de versants" est élevé dans l'Ouest (relief volcanique, altérites épaisses, fortes pluies, forte densité rurale/urbaine). La prévention passe par la cartographie d'aléa (Zonage), le reboisement, le drainage et l'interdiction de construire en pied de pente active.
\end{exemplebox}

\begin{definition}[Inondation]
Submersion temporaire d'une zone habituellement sèche par des eaux douces (crue rivière, ruissellement urbain) ou marines (tempête, tsunami).
\textbf{Facteurs naturels :} Pluies intenses (mousson, orages), relief (bassins versants pentus $\rightarrow$ concentration rapide), nature du sol (imperméable/argileux ou saturé).
\textbf{Facteurs anthropiques aggravants (majeurs en zone urbaine) :} \textbf{Imperméabilisation} des sols (béton, bitume $\rightarrow$ ruissellement coefficient $\uparrow$, temps de concentration $\downarrow$), \textbf{Occupation des lits majeurs} (zones d'expansion de crue naturelles) et zones humides, \textbf{Colmatage/rétrécissement} des lits mineurs (déchets, sable, constructions), \textbf{Déforestation} amont (érosion, colmatage aval).
\end{definition}

\begin{exemplebox}[Inondations récurrentes à Yaoundé et Douala (Cameroun)]
\textbf{Yaoundé (ville aux 7 collines) :} Réseau hydrographique dense (Mfoundi, Mingoa, Ekozoa). Urbanisation galopante ($\times 10$ pop. en 40 ans) $\rightarrow$ occupation des bas-fonds (marais asséchés), imperméabilisation massive, déchets bouchant caniveaux. Résultat : crues éclaires meurtrières quasi annuelles (ex. 2019, 2022), érosion régressive des berges.
\textbf{Douala (ville plate, estuaire Wouri) :} Altitude moyenne $< 10$ m. Risques combinés : crues fluviales (Wouri, Dibamba), montées marines (marées équinoxiales, élévation niveau mer), ruissellement urbain (sols latéritiques imperméables + béton). Quartiers informels (bidonvilles) en zone inondable = vulnérabilité extrême.
\textbf{Lagune de Maga (Extrême-Nord) :} Inondations saisonnières du Logone/Chari. Aménagement hydro-agricole (barrage de Maga, digues) a modifié l'hydraulique naturelle, créant des zones d'inondation permanentes ou soudaines pour les populations riveraines et les réfugiés.
\end{exemplebox}

\begin{aretenir}[Lien mouvements des plaques $\leftrightarrow$ Catastrophes]
\begin{itemize}
    \item \textbf{Lien DIRECT (Limites de plaques) :} \textbf{Séismes} (toutes limites), \textbf{Tsunamis} (subduction, failles transformantes océaniques, volcanisme côtier), \textbf{Volcanisme} (subduction = explosif ; accrétion/points chauds = effusif).
    \item \textbf{Lien INDIRECT / AGGRAVANT (Relief tectonique) :} La tectonique crée le \textbf{relief jeune et instable} (chaînes alpines, arcs volcaniques, rifts, failles actives). Ce relief favorise : pentes fortes $\rightarrow$ \textbf{glissements de terrain} ; orages orographiques $\rightarrow$ pluies intenses $\rightarrow$ \textbf{inondations} / glissements ; sismicité $\rightarrow$ déclencheur majeur glissements.
    \item \textbf{Facteur humain déterminant :} La \textbf{vulnérabilité} transforme l'aléa en catastrophe. Au Cameroun : croissance urbaine non planifiée, pauvreté, déforestation, manque de culture du risque.
\end{itemize}
\end{aretenir}

\courssub{Gestion des catastrophes naturelles : prévision, prédiction, prévention}

\begin{definition}[Prévision, Prédiction, Prévention : définitions opérationnelles]
\begin{itemize}
    \item \textbf{Prévision (Forecasting) :} Estimation \textbf{probabiliste} de l'occurrence d'un aléa sur le long/moyen terme. \textit{Ex :} "Probabilité de 60\% d'un séisme $M>6$ sur la faille de San Andreas d'ici 30 ans" ; "Zone à risque glissement (cartographie d'aléa)". Ne donne ni date, ni heure précise.
    \item \textbf{Prédiction (Prediction) :} Annonce \textbf{déterministe} (date, lieu, magnitude) d'un événement \textbf{imminent} (heures/jours avant). \textit{Ex :} Évacuation avant éruption du Pinatubo 1991 (signaux précurseurs : sismicité, déformation, gaz). \textbf{Actuellement impossible pour les séismes} (pas de précurseur fiable universel). Possible pour certains volcans et crues lentes.
    \item \textbf{Prévention (Mitigation) :} Ensemble des mesures \textbf{structurelles} (ouvrages : digues, parapluies sismiques, tuyaux Nyos, reboisement) et \textbf{non structurelles} (réglementation : PLU, normes parasismiques RPA-CM/Eurocode 8 ; éducation : exercices, culture du risque ; assurance ; systèmes d'alerte précoce) pour \textbf{réduire la vulnérabilité} et les conséquences.
\end{itemize}
\end{definition}

\begin{methode}[Analyser un document (texte, carte, photo, données) sur une catastrophe]
\begin{enumerate}
    \item \textbf{Identifier l'aléa :} Nature (séisme, tsunami, éruption, glissement, inondation), magnitude/intensité, date, localisation précise (coordonnées, région).
    \item \textbf{Contextualiser géodynamiquement :} Quelle limite de plaque ? (Subduction, collision, rift, faille transformante, point chaud ?). Mécanisme générateur (rupture fragile, fusion flux, soulèvement coseismique, instabilité pente).
    \item \textbf{Lister les facteurs aggravants (Vulnérabilité) :}
    \begin{itemize}
        \item \textit{Naturels :} topographie, géologie (amplification sismique, sols instables), climat (saison des pluies).
        \item \textit{Anthropiques :} densité de population, type d'habitat (normes parasismiques ?), urbanisation (zones à risque, imperméabilisation), déforestation, absence de plan de prévention/alerte, pauvreté.
    \end{itemize}
    \item \textbf{Quantifier les pertes (Bilan) :} Morts, blessés, déplacés, dégâts matériels (logements, infrastructures, agriculture), coûts économiques, impact environnemental.
    \item \textbf{Proposer/Évaluer des mesures de gestion (Prévention/Prévision/Protection) :} Existant ? Efficace ? Manquant ? (Voir définitions ci-dessus).
\end{enumerate}
\end{methode}

\begin{exoresolu}[Analyse de la catastrophe du 26 décembre 2004 (Sumatra)]
\textbf{Énoncé :} À partir des données : $M_w = 9,1$, épicentre $3,3^\circ$ N, $95,9^\circ$ E, profondeur $30$ km, faille de subduction Sunda, rupture $1300$ km, tsunami trans-océanique, $220\,000+$ morts. Appliquez la méthode.
\textbf{Solution détaillée :}
\begin{enumerate}
    \item \textbf{Aléa :} Séisme de subduction méga-thrust ($M_w 9,1$, très faible fréquence, énergie $\approx 10^{23}$ J) + Tsunami trans-océanique (hauteur $> 30$ m localement).
    \item \textbf{Contexte géodynamique :} Convergence oblique plaque Indienne / Plaque Birmane (vitesse $\sim 6$ cm/an). Subduction le long de la fosse de Sunda. Rupture coseismique sur segment verrouillé $\rightarrow$ soulèvement brutal plancher océanique ($\sim 5-10$ m) $\rightarrow$ déplacement colonne d'eau.
    \item \textbf{Facteurs aggravants :} \textit{Naturels :} Proximité côtes (Sumatra, Aceh) $\rightarrow$ temps d'alerte $< 15-30$ min. Bathymétrie côtière amplifiant vague. \textit{Anthropiques :} Densité côtière élevée (pêche, tourisme). \textbf{Absence totale de système d'alerte tsunami dans l'océan Indien} (contrairement au Pacifique). Construction non adaptée (bois, bambou). Mangroves détruites (perte tampon naturel).
    \item \textbf{Bilan :} $\sim 227\,898$ morts (Indonésie $\sim 167\,000$), millions de sans-abri, destruction infrastructures (ports, routes, électricité), contamination eau/sol (sel, cadavres), impact psychologique durable.
    \item \textbf{Gestion / Leçons :} Création du \textbf{Système d'Alerte Tsunami de l'Océan Indien (IOTWS)} opérationnel 2006. Exercices d'évacuation réguliers. Replantation mangroves. Aménagement côtier (zones tampons). Au Cameroun, cela inspire la mise en place d'alertes pour le Mont Cameroun et le Lac Nyos.
\end{enumerate}
\end{exoresolu}

\begin{exercice}[Application : Le risque sismique à Kribi (Cameroun)]
Kribi, ville côtière du Sud-Cameroun, est située près de la zone de cisaillement de Kribi-Campo (faille transformante/transpressionnelle liée à l'ouverture de l'Atlantique Sud). Des séismes modérés ($M \approx 4-5$) y sont enregistrés (ex. 2005, 2012). La ville s'étend sur des collines d'altérites latéritiques et des mangroves côtières. Le port en eau profonde et l'usine GNL (Gaz Naturel Liquéfié) y sont implantés.
\begin{enumerate}
    \item Qualifiez l'aléa sismique (type de limite de plaque, magnitude max probable, profondeur).
    \item Identifiez 3 facteurs de vulnérabilité spécifiques à Kribi (urbain, industriel, environnemental).
    \item Proposez 2 mesures de prévention structurelle et 2 mesures non structurelles prioritaires.
\end{enumerate}
\end{exercice}

\begin{corrige}[Exercice n°1 : Risque sismique à Kribi]
\begin{enumerate}
    \item \textbf{Aléa :} Faille transformante/transpressionnelle (coulissage dextre + compression) héritée de l'ouverture de l'Atlantique. Séismes superficiels ($< 30$ km). Magnitude max probable estimée $M_w \approx 6,0-6,5$ (longueur segments failles $\sim 50-100$ km). Aléa \textbf{modéré} (probabilité faible à modérée, mais intensité locale forte possible du fait de la faible profondeur).
    \item \textbf{Vulnérabilités :} 1) \textit{Urbain :} Habitat informel dense sur pentes instables (altérites) $\rightarrow$ amplification sismique + risque glissements sismo-induits. 2) \textit{Industriel :} Usine GNL (risque technologique "Natech" : séisme $\rightarrow$ fuite GNL $\rightarrow$ explosion/incendie). 3) \textit{Environnemental :} Mangroves/sables meubles côtiers $\rightarrow$ amplification de site (facteur 2-5), liquéfaction possible.
    \item \textbf{Prévention :} \textit{Structurelle :} a) Normes parasismiques obligatoires (Eurocode 8 / RPA-CM) pour TOUTES nouvelles constructions (port, GNL, habitats). b) Stabilisation pentes (drainage, murs soutènement, reboisement) zones d'habitat existant. \textit{Non structurelle :} a) Microzonage sismique détaillé (carte amplification/liquéfaction) pour guider PLU (Plan Local d'Urbanisme). b) Exercices simulation séisme/tsunami (risque tsunami faible mais non nul si glissement sous-marin) pour population et personnel industriel ; Plan d'Opération Interne (POI) GNL testé annuellement.
\end{enumerate}
\end{corrige}

\coursec{Gestion des catastrophes naturelles : prévision, prédiction, prévention}

Dans la section précédente, tu as vu que les mouvements des plaques lithosphériques engendrent des catastrophes naturelles : séismes, éruptions volcaniques, tsunamis, glissements de terrain et inondations. Une question se pose alors immédiatement au citoyen comme au décideur : peut-on éviter ces catastrophes ? La réponse honnête de la science est nuancée : on ne peut pas empêcher un séisme ni une éruption, mais on peut \cle{réduire considérablement leurs conséquences}. C'est tout l'objet de la \cle{gestion des risques naturels}, qui repose sur trois piliers qu'il faut savoir distinguer avec précision : la \cle{prédiction}, la \cle{prévision} et la \cle{prévention}, complétés par la \cle{sensibilisation} des populations.

\courssub{Prédiction et prévision : deux notions à ne pas confondre}

\textbf{Observation préalable.}\par
En décembre 2004, un séisme de magnitude 9,1 au large de Sumatra déclenche un tsunami qui tue environ 230 000 personnes autour de l'océan Indien. Aucun signal d'alerte n'avait été diffusé. En revanche, au Japon, pays très fréquemment frappé, les séismes de magnitude comparable font souvent beaucoup moins de victimes. D'où vient cette différence ? La science sait-elle annoncer un séisme à l'avance ?

\begin{definition}[Prédiction]
La \cle{prédiction} consiste à annoncer, de manière certaine, le \textbf{lieu}, la \textbf{date} et l'\textbf{ampleur} (magnitude) d'un événement catastrophique avant qu'il ne se produise.
\end{definition}

\begin{aretenir}[Un savoir admis : la prédiction des séismes est impossible]
À ce jour, \textbf{aucune méthode fiable ne permet de prédire un séisme}. Malgré des décennies de recherches sur d'éventuels signes précurseurs (variations du radon dans les sols, comportement animal, signaux électromagnétiques, petites secousses avant-coureuses), aucun précurseur universel et reproductible n'a été identifié. En revanche, l'éruption volcanique est parfois \emph{presque} prévisible à court terme grâce à la surveillance, car le magma qui monte produit des signaux détectables (sismicité, gonflement du volcan, gaz).
\end{aretenir}

\begin{definition}[Prévision]
La \cle{prévision} consiste à estimer la \textbf{probabilité d'occurrence} d'un événement dans une \textbf{zone} et une \textbf{période} données, à partir de données \textbf{statistiques} (sismicité historique), \textbf{géologiques} (cartographie des failles actives) et de \textbf{surveillance instrumentale} (déformations du sol mesurées par GPS, enregistrements des sismomètres).
\end{definition}

L'intuition est simple : si une faille a produit des séismes majeurs tous les 100 à 150 ans au cours des derniers millénaires, et que le dernier date de 140 ans, on peut dire que la probabilité d'un séisme y est élevée dans les prochaines décennies — sans pouvoir dire ni le jour, ni l'heure, ni la magnitude exacte. C'est une démarche \emph{probabiliste}, analogue aux prévisions météorologiques à long terme.

\begin{experience}[Document : la prévision sismique en pratique]
\textbf{Données disponibles pour une zone de subduction :}
\begin{itemize}
\item \textbf{Sismicité historique} : catalogues des séismes passés (dates, magnitudes, épicentres) sur plusieurs siècles ;
\item \textbf{Géologie} : repérage des failles actives, mesure des déplacements anciens (paléosismologie : fouilles de tranchées à travers les failles) ;
\item \textbf{Surveillance} : réseaux de \textbf{sismomètres} enregistrant les micro-séismes, stations \textbf{GPS} mesurant l'accumulation des déformations (5 cm/an), extensomètres et inclinomètres.
\end{itemize}
\textbf{Interprétation :} quand le GPS montre qu'une zone de faille est « verrouillée » alors que les plaques avancent de 5 cm/an, la contrainte s'accumule : la faille rompra un jour. On établit alors des \textbf{cartes d'aléa sismique} (zones de forte, moyenne ou faible probabilité), qui servent de base aux normes de construction.
\textbf{Conclusion :} la prévision délimite les zones à risque et les périodes statistiquement favorables, mais ne donne jamais de date.
\end{experience}

\begin{attention}[Erreur fréquente au Bac]
Ne jamais écrire « les scientifiques peuvent prédire les séismes » ou « grâce aux sismomètres, on prévoit la date du prochain séisme ». C'est \textbf{faux}. Le sismomètre \emph{enregistre} les ondes quand le séisme a déjà commencé ; il ne prédit rien. Seule la \textbf{prévision probabiliste} est possible pour les séismes. Réserve le mot « prédiction » à son sens strict (lieu + date + ampleur annoncés à l'avance), qui reste hors de portée de la science actuelle.
\end{attention}

\courssub{La surveillance : des volcans aux tsunamis}

Si l'on ne peut pas prédire, on peut \cle{surveiller}. La surveillance rapprochée permet parfois de détecter les signes annonciateurs d'une éruption ou de déclencher une alerte quelques minutes après un séisme.

\textbf{La surveillance volcanique.} Un volcan actif est surveillé par plusieurs paramètres mesurés en continu :
\begin{itemize}
\item la \textbf{sismicité} : la remontée du magma fracture les roches et provoque des essaims de petits séismes ;
\item les \textbf{déformations} : le volcan gonfle sous la pression du magma (mesuré par GPS et inclinomètres) ;
\item les \textbf{émissions de gaz} : l'augmentation du \ce{SO2} ou du rapport \ce{CO2}/\ce{SO2} traduit l'arrivée de magma frais ;
\item la \textbf{température} des fumerolles et des sources chaudes.
\end{itemize}
Lorsque plusieurs de ces paramètres dérivent simultanément, les volcanologues peuvent recommander une \textbf{évacuation préventive}. Au Cameroun, le mont Cameroun (\SI{4095}{m}), volcan actif dont les éruptions récentes (1982, 1999, 2000) ont menacé les plantations et les villages de la plaine de Buea, fait l'objet d'une surveillance par l'Institut de recherches géologiques et minières (IRGM) ; les coulées de lave de 1999 et 2000 ont pu être suivies et les populations informées à temps.

\textbf{Les systèmes d'alerte aux tsunamis.} Un tsunami se propage dans l'océan à environ \SI{800}{km/h} en pleine mer, mais met des dizaines de minutes à quelques heures pour atteindre les côtes lointaines. Ce délai est exploitable : un réseau de \textbf{sismographes} détecte le séisme sous-marin et localise son foyer ; si la magnitude et la position sont compatibles avec un tsunami, des \textbf{bouées océaniques} (système DART : capteurs de pression posés sur le fond, reliés par acoustique à une bouée de surface puis par satellite aux centres d'alerte) confirment le passage de l'onde ; une alerte est alors diffusée aux populations côtières.

\begin{exemplebox}[Un exemple chiffré : l'océan Indien après 2004]
Avant 2004, l'océan Indien ne disposait d'aucun système d'alerte aux tsunamis : l'onde générée au large de Sumatra a frappé les côtes sans aucun avertissement. Après la catastrophe, un système d'alerte a été installé (sismographes, bouées DART, centres régionaux). Il permet désormais de diffuser une \textbf{alerte en moins de 15 minutes} après un séisme sous-marin majeur, laissant aux populations côtières proches de l'épicentre un temps précieux pour gagner les hauteurs, et plusieurs heures aux pays lointains. Le système du Pacifique, plus ancien (créé après le tsunami de 1946 à Hawaï), fonctionne sur le même principe.
\end{exemplebox}

\courssub{La prévention : réduire la vulnérabilité avant la catastrophe}

\begin{definition}[Prévention]
La \cle{prévention} est l'ensemble des \textbf{mesures prises avant la catastrophe} pour réduire la \textbf{vulnérabilité} des populations et des biens, c'est-à-dire leur exposition aux dégâts.
\end{definition}

Retiens la formule synthétique : \textbf{Risque = Aléa $\times$ Vulnérabilité}. L'aléa (le phénomène naturel) ne dépend pas de nous ; en revanche, la vulnérabilité peut être fortement réduite par l'action humaine. Les principales mesures de prévention sont :

\begin{itemize}
\item les \textbf{normes de construction parasismique} : bâtiments renforcés (chaînages, fondations adaptées, matériaux souples) capables de résister aux secousses ; c'est l'effondrement des bâtiments, non la secousse elle-même, qui tue ;
\item l'\textbf{interdiction de construire en zones exposées} : zones inondables, pied de falaises instables, pentes raides, littoral à risque de tsunami, axes de coulées de lave ;
\item le \textbf{reboisement des versants} : les racines fixent le sol et limitent les glissements de terrain et l'érosion (mesure cruciale dans les montagnes de l'Ouest camerounais, où la déforestation aggrave les glissements de terrain meurtriers, comme ceux de Bafoussam en 2019) ;
\item l'\textbf{aménagement des berges} et des cours d'eau (digues, curage, caniveaux) pour limiter les inondations, problème récurrent à Douala et Yaoundé en saison des pluies ;
\item l'élaboration de \textbf{plans d'évacuation} : itinéraires balisés, points de rassemblement, exercices réguliers.
\end{itemize}

\courssub{La sensibilisation : le maillon citoyen de la gestion du risque}

\begin{definition}[Sensibilisation]
La \cle{sensibilisation} est l'ensemble des actions d'\textbf{éducation et d'information} visant à faire connaître à la population les risques encourus et les \textbf{conduites à tenir} avant, pendant et après une catastrophe.
\end{definition}

Une alerte qui n'est pas comprise ne sert à rien : lors du tsunami de 2004, beaucoup de victimes se sont approchées du rivage, intriguées par le retrait anormal de la mer, signe pourtant classique de l'arrivée imminente de l'onde. La sensibilisation passe par les \textbf{écoles}, les \textbf{médias} (radio, télévision), les \textbf{affiches}, les \textbf{exercices d'évacuation} et les chefs de quartier. Au Cameroun, la gestion des catastrophes relève de la \textbf{Direction de la Protection civile} (MINAT), qui coordonne les secours et l'information des populations, en lien avec les autorités administratives et traditionnelles.

\begin{savaistu}[Le Japon, champion de la culture du risque]
Le Japon organise chaque année (notamment le 1\textsuperscript{er} septembre, « Journée de la prévention des désastres », anniversaire du séisme de Tokyo de 1923) des \textbf{exercices nationaux d'évacuation sismique} impliquant des \textbf{millions de citoyens} : écoliers, employés, familles s'entraînent aux gestes qui sauvent. Résultat : malgré une sismicité parmi les plus fortes du monde, la combinaison « normes parasismiques + population entraînée » réduit fortement le nombre de victimes. La leçon est claire : la prévention et la sensibilisation sauvent plus de vies que n'importe quelle technologie de prédiction.
\end{savaistu}

\begin{methode}[Élaborer un outil de sensibilisation (affiche, dépliant, message radio, sketch)]
\begin{enumerate}
\item \textbf{Cibler le public} : écoliers, habitants d'un quartier inondable, pêcheurs du littoral de Limbe\ldots{} Le langage et les images dépendent du public.
\item \textbf{Choisir un message simple et illustré} : une seule idée principale par support, un slogan court, des dessins plutôt que de longs textes.
\item \textbf{Indiquer les conduites à tenir} structurées en trois temps : \emph{avant} (préparer), \emph{pendant} (se protéger), \emph{après} (secourir et éviter les risques secondaires).
\item \textbf{Adapter le support au contexte local} : langue parlée, radio communautaire en zone rurale, sketch théâtral au marché, affiche à l'école et au centre de santé.
\item \textbf{Faire valider} les informations par une source fiable (Protection civile, enseignants, personnels de santé) avant diffusion.
\end{enumerate}
\end{methode}

\begin{aretenir}[Conduites à tenir : l'essentiel]
\textbf{Pendant un séisme :} ne pas courir vers les sorties ; s'abriter \textbf{sous une table solide} ou contre un mur porteur, à l'écart des fenêtres ; à l'extérieur, \textbf{s'éloigner des façades}, des fils électriques et des arbres. \textbf{Après :} couper l'eau, le gaz et l'électricité, ne pas utiliser les ascenseurs, se méfier des répliques.
\textbf{Pendant une inondation :} \textbf{rejoindre les hauteurs}, ne jamais traverser à pied ou en voiture une eau courante (30 cm d'eau en mouvement peuvent emporter une personne), couper l'électricité.
\textbf{Lors d'une alerte tsunami :} si la mer se retire anormalement ou si un séisme long est ressenti sur la côte, \textbf{s'éloigner immédiatement du littoral} et gagner les hauteurs, sans attendre l'alerte officielle.
\end{aretenir}

\begin{popfigure}
\begin{center}
\begin{tikzpicture}[font=\small, >=Stealth]
\tikzset{
  etape/.style={rectangle, rounded corners=3pt, draw=popblue, fill=popblueL, minimum width=2.4cm, minimum height=0.9cm, align=center},
  etape2/.style={rectangle, rounded corners=3pt, draw=poporange, fill=poporangeL, minimum width=2.4cm, minimum height=0.9cm, align=center}
}
\node[etape, fill=popgoldL, draw=popgold, align=center] (alea) at (0,0){\textbf{Aléa}\\ (séisme, éruption,\\ tsunami\ldots)};
\node[etape, align=center] (surv) at (3.4,0){\textbf{Surveillance}\\ sismomètres, GPS,\\ gaz, bouées};
\node[etape, align=center] (prev) at (6.8,0){\textbf{Prévision}\\ probabilité,\\ cartes d'aléa};
\node[etape2, align=center] (alerte) at (10.2,0){\textbf{Alerte}\\ diffusion rapide\\ aux populations};
\node[etape2, draw=popgreen, fill=popgreenL, align=center] (preven) at (6.8,-2.2){\textbf{Prévention}\\ normes, zonage,\\ reboisement};
\node[etape2, draw=poppurple, fill=poppurpleL, align=center] (sens) at (3.4,-2.2){\textbf{Sensibilisation}\\ écoles, médias,\\ exercices};
\draw[->, thick, popblue] (alea) -- (surv);
\draw[->, thick, popblue] (surv) -- (prev);
\draw[->, thick, poporange] (prev) -- (alerte);
\draw[->, thick, popgreen] (prev) -- (preven);
\draw[->, thick, poppurple] (preven) -- (sens);
\draw[->, thick, poppurple, dashed] (sens.west) .. controls (-1.2,-2.2) and (-1.2,-0.6) .. (alea.south west) node[midway, left, align=center, text=poppurple]{population\\ préparée =\\ vulnérabilité réduite};
\end{tikzpicture}
\end{center}
\legende{Schéma fonctionnel de la chaîne de gestion du risque naturel. L'aléa est surveillé (instruments) ; la prévision probabiliste alimente à la fois l'alerte à court terme et la prévention à long terme ; la sensibilisation boucle le dispositif en réduisant la vulnérabilité des populations.}
\end{popfigure}

\begin{popfigure}
\begin{center}
\begin{tikzpicture}[font=\small]
\draw[fill=popcream, draw=popdark, thick, rounded corners=4pt] (0,0) rectangle (11.5,7.2);
\node[anchor=north, font=\large\bfseries, text=popink] at (5.75,7.0) {AFFICHE — CONDUITES À TENIR EN CAS DE SÉISME};
\draw[popdark] (1,6.35) -- (10.5,6.35);
\node[draw=popblue, fill=popblueL, rounded corners=2pt, minimum width=3.1cm, minimum height=2.2cm, align=center] at (2.6,4.6) {\textbf{AVANT}\\ Préparer un sac\\ d'urgence ; fixer les\\ meubles hauts ;\\ repérer les issues};
\node[draw=poporange, fill=poporangeL, rounded corners=2pt, minimum width=3.1cm, minimum height=2.2cm, align=center] at (5.75,4.6) {\textbf{PENDANT}\\ S'abriter sous une\\ table ; s'éloigner\\ des façades et\\ des vitres};
\node[draw=popgreen, fill=popgreenL, rounded corners=2pt, minimum width=3.1cm, minimum height=2.2cm, align=center] at (8.9,4.6) {\textbf{APRÈS}\\ Couper eau, gaz,\\ électricité ; écouter\\ la radio ; attention\\ aux répliques};
% petit pictogramme : personne sous une table
\draw[thick, popdark] (4.9,2.6) -- (4.9,1.7) -- (6.6,1.7) -- (6.6,2.6);
\draw[thick, popdark] (4.9,2.6) -- (6.6,2.6);
\draw[fill=poppinkL, draw=poppink] (5.75,2.25) circle (0.14);
\draw[thick, poppink] (5.75,2.11) -- (5.75,1.82);
\draw[thick, poppink] (5.75,2.0) -- (5.45,1.9);
\draw[thick, poppink] (5.75,2.0) -- (6.05,1.9);
\draw[thick, poppink] (5.75,1.82) -- (5.5,1.75);
\draw[thick, poppink] (5.75,1.82) -- (6.0,1.75);
\node[align=center, text=popmuted] at (5.75,1.2) {Reste calme, protège ta tête,\\ ne cours pas vers les sorties.};
\node[anchor=south, font=\bfseries, text=popdark] at (5.75,0.15) {Direction de la Protection civile — MINAT — République du Cameroun};
\end{tikzpicture}
\end{center}
\legende{Gabarit d'affiche de sensibilisation : un titre clair, trois cases « avant / pendant / après », un pictogramme central (personne abritée sous une table) et la mention de l'autorité compétente. Ce type de support peut être reproduit en maquette ou en dépliant.}
\end{popfigure}

\begin{exoresolu}[Savoir distinguer prédiction, prévision et prévention]
\textbf{Énoncé.} Pour chacune des affirmations suivantes, indique s'il s'agit de prédiction, de prévision, de prévention ou de sensibilisation, en justifiant brièvement : (a) « Une faille active de Californie a 60 \% de chances de produire un séisme de magnitude supérieure à 6,7 dans les 30 prochaines années. » (b) « Les immeubles neufs de cette ville doivent respecter des normes parasismiques. » (c) « Un séisme de magnitude 7 frappera la ville le 15 mars à 14 h. » (d) « Les élèves du lycée participent à un exercice d'évacuation. »
\tcblower
\textbf{Solution.} (a) \textbf{Prévision} : on donne une probabilité (60 \%), une zone (Californie) et une période (30 ans), sans date précise — c'est exactement la définition de la prévision probabiliste. (b) \textbf{Prévention} : les normes parasismiques sont une mesure prise avant toute catastrophe pour réduire la vulnérabilité des bâtiments. (c) Il s'agit d'une \textbf{prédiction} (lieu, date et ampleur annoncés) : or une telle annonce est \textbf{scientifiquement impossible} à l'heure actuelle ; cette affirmation relève de la rumeur ou de la charlatanerie. (d) \textbf{Sensibilisation} : l'exercice d'évacuation éduque la population aux conduites à tenir.
\end{exoresolu}

\begin{exercice}[À toi de jouer]
Ta commune est située sur le flanc du mont Cameroun, dans une zone déjà traversée par des coulées de lave. Le maire te demande, à toi et à ta classe de Terminale D, de concevoir un dépliant de sensibilisation destiné aux habitants. Présente : (1) le public cible et le support choisi ; (2) les trois mesures de prévention que tu recommandes ; (3) les conduites à tenir avant, pendant et après une éruption.
\end{exercice}

\begin{corrige}[Exercice : éléments de correction]
(1) Public : habitants des villages exposés (Buea, Limbe et leurs environs), y compris les non-lecteurs ; support : dépliant illustré en français et en langues locales, complété par des messages à la radio communautaire. (2) Prévention : interdire les constructions sur les axes historiques des coulées ; aménager des itinéraires et points de rassemblement d'évacuation ; maintenir la surveillance volcanologique (sismicité, déformations, gaz) et informer régulièrement la population. (3) Avant : préparer un sac d'urgence (eau, lampe, documents, masque contre les cendres), connaître l'itinéraire d'évacuation. Pendant : suivre les consignes des autorités, se protéger des cendres (masque, lunettes), évacuer si ordonné sans retourner chercher des biens. Après : ne revenir qu'après autorisation, éviter les zones de cendres instables et les lahars (coulées de boue) en cas de pluie, écouter la radio pour les consignes sanitaires.
\end{corrige}

\begin{aretenir}[L'essentiel de la section]
\begin{itemize}
\item La \textbf{prédiction} (lieu + date + ampleur) des séismes est \textbf{impossible} : c'est un savoir admis.
\item La \textbf{prévision} est une estimation \textbf{probabiliste} fondée sur la sismicité historique, la géologie des failles et la surveillance (GPS, sismomètres).
\item La \textbf{surveillance} des volcans (sismicité, déformations, gaz, température) et les \textbf{systèmes d'alerte tsunami} (sismographes + bouées, alerte en moins de 15 minutes dans l'océan Indien) permettent parfois des évacuations à temps.
\item La \textbf{prévention} réduit la vulnérabilité : normes parasismiques, zonage, reboisement, aménagement des berges, plans d'évacuation.
\item La \textbf{sensibilisation} (écoles, médias, exercices, Protection civile) transforme le citoyen en acteur de sa propre sécurité : \textbf{Risque = Aléa $\times$ Vulnérabilité}.
\end{itemize}
\end{aretenir}

\newpage

\coursec{Activité d'intégration}

\coursec{Leçon 15 : Réduction des risques liés aux catastrophes naturelles}

\courssub{Objectifs de l'activité}
\par
Cette activité d'intégration vise à mobiliser l'ensemble des savoirs et savoir-faire du Module III pour résoudre une situation complexe et authentique : l'élaboration d'un plan de protection pour la ville de Limbé. À l'issue de cette activité, vous devez être capable de :
\begin{itemize}
    \item Analyser un profil bathymétrique pour identifier les structures de la dorsale médio-atlantique et en déduire la cinématique des plaques.
    \item Exploiter une courbe \emph{âge du plancher océanique / distance à l'axe de la dorsale} pour calculer un taux d'expansion (vitesse de divergence).
    \item Interpréter un diagramme pression-température (solidus anhydre vs hydraté de la péridotite) pour expliquer l'origine des magmas de la Ligne du Cameroun (volcanisme intraplaque).
    \item Synthétiser les risques naturels majeurs (volcaniques, sismiques, limniques, gravitaires, hydrologiques) affectant le Cameroun.
    \item Concevoir un outil de sensibilisation (affiche) clair, rigoureux et adapté à la population locale, intégrant prévision, prédiction et prévention.
\end{itemize}

\courssub{Situation}
\par
\textbf{Contexte authentique.} La ville de \textbf{Limbé} (ex-Victoria), chef-lieu du département du Fako (Région du Sud-Ouest), compte environ \num{150000} habitants (estimation 2023). Elle est située sur la côte atlantique, au pied du flanc sud-ouest du \textbf{mont Cameroun} (volcan actif, \SI{4095}{m} d'altitude, dernière éruption en 2000), et à proximité de la \textbf{fosse de l'Atlantique équatorial}. La ville concentre des enjeux économiques majeurs : raffinerie (SONARA), port en eau profonde, plantations agro-industrielles (CDC), et tourisme balnéaire.

\par
Le \textbf{Gouvernement de la République du Cameroun}, via le Ministère de l'Administration Territoriale (MINAT) et la Protection Civile, mandate une \textbf{Commission d'Experts Pluridisciplinaire} (vous) pour produire une \emph{Affiche de Sensibilisation aux Risques Majeurs} à destination des habitants de Limbé et des zones riveraines (Bota, Mile 4, Idenau, Batoke). Cette affiche doit être scientifiquement fondée, lisible par des non-spécialistes, et comporter des consignes de conduite à tenir avant, pendant et après une catastrophe.

\par
\textbf{Dossier documentaire fourni à chaque groupe d'experts :}

\par
\textbf{Document 1 : Profil bathymétrique simplifié de l'Atlantique équatorial (traversée Ouest-Est, latitude 4° N).}
\par
Le profil montre (de l'Ouest vers l'Est) :
\begin{itemize}
    \item \textbf{Zone A (0--300 km) :} Plaine abyssale plate à \SI{-4800}{m}, sédiments épais.
    \item \textbf{Zone B (300--450 km) :} Relief accidenté, monts sous-marins alignés (chaîne de fractures), failles transformantes actives (sismicité superficielle).
    \item \textbf{Zone C (450--550 km) :} \textbf{Dorsale médio-atlantique}. Vallée de rift centrale (\SI{-3500}{m} au fond, flancs à \SI{-2500}{m}), axe magmatique (coulées de lave basaltiques récentes, anomalie magnétique positive forte).
    \item \textbf{Zone D (550--700 km) :} Flanc est de la dorsale, topographie plus lisse, plaine abyssale à \SI{-4800}{m}.
    \item \textbf{Zone E (700--900 km) :} Montée du plateau continental, pente continentale raide (\SI{-200}{m} à \SI{-3000}{m}), marge passive camerounaise.
    \item \textbf{Zone F (>900 km) :} Plateau continental (\SI{-50}{m} à \SI{-200}{m}), côte de Limbé (0 m).
\end{itemize}

\begin{popfigure}
\begin{tikzpicture}[scale=0.7]
% Axes
\draw[->, thick] (0,0) -- (12,0) node[right] {Distance (km)};
\draw[->, thick] (0,-5.5) -- (0,1) node[above] {Profondeur (m)};
% Grid
\foreach \x in {0,2,...,12} \draw[popline, thin] (\x,-5.5) -- (\x,0.5);
\foreach \y in {-5,-4,...,0} \draw[popline, thin] (0,\y) -- (12,\y);
% Labels axes
\foreach \x/\label in {0/0, 2/200, 4/400, 6/600, 8/800, 10/1000} \node[below] at (\x,-5.7) {\small \label};
\foreach \y/\label in {0/0, -1/-1000, -2/-2000, -3/-3000, -4/-4000, -5/-5000} \node[left] at (-0.3,\y) {\small \label};

% Profile drawing
% Zone A: Abyssal plain West
\draw[thick, popblue] (0,-4.8) -- (3,-4.8);
% Zone B: Fracture zone / rough
\draw[thick, popblue] (3,-4.8) .. controls (3.5,-4.5) and (4,-5.2) .. (4.5,-4.8) .. controls (5,-4.4) and (5.5,-5.1) .. (5.5,-4.8);
% Zone C: Ridge (MAR)
\draw[thick, poporange] (5.5,-4.8) .. controls (5.8,-4.0) and (6.2,-4.0) .. (6.5,-3.5) node[midway, above=0.3cm, poporange] {\small Axe dorsal};
\draw[thick, poporange] (6.5,-3.5) .. controls (6.8,-4.0) and (7.2,-4.0) .. (7.5,-4.8);
% Zone D: Flank East
\draw[thick, popblue] (7.5,-4.8) -- (9,-4.8);
% Zone E: Continental slope
\draw[thick, popgreen] (9,-4.8) .. controls (9.5,-3.5) and (10,-2.0) .. (10.5,-0.5);
% Zone F: Shelf
\draw[thick, popgreen] (10.5,-0.5) -- (12,0);

% Zone labels
\node[align=center, font=\small] at (1.5, 0.5) {Zone A\\Plaine abyssale\\Ouest};
\node[align=center, font=\small] at (4.5, 0.5) {Zone B\\Zones de fracture\\Transformantes};
\node[align=center, font=\small, poporange] at (6.5, 1.2) {Zone C\\Dorsale Médio-Atlantique\\Vallée de Rift};
\node[align=center, font=\small] at (8.2, 0.5) {Zone D\\Flanc Est};
\node[align=center, font=\small, popgreen] at (10.5, 0.5) {Zone E/F\\Marge Passive\\Cameroun (Limbé)};
\end{tikzpicture}
\legende{Profil bathymétrique schématique Ouest-Est (Latitude ~4°N) traversant l'Atlantique équatorial. Échelle verticale exagérée.}
\end{popfigure}

\par
\textbf{Document 2 : Données d'âge du plancher océanique (échantillonnage par forage DSDP/ODP) sur le flanc Est de la dorsale (Zone D).}
\begin{center}
\begin{tabularx}{\linewidth}{|c|X|c|}\hline
\textbf{Échantillon} & \textbf{Distance à l'axe de la dorsale (km)} & \textbf{Âge radiométrique (Ma)} \\ \hline
1 (Axe) & 0 & 0 \\ \hline
2 & 50 & 2,0 \\ \hline
3 & 125 & 5,0 \\ \hline
4 & 250 & 10,0 \\ \hline
5 & 375 & 15,0 \\ \hline
\end{tabularx}
\end{center}

\par
\textbf{Document 3 : Diagramme Pression-Température (P-T) pour la fusion de la péridotite (manteau supérieur).}
\par
Le diagramme ci-dessous représente :
\begin{itemize}
    \item La \textbf{courbe du solidus anhydre} (péridotite sèche) : $T_{sol}^{anh} (^\circ\mathrm{C}) \approx 1085 + 132,9 P - 5,1 P^2$ (avec $P$ en GPa, valide 1--3 GPa).
    \item La \textbf{courbe du solidus hydraté} (péridotite + 0,1\% \ce{H2O} / \ce{CO2}) : approximation linéaire $T_{sol}^{hyd} \approx 950 + 120 P$ (valide 1--3 GPa).
    \item L'\textbf{adiabatique mantellique} (gradient $\approx \SI{0,3}{^\circ C/km}$, soit $\approx \SI{9}{^\circ C/GPa}$) pour un manteau potentiel à \SI{1350}{^\circ C} : $T_{ad}^{mantle} = 1350 + 9P$.
    \item L'\textbf{adiabatique d'un panache mantellique (hotspot)} pour un manteau potentiel à \SI{1550}{^\circ C} : $T_{ad}^{plume} = 1550 + 9P$.
\end{itemize}

\begin{popfigure}
\begin{tikzpicture}[scale=0.8]
% Axes
\draw[->, thick] (0,0) -- (8.5,0) node[right] {Température (°C)};
\draw[->, thick] (0,0) -- (0,5) node[above] {Pression (GPa)};
% Grid
\foreach \x in {1000,1100,...,1700} \draw[popline, very thin] (\x/200,0) -- (\x/200,4.5);
\foreach \y in {0.5,1,...,4} \draw[popline, very thin] (5,\y) -- (8.5,\y);
% Labels
\foreach \x in {1000,1200,1400,1600} \node[below] at (\x/200,-0.2) {\small \x};
\foreach \y in {1,2,3,4} \node[left] at (-0.2,\y) {\small \y};

% Solidus Anhydre (parabolic approx)
\draw[thick, popblue, domain=1:3.5, samples=50, variable=\p] plot ({ (1085 + 132.9*\p - 5.1*\p*\p)/200 }, {\p});
\node[popblue, rotate=30] at (5.5, 2.5) {Solidus anhydre};

% Solidus Hydraté (linear approx)
\draw[thick, poporange, domain=1:3.5, samples=20, variable=\p] plot ({ (950 + 120*\p)/200 }, {\p});
\node[poporange, rotate=30] at (5.0, 2.0) {Solidus hydraté (0,1\% volatils)};

% Mantle Adiabat (Tpot = 1350)
\draw[thick, popgreen, dashed, domain=0:4, samples=20, variable=\p] plot ({ (1350 + 9*\p)/200 }, {\p});
\node[popgreen, rotate=20] at (7.2, 3.5) {Adiabatique mantellique (1350°C)};

% Plume Adiabat (Tpot = 1550)
\draw[thick, poppurple, dashed, domain=0:4, samples=20, variable=\p] plot ({ (1550 + 9*\p)/200 }, {\p});
\node[poppurple, rotate=20] at (8.0, 3.5) {Adiabatique panache (1550°C)};

% Depth scale on right
\foreach \y/\d in {1/30, 2/60, 3/90, 4/120} \node[right, font=\small] at (8.6,\y) {\d km};
\node[right, font=\small, rotate=90] at (8.9, 2.5) {Profondeur approx.};
\end{tikzpicture}
\legende{Diagramme P-T de fusion de la péridotite. L'intersection adiabatique/solidus détermine la profondeur de début de fusion.}
\end{popfigure}

\par
\textbf{Document 4 : Catalogue des catastrophes naturelles récentes au Cameroun (extrait).}
\begin{center}
\begin{tabularx}{\linewidth}{|l|X|X|}\hline
\textbf{Événement} & \textbf{Localisation / Date} & \textbf{Caractéristiques \& Bilan} \\ \hline
Éruption volcanique & Mont Cameroun, 1999 (mars) \& 2000 (mai-juin) & Coulées de lave basaltique fluide (type hawaïen) atteignant \SI{4}{km} de la ville de Buea et l'océan. Fissures éruptives à \SI{2650}{m} et \SI{1400}{m} d'altitude. Émissions de \ce{SO2}, \ce{CO2}. Aucune victime directe, destruction cultures/habitats. \\ \hline
Catastrophe limnique & Lac Nyos (Région Nord-Ouest), 21 août 1986 & Dégazage brutal de $\approx \num{1,6} \times 10^6$ tonnes de \ce{CO2} dissous. Nuée dense (densité $> air$) s'écoulant dans vallées sur \SI{25}{km}. \textbf{1746 morts} (asphyxie), 3500 têtes bétail. Déclencheur probable : glissement de terrain interne ou injection magmatique. \\ \hline
Glissements de terrain & Bafoussam (Région Ouest), 28 octobre 2019 & Après pluies intenses (\SI{>100}{mm/24h}). Rupture de versant sur sols latéritiques épaissis (\SI{>20}{m}), pente \SI{>35}{\degree}. \textbf{42 morts}, centaines de sans-abri. Facteurs : déforestation, urbanisation non régulée, saturation en eau. \\ \hline
Inondations urbaines & Yaoundé (Région Centre), Octobre 2019 / 2020 / 2022 & Crues éclairs des rivières Mfoundi, Mingoa, Biyeme. Pluviométrie \SI{>150}{mm/jour}. Occupation lit majeur, colmatage drains, déchets plastiques. Dégâts matériels massifs, pertes en vies humaines (électrocution, noyade), épidémies choléra post-crue. \\ \hline
Sismicité & Région de l'Adamaoua / Monts Mandara / Côte (Kribi, Limbé) & Sismicité intraplaque modérée (M < 5,5). Failles réactivées (ligne du Cameroun, chaîne de fractures océaniques). Risque amplifié par vulnérabilité bâti (parpaings non ferraillés). \\ \hline
\end{tabularx}
\end{center}

\courssub{Consignes}
\par
En groupe de 4 « experts » (Géologue marin, Volcanologue/Géochimiste, Géomorphologue/Risques, Chargé de communication), répondez aux questions suivantes pour préparer l'affiche. Chaque membre est responsable d'une partie mais la note est collective.

\begin{enumerate}
    \item \textbf{(Géologue marin)} À partir du Document 1 :
    \begin{enumerate}
        \item Identifiez et nommez les structures géologiques présentes dans les Zones A, B, C, D, E, F.
        \item Quel type de limite de plaque sépare la Zone A de la Zone D ? Justifiez par la morphologie et la symétrie (ou asymétrie) du profil.
        \item Nommez les deux plaques lithosphériques concernées de part et d'autre de la Zone C.
    \end{enumerate}
    \item \textbf{(Géologue marin / Quantitatif)} À partir du Document 2 :
    \begin{enumerate}
        \item Représentez graphiquement l'âge en fonction de la distance (papier millimétré ou tableur).
        \item Calculez le \textbf{taux d'expansion demi-plaque} (en cm/an) à partir de la pente de la droite de régression. Précisez la méthode de calcul.
        \item En déduire le \textbf{taux d'expansion pleine} (vitesse de divergence des deux plaques). Ce taux est-il compatible avec une dorsale à expansion lente, intermédiaire ou rapide ?
    \end{enumerate}
    \item \textbf{(Volcanologue/Géochimiste)} À partir du Document 3 :
    \begin{enumerate}
        \item Expliquez pourquoi l'adiabatique mantellique « normale » (1350°C) ne coupe pas le solidus anhydre à haute pression ($> \SI{2}{GPa}$), alors qu'elle génère du magma (MORB) sous la dorsale. \textit{Indice : décompression adiabatique.}
        \item Vérifiez qu'à la pression de \SI{2,5}{GPa} (profondeur $\approx \SI{80}{km}$), l'adiabatique du panache mantellique (1550°C) est bien au-dessus du solidus hydraté. Calculez les deux températures à cette pression et en déduisez que la fusion a débuté à plus grande profondeur (pression plus élevée).
        \item Concluez sur l'origine du magma de la Ligne du Cameroun (volcanisme intraplaque) : rôle de la température (panache) et de la composition (volatils dans la lithosphère métasomatisée). Pourquoi ce volcanisme est-il alcalin (basanites, phonolites) et non tholéiitique comme les MORB ?
    \end{enumerate}
    \item \textbf{(Géomorphologue/Risques)} À partir du Document 4 et de vos connaissances :
    \begin{enumerate}
        \item Classez les 5 événements en 4 familles de risques : \emph{Volcanique}, \emph{Limnique}, \emph{Gravitaire}, \emph{Hydrologique/Météorologique}, \emph{Sismique}.
        \item Pour l'éruption du Mont Cameroun (1999/2000) : identifiez le type de dynamisme, la dangerosité principale pour Limbé (coulées, gaz, lahars), et le temps de préavis probable (prédiction vs prévision).
        \item Pour le Lac Nyos : expliquez le mécanisme physique de la catastrophe (stabilité de la colonne d'eau, rôle de la densité \ce{CO2}/eau). Quelle mesure de prévention technique a été installée depuis 2001 ?
        \item Pour les glissements de Bafoussam et inondations de Yaoundé : citez 2 facteurs anthropiques aggravants communs.
    \end{enumerate}
    \item \textbf{(Synthèse collective - Production de l'affiche)} Réunissez vos résultats pour concevoir l'\textbf{Affiche de Sensibilisation « Limbé : Connaître les risques pour mieux s'en protéger »}. L'affiche doit comporter obligatoirement :
    \begin{itemize}
        \item Une carte schématique simplifiée (Limbé, Mont Cameroun, Océan, Fosse, Dorsale) avec flèches de mouvement des plaques.
        \item Un encadré « \textbf{Les 4 risques majeurs à Limbé} » (Volcanique, Sismique, Tsunami, Glissement/Inondation) avec pour chacun : \emph{Origine géologique}, \emph{Signes précurseurs (Prédiction)}, \emph{Conduite à tenir (Prévention)}.
        \item Un encadré « \textbf{Numéros d'urgence \& Refuges} » (réalistes : Protection Civile, Hôpital, Mairie, Sites d'hébergement en hauteur).
        \item Un slogan accrocheur et une illustration pédagogique (schéma coupe volcan, ou panache mantellique).
    \end{itemize}
    \item \textbf{Restitution orale} : Présentation de l'affiche (5 min/groupe) suivie d'une évaluation croisée (grille fournie par l'enseignant : rigueur scientifique, clarté, utilité pratique, esthétique).
\end{enumerate}

\courssub{Ressources à mobiliser}
\begin{propriete}[Taux d'expansion océanique]
Si $d$ est la distance à l'axe (km) et $t$ l'âge (Ma), le taux d'expansion demi-plaque $v_{1/2} = \frac{d}{t}$ (en km/Ma). Conversion : \textbf{1 km/Ma = 0.1cm/an}. Le taux plein $v = 2 \times v_{1/2}$.
\begin{itemize}
    \item Lente : $v < \SI{4}{cm/an}$ (ex: Atlantique Nord ~2--3 cm/an).
    \item Intermédiaire : \SI{4}{cm/an} $< v <$ \SI{9}{cm/an}.
    \item Rapide : $v > \SI{9}{cm/an}$ (ex: Pacifique Est ~14 cm/an).
\end{itemize}
\end{propriete}

\begin{propriete}[Fusion par décompression (Rift / Dorsale)]
La remontée adiabatique du manteau (conservation de l'entropie) abaisse la pression sans perte de chaleur sensible. La température du manteau reste \emph{supérieure} au solidus à basse pression $\rightarrow$ \textbf{fusion partielle} (\% fusion $\approx$ 10--20\% pour MORB). Équation adiabatique : $\frac{dT}{dP} = \frac{\alpha T}{\rho C_p} \approx \SI{9}{^\circ C/GPa}$.
\end{propriete}

\begin{propriete}[Fusion par flux de volatils / Panache (Intraplaque / Ligne du Cameroun)]
L'apport de volatils (\ce{H2O}, \ce{CO2}) abaisse drastiquement le solidus (solidus hydraté). Un panache chaud (excès $T \approx \SI{200}{^\circ C}$) remontant sous une lithosphère métasomatisée (enrichie en volatils) croise le solidus hydraté à plus grande profondeur ($\sim$ 80--100 km) $\rightarrow$ faible degré de fusion (\% fusion $< 5\%$) $\rightarrow$ magmas \textbf{alcalins} (riches en \ce{Na2O}+\ce{K2O}, \ce{La}/\ce{Yb} élevé).
\end{propriete}

\begin{definition}[Prévision vs Prédiction vs Prévention]
\begin{itemize}
    \item \textbf{Prévision :} Probabilité statistique d'occurrence sur le long terme (ex: carte d'aléa, période de retour). « \emph{Où et quand (statistiquement) ?} »
    \item \textbf{Prédiction :} Annonce précise (date, lieu, magnitude) à court terme grâce à des précurseurs (sismicité, déformation, gaz). « \emph{Alerte imminente.} » Difficile pour séismes, possible pour volcans/lavas/crues.
    \item \textbf{Prévention :} Mesures structurelles (règles parasismiques, digues, dégazage Nyos) et non structurelles (plan d'évacuation, exercices, assurance, information) pour réduire la \textbf{vulnérabilité}.
\end{itemize}
\end{definition}

\begin{methode}[Calcul du taux d'expansion à partir de points expérimentaux]
\begin{enumerate}
    \item Reporter les points $(d_i, t_i)$ dans un repère orthonormé.
    \item Tracer la droite de régression linéaire (ou droite passant par l'origine et le point le plus éloigné si relation proportionnelle).
    \item Calculer la pente $p = \frac{\Delta d}{\Delta t}$ (en km/Ma).
    \item Convertir : $v_{1/2} (\text{cm/an}) = p \times 0,1$.
    \item $v_{\text{plein}} = 2 \times v_{1/2}$.
\end{enumerate}
\end{methode}

\courssub{Démarche et corrigé commenté}
\begin{exoresolu}[Corrigé de l'activité d'intégration : Commission d'experts Limbé]
\tcblower
\textbf{1. Analyse du profil bathymétrique (Doc 1)}
\begin{itemize}
    \item \textbf{Zone A \& D :} \textbf{Plaines abyssales} (croûte océanique vieille, dense, isotherme, recouverte de sédiments pélagiques). Profondeur \SI{\approx 4800}{m}.
    \item \textbf{Zone B :} \textbf{Zones de fracture (failles transformantes)}. Vallées et crêtes orientées perpendiculairement à la dorsale, décalant l'axe. Sismicité superficielle ($< \SI{15}{km}$), mécanismes au foyer de type \emph{coulissage}. Permettent l'expansion différentielle sur une sphère.
    \item \textbf{Zone C :} \textbf{Dorsale médio-atlantique (Zone d'accrétion / Rift océanique).} Vallée de rift centrale (graben effondré), axe magmatique (coulées de coussins, cheminées hydrothermales). Anomalie thermique positive, croûte fine (\SI{\approx 6}{km}), flux de chaleur élevé.
    \item \textbf{Zone E :} \textbf{Pente continentale} (marque la transition croûte océanique / croûte continentale amincie).
    \item \textbf{Zone F :} \textbf{Plateau continental} (croûte continentale, sédiments épais), côte de Limbé.
    \item \textbf{Limite de plaque :} \textbf{Divergente (accrétion)}. La symétrie grossière du profil de part et d'autre de la Zone C (plaines abyssales à même profondeur, flancs symétriques) confirme l'expansion océanique symétrique.
    \item \textbf{Plaques :} \textbf{Plaque sud-américaine} (Ouest / Zone A) et \textbf{Plaque africaine} (Est / Zones D, E, F).
\end{itemize}

\textbf{2. Calcul du taux d'expansion (Doc 2)}
\begin{itemize}
    \item \textbf{Graphique :} Points alignés passant par l'origine. Relation proportionnelle $d = k \cdot t$.
    \item \textbf{Calcul :} Utilisons les points extrêmes pour minimiser l'erreur relative : $(d=375 \text{ km}, t=15 \text{ Ma})$.
    \[ v_{1/2} = \frac{375 \text{ km}}{15 \text{ Ma}} = 25 \text{ km/Ma} \]
    Conversion : $25 \text{ km/Ma} \times 0,1 = \mathbf{2,5 \text{ cm/an}}$ (taux demi-plaque).
    \item \textbf{Taux plein :} $v = 2 \times 2,5 = \mathbf{5,0 \text{ cm/an}}$.
    \item \textbf{Classification :} \SI{5,0}{cm/an} $\rightarrow$ \textbf{Expansion intermédiaire} (proche de la limite lente). Compatible avec la dorsale médio-atlantique équatoriale.
\end{itemize}

\textbf{3. Origine des magmas : Ligne du Cameroun vs Dorsale (Doc 3)}
\begin{itemize}
    \item \textbf{Dorsale (MORB) :} L'adiabatique mantellique ($T_{pot}=1350^\circ\mathrm{C}$) ne coupe pas le solidus anhydre à haute pression ($> \SI{2,3}{GPa}$). Mais la \textbf{remontée adiabatique} sous la dorsale suit une trajectoire $T(P)$ qui, à partir de \SI{\approx 70}{km} (\SI{\approx 2,3}{GPa}), se retrouve \emph{au-dessus} du solidus anhydre car la pression chute plus vite que la température (décompression). $\rightarrow$ Fusion partielle étendue (\SI{10-20}{\%}), magmas tholéiitiques (MORB).
    \item \textbf{Panache / Ligne du Cameroun :}
    \begin{itemize}
        \item Adiabatique panache : $T_{ad}^{plume}(P) = 1550 + 9P$ ($P$ en GPa, $T$ en $^\circ\mathrm{C}$).
        \item Solidus hydraté (linéaire 1-3 GPa) : $T_{sol}^{hyd}(P) \approx 950 + 120P$.
        \item Intersection (début de fusion) : $1550 + 9P = 950 + 120P \Rightarrow 111P = 600 \Rightarrow P \approx \mathbf{5,4 \text{ GPa}}$ ($\sim \SI{160}{km}$).
        \item \textbf{À 2,5 GPa (80 km) :}
        $T_{ad}^{plume} = 1550 + 9 \times 2,5 = \mathbf{1572,5 ^\circ\mathrm{C}}$.
        $T_{sol}^{hyd} = 950 + 120 \times 2,5 = \mathbf{1250 ^\circ\mathrm{C}}$.
        L'adiabatique du panache est \textbf{nettement au-dessus} du solidus hydraté (\SI{322,5}{^\circ C} d'excès). La fusion a donc débuté plus profondément (vers 160 km) et se poursuit lors de la remontée.
    \end{itemize}
    \item \textbf{Nature alcaline :} Faible degré de fusion partielle ($F < 5\%$) $\rightarrow$ fusion enrichie en éléments incompatibles (K, Na, La, Ce) $\rightarrow$ \textbf{Basanites, Tephrites, Phonolites} (série alcaline). À la dorsale, $F \approx 15\%$ $\rightarrow$ appauvrissement en incompatibles $\rightarrow$ \textbf{Tholéiites} (MORB).
\end{itemize}

\textbf{4. Analyse des risques camerounais (Doc 4)}
\begin{itemize}
    \item \textbf{Classification :}
    \begin{tabular}{ll}
    Volcanique : & Mont Cameroun (1999, 2000). \\
    Limnique : & Lac Nyos (1986) — \emph{phénomène limnique déclenché par gravitaire/magmatique}. \\
    Gravitaire : & Glissements Bafoussam (2019). \\
    Hydrologique / Météorologique : & Inondations Yaoundé (2019, 2020, 2022). \\
    Sismique : & Sismicité intraplaque (Adamaoua, Mandara, Côte).
    \end{tabular}
    \item \textbf{Mont Cameroun :} Dynamisme \textbf{hawaïen / strombolien} (lave basaltique fluide, fontaines de lave, coulées rapides \SI{>10}{km/h} sur pente raide). Dangerosité Limbé : \textbf{Coulées de lave} (atteignent l'océan en \SI{\approx 6-12}{h} depuis flanc SW), \textbf{Gaz} (\ce{CO2} accumulations en bas de pente, \ce{SO2} pluies acides), \textbf{Lahars} (remobilisation cendres par pluies tropicales). Prédiction : Réseau sismique (OVSM/IRGM) + GPS + gaz $\rightarrow$ préavis \textbf{jours à semaines} (ex: 1999 : crise sismique 2 semaines avant).
    \item \textbf{Lac Nyos :} Lac de cratère méromictique (stratifié permanent). \ce{CO2} d'origine magmatique (dégazage profond) piégé en fond (haute pression, basse $T$). \textbf{Déclencheur} : glissement de berge interne ou injection magmatique $\rightarrow$ mélange eaux profondes/superficielles $\rightarrow$ chute pression $\rightarrow$ exsolution brutale \ce{CO2} (ébullition) $\rightarrow$ vague de dégazage. \ce{CO2} plus dense que l'air ($\rho_{\ce{CO2}} \approx 1,98 \text{ kg/m}^3$ vs $\rho_{air} \approx 1,2$) $\rightarrow$ écoulement gravitaire dans vallées. \textbf{Prévention technique :} Installation depuis 2001 de \textbf{tuyaux de dégazage auto-amorçants} (pompage eaux profondes $\rightarrow$ libération \ce{CO2} contrôlée en surface). 3 tuyaux opérationnels (2019).
    \item \textbf{Facteurs anthropiques aggravants (Bafoussam \& Yaoundé) :}
    \begin{enumerate}
        \item \textbf{Occupation anarchique des zones à risque} : construction sur versants instables (pente $>30^\circ$, sols latéritiques) et en lit majeur de rivières (zone d'expansion de crue).
        \item \textbf{Dégradation de l'environnement / Gestion des déchets} : Déforestation (perte cohésion racinaire, ruissellement accru) + colmatage réseaux d'assainissement par déchets plastiques (ralentissement écoulement, montée eaux plus rapide).
    \end{enumerate}
\end{itemize}

\textbf{5. Production de l'Affiche « Limbé : Connaître les risques pour mieux s'en protéger » (Synthèse attendue)}

\par
\textbf{Structure type de l'affiche (contenu minimal exigible) :}

\begin{center}
\begin{tabularx}{\linewidth}{|X|X|}\hline
\textbf{HAUT GAUCHE : CARTE SCHÉMATIQUE} & \textbf{HAUT DROIT : SLOGAN} \\
\begin{tikzpicture}[scale=0.6]
% Coastline
\draw[thick, popblue] (0,0) .. controls (1,0.2) and (2,-0.2) .. (3,0) .. controls (4,0.2) and (5,-0.1) .. (6,0);
\draw[thick, popblue] (0,-0.1) .. controls (1,0.1) and (2,-0.3) .. (3,-0.1) .. controls (4,0.1) and (5,-0.2) .. (6,-0.1);
% Ocean label
\node[popblue, font=\small] at (3, 0.8) {Océan Atlantique};
% Limbé
\fill[popred] (2.5, -0.05) circle (2pt) node[below right, font=\small, popdark] {Limbé ★};
% Mount Cameroon
\draw[fill=poporange!50!popdark, thick] (1.5,-1.5) -- (2.0,0.2) -- (2.5,-1.5) -- cycle;
\node[popdark, font=\small] at (2.0, 0.5) {Mont Cameroun ▲};
% Ridge far West
\draw[poporange, decorate, decoration={snake, amplitude=1pt, segment length=4pt}] (0,1.5) -- (6,1.5);
\node[poporange, font=\small, rotate=0] at (3, 1.8) {Dorsale Médio-Atlantique};
% Plate arrows
\draw[->, thick, popgreen] (0.5, 1.2) -- (0.5, 0.5) node[midway, left] {Plaque Am. Sud};
\draw[->, thick, popgreen] (5.5, 1.2) -- (5.5, 0.5) node[midway, right] {Plaque Afrique};
% Transform faults
\draw[dashed, popmuted] (1,1.5) -- (1,-1.5);
\draw[dashed, popmuted] (4,1.5) -- (4,-1.5);
\node[popmuted, font=\tiny, rotate=90] at (0.7, 0) {Faille transformante};
\end{tikzpicture} &
\fbox{\parbox{\linewidth}{\centering \textbf{« Limbé, ville résiliente : \\ Je m'informe, Je me prépare, Je survive ! »}}} \\
\footnotesize Carte : Limbé (★), Mont Cameroun (▲), Océan, Dorsale (~~), Failles transformantes (---). Flèches plaques : Afrique (→), Am. Sud (←). & \\ \hline
\multicolumn{2}{|c|}{\textbf{ENCADRÉ CENTRAL : LES 4 RISQUES MAJEURS À LIMBÉ}} \\ \hline
\textbf{1. VOLCANIQUE (Mont Cameroun)} & \textbf{2. SISMIQUE (Failles transformantes / Panache)} \\
\emph{Origine :} Panache mantellique, faille crustale (Ligne du Cameroun). & \emph{Origine :} Réactivation failles (ligne Cameroun, chaîne fracture océanique). \\
\emph{Précurseurs (Prédiction) :} Essaims sismiques (jours/semaines), déformation GPS, hausse \ce{SO2}/\ce{CO2}. & \emph{Précurseurs :} Pas de prédiction fiable (court terme). Aléa modéré (M<5,5). \\
\emph{Conduite (Prévention) :} \textbf{ÉVACUATION ORDONNÉE} vers zones hautes (Buea, Muyuka) \textbf{AVANT} arrivée coulées. Masque humide (gaz). Ne pas traverser coulées. & \emph{Conduite :} \textbf{À l'intérieur :} Se mettre sous table solide (PLIER-COUVRIR-TENIR). \textbf{À l'extérieur :} S'éloigner bâtiments, fils électriques. Sac d'urgence prêt. \\ \hline
\textbf{3. TSUNAMI (Séisme dorsal / Glissement sous-marin)} & \textbf{4. GLISSEMENT / INONDATION (Pluies extrêmes)} \\
\emph{Origine :} Séisme dorsale (M>7) ou effondrement flanc volcan sous-marin. & \emph{Origine :} Sols latéritiques saturés + pentes raides + urbanisation / déchets. \\
\emph{Précurseurs :} \textbf{Recul brutal de la mer} (signe majeur), secousse forte longue (>20s). Alerte PTWC relayée. & \emph{Précurseurs :} Vigilance météo (METEO CAMEROUN), pluie > 50mm/h, ruissellement inhabituel, fissures sol. \\
\emph{Conduite :} \textbf{FUYEZ VERS LES HAUTEURS IMMÉDIATEMENT} (ne pas attendre sirène). \textbf{Ne pas aller voir la vague.} Rester en hauteur jusqu'à levée alerte. & \emph{Conduite :} \textbf{Ne pas traverser routes inondées} (30 cm eau = voiture emportée). Couper électricité/gaz. Rejoindre centre d'hébergement (Mairie, École, Église). \\ \hline
\multicolumn{2}{|c|}{\textbf{BAS : NUMÉROS D'URGENCE \& REFUGES}} \\ \hline
\multicolumn{2}{|c|}{\parbox{\linewidth}{\centering
\textbf{📞 URGENCES :} 112 (Standard européen/GSM) | 119 (Sapeurs-Pompiers) | 118 (Police) | Hôpital Laquintinie Douala : +237 233 43 00 00 \\
\textbf{🏠 SITES DE REPLI (Hauteur > 50m) :} Lycée de Limbé (Mile 1) | Stade municipal | Campus UB (Molyko, Buea) | Église Presbytérienne (Down Beach) \\
\textbf{🎒 MON SAC D'URGENCE (Prêt à l'entrée) :} Eau (2L/jour/pers), Conserves, Lampe torche + piles, Radio piles, Médicaments, Papiers d'identité (sac étanche), Sifflet, Couteau, Couverture survie.
}} \\ \hline
\end{tabularx}
\end{center}

\par
\textbf{Illustration pédagogique suggérée (au centre ou bas) :} Coupe schématique Nord-Sud du Mont Cameroun montrant : Chambre magmatique profonde (panache), conduit d'alimentation, fissures éruptives 1999/2000, coulées vers Limbé, panache de gaz \ce{CO2} (densité > air) s'accumulant en bas de pente (risque asphyxie silencieuse), zone d'évacuation (flèches vers Buea/Muyuka).

\textbf{6. Grille d'évaluation croisée (Extrait pour auto-évaluation)}
\begin{center}
\begin{tabularx}{\linewidth}{|l|X|c|}\hline
\textbf{Critère} & \textbf{Indicateurs de réussite} & \textbf{/5} \\ \hline
Rigueur scientifique & Structures dorsale nommées, taux expansion calculé (2,5 cm/an demi), origine magma (panache + volatils) expliquée, mécanismes risques justifiés. & \\
Clarté \& Lisibilité & Titre visible, carte lisible, 4 risques distincts, pictogrammes, police grande taille, couleurs contrastées (Rouge danger, Vert sécurité, Bleu info). & \\
Utilité pratique & Numéros réels, refuges nommés localement, conduites à tenir concrètes (ex: « Ne pas traverser coulée »), sac d'urgence listé. & \\
Esthétique / Originalité & Mise en page aérée, slogan accrocheur, illustration schématique explicative (coupe volcan), langue française correcte (+ termes pidgin/anglais si pertinent). & \\
Travail d'équipe & Répartition rôles respectée, cohérence entre parties, présentation orale fluide (5 min), réponses aux questions pertinentes. & \\ \hline
\end{tabularx}
\end{center}
\end{exoresolu}

\courssub{Bilan de l'activité}
\par
Cette activité d'intégration a permis de mettre en évidence la \textbf{continuité entre la géodynamique profonde (panache mantellique, expansion océanique) et les risques de surface vécus par les populations camerounaises}.
\begin{aretenir}[Synthèse des apprentissages mobilisés]
\begin{itemize}
    \item \textbf{Géodynamique :} La morphologie des fonds océaniques (Doc 1) enregistre l'histoire de l'expansion (taux \SI{5}{cm/an} plein, Doc 2). La dorsale médio-atlantique est une limite divergente ; la Ligne du Cameroun est une structure intraplaque liée à un panache mantellique interagissant avec une lithosphère métasomatisée (volatils $\rightarrow$ solidus hydraté abaissé, Doc 3).
    \item \textbf{Pétrologie / Géochimie :} Le même manteau (péridotite) produit des magmas différents selon le contexte thermo-barométrique : \textbf{MORB tholéiitiques} (fusion décompression extensive, dorsale) vs \textbf{Basanites/Phonolites alcalines} (fusion faible degré, panache + volatils, Ligne du Cameroun).
    \item \textbf{Gestion des risques :} La distinction \textbf{Prévision / Prédiction / Prévention} est opératoire.
    \begin{itemize}
        \item \emph{Volcanisme :} Prédiction possible (précurseurs) $\rightarrow$ Évacuation.
        \item \emph{Séismes :} Prédiction impossible $\rightarrow$ Prévention parasismique + Conduite « Plier-Couvrir-Tenir ».
        \item \emph{Tsunami :} Prédiction courte (alerte sismique + recul mer) $\rightarrow$ Fuite immédiate en hauteur.
        \item \emph{Glissements/Inondations :} Prévision météo + Aménagement (non-construction zones à risque, drains, reboisement) + Alerte précoce.
    \end{itemize}
    \item \textbf{Citoyenneté scientifique :} Transformer des données brutes (bathymétrie, géochronologie, diagrammes P-T, catalogues catastrophes) en \textbf{message clair, actionnable, localisé} (affiche Limbé) est la finalité de l'éducation au développement durable. La résilience se construit par la connaissance partagée.
\end{itemize}
\end{aretenir}

\begin{attention}[Pièges fréquents à éviter dans le corrigé / l'affiche]
\begin{itemize}
    \item \textbf{Confondre taux demi-plaque et taux plein :} Ne pas oublier le facteur 2. Un taux de \SI{2,5}{cm/an} demi-plaque donne \SI{5}{cm/an} plein (intermédiaire), pas « lent ».
    \item \textbf{Mélanger dorsale et point chaud :} La dorsale (Zone C) produit des MORB par décompression \emph{sans} apport de volatils majeur (solidus anhydre). La Ligne du Cameroun nécessite l'apport de volatils (solidus hydraté) \emph{ET} un excès thermique (panache) pour fondre à faible profondeur.
    \item \textbf{Oublier le \ce{CO2} au Mont Cameroun :} Le risque gaz (\ce{CO2} dense, asphyxie silencieuse en bas de pente, type Nyos mais diffus) est sous-estimé par les élèves focalisés sur la lave. L'affiche doit le mentionner.
    \item \textbf{Conduites à tenir génériques :} « Fuir » ne suffit pas. Préciser : « Vers où ? » (Hauteurs identifiées), « Comment ? » (Voies dégagées, piéton), « Avec quoi ? » (Sac d'urgence).
    \item \textbf{Unités :} Vitesse en cm/an (pas km/Ma sans conversion), Pression en GPa, Température en $^\circ\mathrm{C}$, Profondeur en km.
\end{itemize}
\end{attention}

\begin{savaistu}[Le saviez-vous ? Le dégazage du Lac Nyos, une réussite technique camerounaise]
\par
Après la catastrophe de 1986, une équipe internationale (dont des chercheurs de l'IRGM Yaoundé et de l'Université de Savoie) a conçu un système de \textbf{dégazage par tuyaux auto-amorçants (siphons)}. Le premier tuyau (diamètre \SI{140}{mm}, longueur \SI{200}{m}) a été installé en \textbf{2001}. L'eau profonde, riche en \ce{CO2} dissous, monte dans le tuyau ; la pression diminuant, le gaz exsolutionne, forme des bulles qui tirent l'eau vers le haut (effet « geyser » continu). Deux autres tuyaux ont été ajoutés en 2011. Le lac est désormais surveillé en temps réel (capteurs \ce{CO2}, température, profil chimique) et le risque d'éruption limnique majeure est \textbf{maîtrisé}, tant que les tuyaux fonctionnent. C'est un exemple unique au monde de \textbf{prévention active d'un risque géologique majeur} par l'ingénierie géochimique.
\end{savaistu}

\newpage

\coursec{Méthodes et exercices résolus}

\coursec{Réduction des risques liés aux catastrophes naturelles}

\courssub{Analyser les données topographiques des fonds océaniques}

\begin{methode}[Lire et interpréter un profil topographique océanique]
\begin{enumerate}
\item \textbf{Repérer les axes et les échelles} : profondeur en ordonnée (en m ou km, souvent négative ou croissante vers le bas), distance en abscisse (en km). Vérifier si l'échelle verticale est exagérée (fréquent pour visualiser le relief).
\item \textbf{Identifier les grandes structures} de l'océan vers le continent : la \cle{dorsale médio-océanique} (relief sous-marin élevé, creusé d'un \cle{rift} axial), la \cle{plaine abyssale} (fond plat vers \num{-4000} à \num{-6000} m), la \cle{fosse océanique} (dépression très profonde, jusqu'à \num{-11000} m), puis le \cle{talus continental}, le \cle{plateau continental} et le continent.
\item \textbf{Localiser le rift} : vallée axiale au sommet de la dorsale, encadrée de reliefs symétriques.
\item \textbf{Interpréter} : la dorsale est une zone d'\cle{accrétion océanique} (formation de lithosphère neuve) ; la fosse est une zone de \cle{subduction} (disparition de lithosphère). La symétrie des reliefs de part et d'autre du rift traduit l'expansion océanique.
\item \textbf{Conclure} en reliant chaque structure au fonctionnement des plaques lithosphériques.
\end{enumerate}
\textbf{Réflexes} : profondeur croissante vers le bas ; dorsale = divergence ; fosse = convergence ; toujours citer des valeurs chiffrées lues sur le document.
\end{methode}

\begin{popfigure}
\begin{tikzpicture}[scale=0.8, every node/.style={font=\small}]
% Axes
\draw[->, popdark] (0,0) coordinate (O) -- (14,0) node[right] {Distance (km)};
\draw[->, popdark] (0,-6) -- (0,1) node[above] {Profondeur (m)};
% Graduations profondeur
\foreach \y/\ylab in {0/0, -1/-1000, -2/-2000, -3/-3000, -4/-4000, -5/-5000, -6/-6000}
  \draw (0,\y) -- (-0.2,\y) node[left] {\ylab};
% Profil Afrique (gauche)
\draw[popblue, very thick] (0,0) -- (1,0) node[midway, above=2pt] {Plateau cont.} 
  .. controls (1.5,-0.5) and (2,-2) .. (2.5,-3) node[midway, right=2pt] {Talus}
  -- (5,-5) node[midway, above=2pt] {Plaine abyssale}
  .. controls (6,-4.5) and (6.5,-3) .. (7,-2) node[midway, above=2pt] {Dorsale}
  .. controls (7.5,-3) and (8,-4.5) .. (8.5,-5) 
  -- (11,-5) 
  .. controls (11.5,-3) and (12,-1) .. (12.5,0) node[midway, right=2pt] {Talus}
  -- (14,0) node[midway, above=2pt] {Plateau cont.};
% Rift axial
\draw[poporange, thick, dashed] (6.8,-2.2) -- (7.2,-2.2) node[right, poporange] {Rift axial};
% Annotations
\node[popblue, anchor=south] at (0.5,0.3) {Afrique};
\node[popblue, anchor=south] at (13.5,0.3) {Amérique du Sud};
\node[poppurple, anchor=north] at (7,-1.5) {Dorsale médio-atlantique};
\node[poppurple, anchor=north] at (3.5,-5.3) {Plaine abyssale (\SI{-5000}{\meter})};
\node[poppurple, anchor=north] at (9.5,-5.3) {Plaine abyssale (\SI{-5000}{\meter})};
\end{tikzpicture}
\legende{Profil topographique schématique Est-Ouest de l'Atlantique Nord (échelle verticale exagérée). On distingue les plateaux continentaux, les talus, les plaines abyssales et la dorsale médio-atlantique creusée de son rift axial.}
\end{popfigure}

\begin{exoresolu}[Profil topographique de l'Atlantique Nord]
On donne un profil topographique Est--Ouest de l'océan Atlantique entre l'Afrique et l'Amérique du Sud. On relève, d'Est en Ouest : plateau continental africain (0 à \num{-200} m), talus continental, plaine abyssale vers \num{-5000} m, relief élevé atteignant \num{-2000} m avec une dépression axiale de \num{20} km de large, puis de nouveau plaine abyssale, talus et plateau continental sud-américain.
\begin{enumerate}
\item Identifiez les structures A (relief élevé), B (dépression axiale) et C (plaine abyssale).
\item Interprétez la structure A en termes de fonctionnement des plaques.
\end{enumerate}
\tcblower
\textbf{1. Identification des structures.}
\begin{itemize}
\item \textbf{A} est un relief sous-marin élevé (\num{-2000} m) situé au milieu de l'océan, à égale distance des deux continents : c'est la \textbf{dorsale médio-atlantique}.
\item \textbf{B}, la dépression axiale au sommet de A, est le \textbf{rift} : vallée d'effondrement axiale.
\item \textbf{C}, fond plat et profond (\num{-5000} m) de part et d'autre de la dorsale, est la \textbf{plaine abyssale}.
\end{itemize}
\textbf{2. Interprétation.} La dorsale est un relief volcanique axial. Sa position médiane et la symétrie des fonds de part et d'autre montrent qu'elle correspond à la zone où la lithosphère océanique se forme : le magma basaltique remonte au niveau du rift, se solidifie et s'écarte symétriquement de part et d'autre de l'axe. C'est une zone de \textbf{divergence} (accrétion océanique) : les plaques africaine et sud-américaine s'éloignent l'une de l'autre à une vitesse de l'ordre de 2 à \SI{3}{cm/an}.
\textbf{Conclusion} : le profil montre une dorsale médio-océanique creusée d'un rift, témoin de l'expansion du plancher océanique atlantique.
\end{exoresolu}

\courssub{Reconstituer la formation du rift et du plancher océanique}

\begin{methode}[Expliquer l'expansion océanique et réaliser une maquette du rift]
\begin{enumerate}
\item \textbf{Point de départ} : sous l'axe de la dorsale, la remontée de matériaux asthénosphériques chauds provoque une \cle{décompression}.
\item \textbf{Fusion} : cette décompression entraîne la fusion partielle des péridotites du manteau, produisant un magma basaltique.
\item \textbf{Injection et solidification} : le magma s'injecte dans la fissure axiale (rift) et cristallise en \cle{gabbros} en profondeur (refroidissement lent) et en \cle{basaltes} (coulées en coussins ou \emph{pillow-lavas}) en surface (refroidissement rapide) : c'est la croûte océanique neuve (modèle ophiolitique).
\item \textbf{Écartement} : la lithosphère neuve est repoussée symétriquement de part et d'autre de l'axe ; elle refroidit, s'épaissit (lithosphère = croûte + manteau lithosphérique) et s'enfonce isostatiquement (la profondeur augmente avec l'âge $\propto \sqrt{\text{âge}}$).
\item \textbf{Maquette} : représenter deux bandes de papier (plaques) sortant d'une fente (rift) ; colorier des bandes symétriques d'âges croissants ; dessiner en coupe : rift axial, chambre magmatique, gabbros, basaltes (coussins), sédiments qui s'épaississent en s'éloignant de l'axe.
\end{enumerate}
\textbf{Preuves à citer} : symétrie des anomalies magnétiques, âge croissant des basaltes avec la distance à l'axe, épaisseur croissante des sédiments, flux de chaleur maximal à l'axe.
\end{methode}

\begin{popfigure}
\begin{tikzpicture}[scale=0.9, every node/.style={font=\small}]
% Coupe dorsale
\draw[popdark, thick] (-5,-0.5) -- (5,-0.5); % Surface mer
% Eau
\fill[popblue!20] (-5,-0.5) rectangle (5,-3);
% Sédiments (épaississement)
\fill[popgold!50] (-5,-3) -- (-5,-3.3) -- (-1,-3.5) -- (-1,-3) -- cycle;
\fill[popgold!50] (1,-3) -- (1,-3.5) -- (5,-3.3) -- (5,-3) -- cycle;
% Basaltes (coussins)
\fill[poporange!50] (-5,-3.3) -- (-1,-3.5) -- (-1,-4) -- (-5,-3.8) -- cycle;
\fill[poporange!50] (1,-3.5) -- (5,-3.3) -- (5,-3.8) -- (1,-4) -- cycle;
% Gabbros
\fill[poppurple!30] (-5,-3.8) -- (-1,-4) -- (-1,-5) -- (-5,-4.5) -- cycle;
\fill[poppurple!30] (1,-4) -- (5,-3.8) -- (5,-4.5) -- (1,-5) -- cycle;
% Chambre magmatique
\fill[poppink!50] (-1,-4) rectangle (1,-5) node[midway, white, font=\bfseries, align=center]{Chambre\\magmatique};
% Rift
\draw[poporange, thick, decorate, decoration={zigzag, amplitude=2pt, segment length=4pt}] (-0.5,-3) -- (0.5,-3) node[midway, above=4pt, poporange, font=\bfseries] {Rift axial};
% Flèches écartement
\draw[<->, popblue, thick] (-3,-0.2) -- (-1,-0.2) node[midway, above] {Écartement};
\draw[<->, popblue, thick] (1,-0.2) -- (3,-0.2) node[midway, above] {Écartement};
% Légendes couches
\node[popgold, anchor=east] at (-5.2,-3.15) {Sédiments};
\node[poporange, anchor=east] at (-5.2,-3.6) {Basaltes (coussins)};
\node[poppurple, anchor=east] at (-5.2,-4.2) {Gabbros};
\node[poppink, anchor=east] at (-5.2,-4.7) {Manteau (péridotites)};
% Anomalies magnétiques (dessous)
\begin{scope}[yshift=-6.5cm]
\draw[popdark, thick] (-5,0) -- (5,0);
\foreach \x/\col in {-4.5/popblue, -3.5/poporange, -2.5/popblue, -1.5/poporange, -0.5/popblue, 0.5/poporange, 1.5/popblue, 2.5/poporange, 3.5/popblue, 4.5/poporange}
  \fill[\col] (\x,-0.3) rectangle (\x+0.9,0.3);
\draw[<->, popdark] (-4.5,-0.6) -- (-0.5,-0.6) node[midway, below] {Symétrie};
\draw[<->, popdark] (0.5,-0.6) -- (4.5,-0.6) node[midway, below] {Symétrie};
\node[popdark, font=\bfseries] at (0,-1) {Anomalies magnétiques (Normale / Inversée)};
\end{scope}
\end{tikzpicture}
\legende{Coupe schématique d'une dorsale médio-océanique (haut) : structure ophiolitique (sédiments, basaltes en coussins, gabbros, péridotites) et symétrie des anomalies magnétiques (bas) prouvant l'expansion océanique.}
\end{popfigure}

\begin{exoresolu}[Âge du plancher océanique et vitesse d'expansion]
Des dragages le long d'un profil perpendiculaire à la dorsale médio-atlantique donnent les âges des basaltes suivants : à \SI{50}{km} de l'axe, 3 Ma ; à \SI{100}{km}, 6 Ma ; à \SI{200}{km}, 12 Ma (les âges sont identiques de l'autre côté de l'axe).
\begin{enumerate}
\item Calculez la vitesse d'expansion (demi-vitesse d'écartement) de chaque côté de la dorsale.
\item Déduisez-en la vitesse d'écartement des deux plaques.
\item Expliquez la symétrie des âges.
\end{enumerate}
\tcblower
\textbf{1. Demi-vitesse d'expansion.} La vitesse est le rapport de la distance à l'axe sur l'âge de la croûte :
\[ v = \frac{d}{t} = \frac{\SI{100}{km}}{\SI{6}{Ma}} = \frac{\num{100000}\ \text{m}}{\num{6e6}\ \text{ans}} \approx \SI{1,67}{cm/an} \approx \SI{1,7}{cm/an} \]
Vérification avec les autres points : $\frac{\SI{50}{km}}{\SI{3}{Ma}} \approx \SI{1,7}{cm/an}$ et $\frac{\SI{200}{km}}{\SI{12}{Ma}} \approx \SI{1,7}{cm/an}$. La vitesse est constante.

\textbf{2. Vitesse d'écartement des plaques.} L'écartement total est la somme des deux demi-vitesses symétriques :
\[ v_{\text{total}} = 2 \times \SI{1,7}{cm/an} = \SI{3,4}{cm/an} \]

\textbf{3. Symétrie des âges.} La croûte océanique se forme en continu au niveau du rift axial, puis chaque moitié s'éloigne de l'axe à la même vitesse, de part et d'autre. Deux points situés à égale distance de l'axe se sont donc formés au même instant : les âges sont symétriques par rapport à l'axe de la dorsale. C'est la preuve directe de l'\textbf{expansion océanique}.
\textbf{Conclusion} : le plancher océanique s'étend à une vitesse totale d'environ \SI{3,4}{cm/an} à partir du rift médio-atlantique.
\end{exoresolu}

\courssub{Naissance des magmas : données expérimentales de fusion des péridotites}

\begin{methode}[Exploiter une courbe de fusion partielle des péridotites (diagramme P-T)]
\begin{enumerate}
\item \textbf{Lire le diagramme} : Température (\si{\celsius}) en ordonnée, Pression (GPa) ou Profondeur (km) en abscisse (souvent inversée : 0 en haut). Repérer le \cle{solidus anhydre} (courbe de début de fusion des péridotites sèches), le \cle{solidus hydraté} (abaissé par la présence d'eau) et le \cle{géotherme} (courbe de la température réelle du manteau selon le contexte).
\item \textbf{Condition de fusion} : il y a fusion partielle là où le géotherme \emph{croise ou dépasse} le solidus correspondant.
\item \textbf{Cas de la dorsale (accrétion)} : la remontée rapide de péridotites chaudes se fait presque sans échange de chaleur (adiabatique) ; la \textbf{pression diminue} (décompression) alors que la température reste élevée : le géotherme local croise le solidus anhydre $\Rightarrow$ fusion partielle par \textbf{décompression} $\Rightarrow$ magma \textbf{basaltique} (pauvre en silice, fluide, volcanisme effusif).
\item \textbf{Cas de la subduction} : la plaque plongeante transporte de l'eau (minéraux hydratés, sédiments). En profondeur, la déshydratation libère de l'eau qui métasomatose le manteau du coin mantellique. L'eau \textbf{abaisse fortement le solidus} (solidus hydraté) $\Rightarrow$ le géotherme du coin mantellique croise ce solidus abaissé $\Rightarrow$ fusion partielle $\Rightarrow$ magma \textbf{andésitique} (riche en silice, visqueux, volcanisme explosif).
\item \textbf{Conclure} : deux mécanismes distincts — décompression (dorsale) et hydratation (subduction) — expliquent la naissance des magmas primaires.
\end{enumerate}
\textbf{Réflexe} : ne jamais écrire « le manteau est fondu » ou « liquide ». Le manteau est \textbf{solide} (comportement ductile à long terme) ; seule une \emph{fusion partielle} (quelques pourcents, typiquement 5 à 20\%) des péridotites produit le magma.
\end{methode}

\begin{popfigure}
\begin{tikzpicture}[scale=0.8, every node/.style={font=\small}]
% Axes
\draw[->, popdark] (0,0) -- (8,0) node[right] {Température (\si{\celsius})};
\draw[->, popdark] (0,0) -- (0,6) node[above] {Profondeur (km)};
% Inverser l'axe Y visuellement (0 en haut)
\foreach \y/\ylab in {0/0, 1/50, 2/100, 3/150, 4/200, 5/250}
  \draw (0,\y) -- (-0.1,\y) node[left] {\ylab};
% Solidus Anhydre (croissant)
\draw[popblue, very thick, domain=0:5, samples=50, smooth] plot (\x, {0.01*\x*\x + 0.5*\x + 0.5}) node[right, popblue] {Solidus anhydre};
% Solidus Hydraté (plus bas)
\draw[poporange, very thick, domain=0:5, samples=50, smooth] plot (\x, {0.01*\x*\x + 0.5*\x + 1.5}) node[right, poporange] {Solidus hydraté};
% Géotherme Dorsale (Adiabatique - pente raide)
\draw[popgreen, very thick, dashed, domain=0:5, samples=50, smooth] plot (\x, {0.005*\x*\x + 0.2*\x + 0.2}) node[left, popgreen] {Géotherme dorsale};
% Géotherme Subduction (Coin mantellique - plus frais)
\draw[poppurple, very thick, dashed, domain=0:5, samples=50, smooth] plot (\x, {0.008*\x*\x + 0.3*\x + 0.8}) node[left, poppurple] {Géotherme coin mantellique};
% Zones de fusion
\fill[popgreen!30] (0,0.5) -- plot[domain=0:3.5, samples=20] (\x, {0.01*\x*\x + 0.5*\x + 0.5}) -- (3.5,0.5) -- cycle;
\node[popgreen!80!black, font=\bfseries\small, align=center] at (1.5,1.5){Fusion\\décompression};
\fill[poppurple!30] (0,1.5) -- plot[domain=0:4.5, samples=20] (\x, {0.01*\x*\x + 0.5*\x + 1.5}) -- (4.5,1.5) -- cycle;
\node[poppurple!80!black, font=\bfseries\small, align=center] at (2,2.5){Fusion\\hydratation};
% Point 100 km (y=2)
\draw[popdark, dotted] (0,2) -- (7,2) node[right] {100 km};
\draw[popdark, dotted] (3,0) -- (3,6);
\node[popdark] at (3,-0.3) {1100°C};
\node[popdark] at (4.5,-0.3) {1300°C};
\end{tikzpicture}
\legende{Diagramme Pression-Température (schématique) illustrant la fusion des péridotites. La fusion par décompression (dorsale) se produit quand l'adiabate croise le solidus anhydre. La fusion par hydratation (subduction) est rendue possible par l'abaissement du solidus (solidus hydraté) sous l'effet de l'eau libérée par la plaque plongeante.}
\end{popfigure}

\begin{exoresolu}[Fusion des péridotites en zone de subduction]
En l'absence d'eau, les péridotites commencent à fondre vers \SI{1300}{\celsius} à \SI{100}{km} de profondeur. En présence d'eau, leur température de début de fusion à cette profondeur tombe à environ \SI{1000}{\celsius}. La température du manteau au-dessus de la plaque plongeante y est de \SI{1100}{\celsius}.
\begin{enumerate}
\item Montrez que le manteau ne peut pas fondre à \SI{100}{km} en l'absence d'eau.
\item Expliquez comment la subduction permet la naissance d'un magma.
\item Citez le type de volcanisme associé et un exemple camerounais de volcan.
\end{enumerate}
\tcblower
\textbf{1. Fusion impossible sans eau.} À \SI{100}{km}, la température du manteau est de \SI{1100}{\celsius}, inférieure à la température de début de fusion des péridotites anhydres (\SI{1300}{\celsius}) : $1100 < 1300$. Le géotherme ne croise pas le solidus anhydre : les péridotites restent \textbf{solides}, aucune fusion n'est possible.

\textbf{2. Rôle de la subduction.} La plaque lithosphérique plongeante transporte des minéraux hydratés (serpentine, amphiboles, argiles). En profondeur, la déshydratation de ces minéraux libère de l'eau qui migre vers le manteau du coin mantellique. Or l'eau \textbf{abaisse le solidus} des péridotites : la température de fusion passe de \SI{1300}{\celsius} à \SI{1000}{\celsius}. Comme $1100 > 1000$, le géotherme dépasse désormais le solidus hydraté : il y a \textbf{fusion partielle} des péridotites. Le magma formé, moins dense, remonte et alimente le volcanisme d'arc.

\textbf{3. Volcanisme associé.} Le magma issu de la subduction est riche en silice, visqueux, chargé en gaz : il produit un volcanisme de type \textbf{explosif} (andésites, dacites, rhyolites, nuées ardentes), caractéristique des zones de subduction (Andes, Japon, Arcipelins). Au Cameroun, le \textbf{Mont Cameroun} (\SI{4095}{m}), volcan actif de la ligne volcanique du Cameroun, illustre le volcanisme local, bien qu'il soit lié à un \textbf{point chaud} (panache mantellique) plutôt qu'à une subduction active actuelle.
\textbf{Conclusion} : c'est l'hydratation du manteau par la plaque plongeante, qui abaisse la température de fusion des péridotites, qui permet la naissance des magmas des zones de subduction.
\end{exoresolu}

\courssub{Décrire les mouvements de convergence et de coulissage des plaques}

\begin{methode}[Identifier un type de frontière de plaques]
\begin{enumerate}
\item \textbf{Repérer les indices du document} : fosse océanique profonde, arc volcanique, chaîne de montagnes, séismes profonds, décrochement avec séismes superficiels.
\item \textbf{Convergence} : deux plaques se rapprochent. Deux cas :
\begin{itemize}
\item \textbf{Subduction} : une plaque océanique (dense) plonge sous une autre plaque (océanique ou continentale) $\Rightarrow$ fosse océanique + volcanisme explosif d'arc (andésites) + séismes de profondeur croissante définissant le \cle{plan de Bénioff} (jusqu'à \SI{670}{km}).
\item \textbf{Collision} : deux continents (croûtes légères) se rencontrent $\Rightarrow$ pas de subduction possible, épaississement crustal par plissement et chevauchement $\Rightarrow$ chaîne de montagnes intracontinentales (ex. Himalaya), séismes de profondeur moyenne, \textbf{pas de volcanisme} (sauf rares cas post-collisionnels).
\end{itemize}
\item \textbf{Coulissage (Frontière transformante)} : deux plaques glissent horizontalement l'une contre l'autre le long d'une \cle{faille transformante} (grande faille décrochante) $\Rightarrow$ séismes superficiels fréquents (foyers $<$ \SI{30}{km}), \textbf{pas de volcanisme}, pas de création/destruction de lithosphère. Compense les décalages entre segments de dorsale.
\item \textbf{Facteurs clés} : la convergence est motrice (poussée des dorsales \emph{ridge push}, traction des plaques plongeantes froides et denses \emph{slab pull}) ; le coulissage accommode la géométrie sphérique de la Terre.
\item \textbf{Conclure} en nommant la frontière : divergente (dorsale), convergente (subduction/collision) ou transformante (coulissage).
\end{enumerate}
\end{methode}

\begin{popfigure}
\begin{tikzpicture}[scale=0.75, every node/.style={font=\small}]
% Subduction
\begin{scope}[xshift=-7cm]
\draw[popdark, thick] (-2,0) -- (2,0) node[midway, above] {Continent};
\draw[popblue, thick] (-2,-0.5) .. controls (-1,-1) and (0,-2) .. (1,-3) .. controls (1.5,-4) and (2,-5) .. (2,-6) node[right] {Plaque océanique plongeante};
\draw[poporange, thick, decorate, decoration={snake, amplitude=1pt, segment length=4pt}] (-1.5,-0.5) -- (-0.5,-0.5) node[above] {Fosse};
\fill[poppink!50] (-1.5,-1) -- (-0.5,-1) -- (0,-2) -- (-1,-2) -- cycle node[midway, white, font=\bfseries\tiny] {Prisme d'accrétion};
\fill[poppurple!50] (0.5,-1) -- (1.5,-1) -- (1,-2) -- (0,-2) -- cycle node[midway, white, font=\bfseries\tiny] {Arc volcanique};
% Séismes Benioff
\foreach \y in {1,2,3,4,5} {
  \foreach \x in {-0.5,0,0.5} {
    \fill[popred] (\x,-\y) circle (1.5pt);
  }
}
\node[popred, font=\bfseries\small] at (2.5,-3) {Plan de Bénioff};
\draw[->, popblue, thick] (1.5,-5.5) -- (1.5,-6.5) node[below] {Slab pull};
\node[font=\bfseries, popdark] at (0,0.5) {SUBDUCTION};
\end{scope}

% Collision
\begin{scope}[xshift=0cm]
\draw[poporange, thick] (-2,0) -- (-0.5,0) .. controls (-0.2,0.5) and (0.2,0.5) .. (0.5,0) -- (2,0) node[midway, above] {Continent};
\draw[poporange, thick] (-2,-0.5) .. controls (-1,-1) and (-0.5,-2) .. (0,-2.5) .. controls (0.5,-2) and (1,-1) .. (2,-0.5);
\draw[poporange, thick] (-2,-1) .. controls (-1,-1.5) and (-0.5,-2.5) .. (0,-3) .. controls (0.5,-2.5) and (1,-1.5) .. (2,-1);
% Chevauchements
\draw[->, popdark, thick] (-1.5,-0.2) -- (-0.5,-0.2);
\draw[->, popdark, thick] (1.5,-0.2) -- (0.5,-0.2);
% Séismes
\foreach \y in {0.5,1,1.5,2,2.5} {
  \fill[popred] (-0.5,-\y) circle (1.5pt);
  \fill[popred] (0.5,-\y) circle (1.5pt);
}
\node[popdark, font=\bfseries\small] at (0,-3.5) {Épaississement crustal (racines)};
\node[font=\bfseries, popdark] at (0,0.5) {COLLISION};
\end{scope}

% Coulissage (Transformante)
\begin{scope}[xshift=7cm]
\draw[popgreen, very thick, ->] (-2,0.2) -- (2,0.2) node[midway, above] {Plaque A};
\draw[popgreen, very thick, <-] (-2,-0.2) -- (2,-0.2) node[midway, below] {Plaque B};
\draw[popdark, thick, decorate, decoration={zigzag, amplitude=2pt, segment length=4pt}] (-2,0) -- (2,0) node[midway, above=4pt] {Faille transformante};
% Séismes superficiels
\foreach \x in {-1.5,-0.5,0.5,1.5} {
  \fill[popred] (\x,0) circle (2pt);
}
\node[popred, font=\bfseries\small] at (0,-1) {Séismes superficiels (< 30 km)};
\node[font=\bfseries, popdark] at (0,0.7) {COULISSAGE};
\end{scope}
\end{tikzpicture}
\legende{Schémas des trois types de frontières de plaques convergentes/transformantes. Haut : Subduction (fosse, prisme, arc volcanique, plan de Bénioff). Milieu : Collision continentale (chevauchements, épaississement, séismes moyens, pas de volcanisme). Bas : Faille transformante (coulissage horizontal, séismes superficiels, pas de volcanisme).}
\end{popfigure}

\begin{exoresolu}[Identifier les frontières de plaques]
Le tableau suivant résume les observations faites en trois régions du globe :
\begin{center}
\begin{tabularx}{\linewidth}{|l|X|X|X|}\hline
\rowcolor{popblueL} \textbf{Indices} & \textbf{Région 1} & \textbf{Région 2} & \textbf{Région 3} \\ \hline
Relief & Fosse de \num{-8000} m & Chaîne de hautes montagnes & Faille décrochante \\ \hline
Volcanisme & Explosif (andésites) & Absent & Absent \\ \hline
Séismes & Jusqu'à \SI{600}{km} de profondeur & Profondeur moyenne & Superficiels ($<$ \SI{30}{km}) \\ \hline
\end{tabularx}
\end{center}
Identifiez le type de frontière de plaques dans chaque région et justifiez.
\tcblower
\textbf{Région 1 : convergence avec subduction.} La fosse profonde (\num{-8000} m) marque l'enfoncement d'une plaque océanique sous une autre. Le volcanisme andésitique explosif résulte de la fusion partielle du manteau hydraté au-dessus de la plaque plongeante. Les séismes profonds (jusqu'à \SI{600}{km}) dessinent le plan de Bénioff, c'est-à-dire la plaque plongeante elle-même. Exemple : fosse du Pérou--Chili (subduction Nazca/Sud-Amérique).

\textbf{Région 2 : convergence avec collision.} L'absence de fosse et de volcanisme, associée à une chaîne de hautes montagnes et à des séismes de profondeur moyenne, traduit la rencontre de deux continents : aucun ne peut plonger (densités trop faibles), la lithosphère s'épaissit par plissement et chevauchement. Exemple : Himalaya (collision Inde--Eurasie).

\textbf{Région 3 : coulissage (faille transformante).} Les deux plaques glissent horizontalement l'une contre l'autre : ni création ni destruction de lithosphère, donc pas de volcanisme ; les contraintes s'accumulent et se relâchent brutalement en séismes superficiels (foyers à moins de \SI{30}{km}). Exemple : faille de San Andreas en Californie (frontière Pacifique/Amérique du Nord).

\textbf{Conclusion} : fosse + volcanisme + séismes profonds = subduction ; montagnes sans volcanisme = collision ; faille décrochante + séismes superficiels = coulissage.
\end{exoresolu}

\courssub{Les catastrophes naturelles liées aux mouvements des plaques}

\begin{methode}[Expliquer l'origine d'une catastrophe naturelle]
\begin{enumerate}
\item \textbf{Identifier la catastrophe} : séisme, tsunami, éruption volcanique, glissement de terrain, inondation, éruption limnique.
\item \textbf{Relier au mouvement des plaques} :
\begin{itemize}
\item \textbf{Séisme} : rupture brutale des roches accumulant les contraintes le long des failles (frontières convergentes ou transformantes) ; foyer (hypocentre) en profondeur, épicentre en surface. Magnitude (énergie, échelle de Richter/moment) vs Intensité (dégâts, échelle MSK/EMS).
\item \textbf{Tsunami} : séisme sous-marin de subduction (magnitude $>$ 7,5) déplaçant verticalement le fond océanique (soulèvement/affaissement) $\Rightarrow$ onde rapide en pleine mer (longueur d'onde \SI{100}{km}, vitesse jusqu'à \SI{800}{km/h}, amplitude faible), amplification et déferlante près des côtes (ralentissement, raccourcissement longueur d'onde, hausse amplitude).
\item \textbf{Volcanisme} : magma des zones d'accrétion (effusif, basaltes fluides, coulées de lave, fontaines) ou de subduction (explosif, andésites visqueuses, nuées ardentes, lahars, retombées cendres).
\item \textbf{Glissements de terrain et inondations} : favorisés par reliefs jeunes (tectonique active), pentes fortes, pluies diluviennes, déforestation, urbanisation anarchique ; parfois déclenchés par un séisme (secousses) ou une éruption (fonte neige/glace $\rightarrow$ lahars).
\item \textbf{Éruption limnique} : dégazage brutal d'un lac volcanique stratifié (accumulation de \ce{CO2} ou \ce{CH4} dissous en profondeur), déclenchée par un glissement de terrain, un séisme ou un renversement thermique. Exemple majeur : lac Nyos (Cameroun, 1986) et lac Monoun (1984).
\end{itemize}
\item \textbf{Évaluer le risque} : \cle{Risque} = \cle{Aléa} (probabilité/intensité du phénomène) $\times$ \cle{Vulnérabilité} (exposition et fragilité des populations/biens). La prévention vise à réduire la vulnérabilité.
\end{enumerate}
\end{methode}

\begin{popfigure}
\begin{tikzpicture}[scale=0.7, every node/.style={font=\small}]
% Lac Nyos coupe
% Relief
\draw[popdark, thick] (-5,0) .. controls (-3,1) and (3,1) .. (5,0) .. controls (4,-1) and (-4,-1) .. (-5,0);
% Eau lac
\fill[popblue!30] (-4,0) .. controls (-3,0.5) and (3,0.5) .. (4,0) -- (4,-3) .. controls (3,-3.5) and (-3,-3.5) .. (-4,-3) -- cycle;
% Fond lac
\draw[poporange, thick] (-4,-3) .. controls (-3,-3.5) and (3,-3.5) .. (4,-3);
% Magma profondeur
\fill[poppink!50] (-5,-5) rectangle (5,-7) node[midway, white, font=\bfseries] {Magma en refroidissement (source CO2)};
% Flèches CO2
\draw[->, popgray, thick] (0,-5) -- (0,-3.2) node[midway, right] {Remontée CO2};
\draw[->, popgray, thick] (0,-3) -- (0,-2) node[midway, right] {Dissolution};
% Eau profonde saturée
\fill[popblue!60] (-3.5,-3) rectangle (3.5,-2.5) node[midway, white, font=\bfseries] {Eau profonde saturée en CO2};
% Tuyau dégazage
\draw[popgreen, very thick] (0,-0.2) -- (0,-3) node[midway, right, popgreen] {Tuyau de dégazage};
\draw[popgreen, thick, ->] (0,-0.2) -- (0,0.5) node[above, popgreen] {Fontaine dégazage};
% Nuage CO2 (catastrophe)
\draw[popred, thick, decorate, decoration={snake, amplitude=2pt, segment length=5pt}] (-1,-0.2) -- (-3,-0.2) -- (-4,-1.5) -- (-5,-2.5);
\draw[popred, thick, decorate, decoration={snake, amplitude=2pt, segment length=5pt}] (1,-0.2) -- (3,-0.2) -- (4,-1.5) -- (5,-2.5);
\node[popred, font=\bfseries] at (0,-1) {Nuage CO2 dense (vallées)};
% Légendes
\node[popblue] at (-4.5,0.3) {Surface lac};
\node[poporange] at (-4.5,-3.3) {Fond lac};
\node[popdark] at (-5.5,-1.5) {Cratère volcanique};
\end{tikzpicture}
\legende{Mécanisme de l'éruption limnique du lac Nyos (1986) et principe du dégazage artificiel par tuyau (siphon auto-entretenu). Le CO2 magmatique se dissout dans les eaux profondes ; un déclencheur provoque le dégazage brutal ; le CO2, plus dense que l'air, s'écoule dans les vallées. Le tuyau permet un dégazage continu et sécurisé.}
\end{popfigure}

\begin{exoresolu}[La catastrophe du lac Nyos (Cameroun, 1986)]
Le 21 août 1986, le lac Nyos, lac de cratère volcanique du Nord-Ouest du Cameroun (ligne volcanique du Cameroun), a brutalement libéré un nuage de dioxyde de carbone qui a asphyxié environ \num{1746} personnes et des milliers de têtes de bétail dans les vallées voisines.
\begin{enumerate}
\item Expliquez l'origine volcanique du \ce{CO2} accumulé au fond du lac.
\item Pourquoi le gaz a-t-il tué principalement dans les vallées ?
\item Proposez une mesure technique de prévention mise en œuvre depuis.
\end{enumerate}
\tcblower
\textbf{1. Origine du \ce{CO2}.} Le lac Nyos occupe le cratère d'un volcan de la ligne volcanique du Cameroun. Sous le lac, le magma en refroidissement dégaze du dioxyde de carbone qui remonte à travers les roches fracturées et se dissout dans les eaux profondes du lac. Sous la pression de la colonne d'eau (stabilité thermique empêchant le mélange), le \ce{CO2} reste dissous et s'accumule pendant des années près du fond, jusqu'à saturation (conditions de stabilité : fond froid, haute pression).

\textbf{2. Mécanisme mortel.} Un déclencheur (glissement de terrain sur la paroi du cratère, perturbation thermique, ou séisme mineur) a provoqué la remontée brutale des eaux profondes saturées : la décompression a libéré le gaz dissous (ébullition explosive, \textbf{éruption limnique}), formant un nuage de \ce{CO2} d'environ \SI{1.6}{km^3}. Le \ce{CO2} est plus dense que l'air ($M(\ce{CO2}) = \SI{44}{g/mol} > \SI{29}{g/mol}$ pour l'air) : le nuage est resté au sol et s'est écoulé par gravité le long des vallées (vitesse \SI{20}{à 50}{km/h}), déplaçant l'air respirable et asphyxiant les habitants et le bétail des zones basses (jusqu'à \SI{25}{km} du lac).

\textbf{3. Prévention technique.} La mesure technique adoptée (dès 2001, équipe française/camerounaise) est le \textbf{dégazage contrôlé par siphon auto-entretenu} : un tuyau vertical (diam. \SI{15}{cm} puis \SI{30}{cm}) plongé dans les eaux profondes (profondeur \SI{200}{m}) établit un siphon qui remonte en continu l'eau chargée en \ce{CO2}. Arrivée en surface, la décompression fait dégazer l'eau en jet (fontaine) inoffensif. Plusieurs tuyaux ont été installés pour augmenter le débit. On y ajoute une surveillance continue (capteurs CO2, température, sismomètres) et un système d'alerte sonore pour les populations riveraines.
\textbf{Conclusion} : la catastrophe du lac Nyos illustre un risque volcanique indirect (gaz dissous dans un lac de cratère) dont la prévention repose sur le dégazage artificiel permanent, la surveillance instrumentée et l'éducation des populations.
\end{exoresolu}

\courssub{Élaborer des outils de sensibilisation à la gestion des catastrophes}

\begin{methode}[Construire un outil de sensibilisation : prévision, prédiction, prévention]
\begin{enumerate}
\item \textbf{Distinguer rigoureusement les trois notions} (clé du Bac) :
\begin{itemize}
\item \textbf{Prévision} : estimation statistique à long terme de la probabilité d'un événement dans une région (cartographie des zones à risque, zonage sismique, retour d'expérience historique). Ex : « Cette zone a 10\% de chance de subir un séisme magnitude $>$ 7 dans les 50 ans ».
\item \textbf{Prédiction} : annonce précise (date, lieu, magnitude) d'un événement \emph{imminent} (heures/jours avant). Elle est \textbf{impossible pour les séismes} (aucun précurseur fiable universel), \textbf{partiellement possible pour les éruptions} (surveillance des signes précurseurs : gonflement édifice, sismicité volcano-tectonique, dégazage \ce{SO2}/\ce{CO2}, variations chimiques sources).
\item \textbf{Prévention} : ensemble des mesures structurelles et non structurelles réduisant la vulnérabilité : constructions parasismiques (normes Eurocode 8 / CM97 au Cameroun), plans d'évacuation, digues de protection, reboisement des pentes (lutte érosion), éducation/formation des populations (exercices réguliers), assurance catastrophe.
\end{itemize}
\item \textbf{Choisir la cible et le support} : affiche, dépliant, sketch, spot radio en langues locales (fulfulde, ewondo, bassa, pidgin...), vidéo courte, adaptés au public (élèves, riverains, autorités).
\item \textbf{Structurer le message} : le danger (aléa local identifié), les signes d'alerte (ex : recul de la mer = tsunami), les conduites à tenir (avant : kit survie ; pendant : se protéger ; après : ne pas rentrer, écouter radio), les numéros d'urgence (112, 118, numéros locaux protection civile).
\item \textbf{Vérifier} l'exactitude scientifique (pas de « prédiction séisme »), la clarté (une idée par phrase, pictogrammes universels), la lisibilité (police grande, contraste couleurs).
\end{enumerate}
\end{methode}

\begin{exoresolu}[Concevoir une affiche de sensibilisation aux séismes et tsunamis]
Une ville côtière est située près d'une fosse de subduction active. Rédigez le contenu d'une affiche de sensibilisation destinée aux élèves, comportant : (a) le danger, (b) deux conduites à tenir pendant un séisme, (c) deux conduites à tenir en cas d'alerte au tsunami, (d) une mesure de prévention à long terme.
\tcblower
\textbf{(a) Le danger.} « Ta ville est proche d'une zone de subduction où une plaque océanique plonge sous le continent. Les contraintes accumulées peuvent se relâcher brutalement : c'est un \textbf{séisme}. Si le fond marin se soulève ou s'affaisse soudainement, une vague géante, le \textbf{tsunami}, peut atteindre la côte en quelques minutes. \textbf{La prédiction de l'heure exacte est impossible.} »

\textbf{(b) Pendant le séisme (se protéger).}
\begin{itemize}
\item \textbf{S'abriter} immédiatement sous une table solide, protéger sa tête et son cou avec les bras, s'éloigner des fenêtres, miroirs, objets suspendus, bibliothèques.
\item \textbf{Ne pas utiliser} les escaliers ni les ascenseurs ; \textbf{ne pas sortir en courant} pendant les secousses (risque de chutes, chutes de débris en façade). Si on est dehors : s'éloigner des bâtiments, poteaux électriques, falaises.
\end{itemize}

\textbf{(c) En cas d'alerte au tsunami (sirène, SMS, recul anormal de la mer, bruit de train).}
\begin{itemize}
\item \textbf{Quitter immédiatement} le bord de mer et gagner un point haut (plus de \SI{20}{m} d'altitude si possible) ou l'intérieur des terres (\SI{> 1}{km}), \textbf{à pied} (voitures = embouteillages mortels).
\item \textbf{Ne pas retourner} sur la côte avant l'annonce officielle de fin d'alerte par les autorités : un tsunami comporte plusieurs vagues successives (la première n'est pas la plus grande), l'intervalle peut être de \SI{10}{à 60}{min}.
\end{itemize}

\textbf{(d) Prévention à long terme (réduire la vulnérabilité).} Construire et rénover selon les normes parasismiques en vigueur (fondations profondes, chaînages, matériaux ductiles), aménager et signaler des itinéraires et zones de refuge en hauteur accessibles à tous (personnes à mobilité réduite comprises), et organiser \textbf{régulièrement des exercices d'évacuation} (simulation séisme + tsunami) dans toutes les écoles et lieux publics.

\textbf{Conclusion} : un bon outil de sensibilisation nomme clairement l'aléa local, donne des consignes simples, mémorisables et adaptées au contexte, et rappelle que la prévention (bâtir solide, s'entraîner, planifier) sauve plus de vies que la prédiction, impossible pour les séismes.
\end{exoresolu}

\begin{attention}[Les 5 erreurs qui coûtent des points au Bac]
\begin{enumerate}
\item \textbf{Confondre prévision, prédiction et prévention.} La \textbf{prédiction des séismes} (date, lieu, magnitude précis) est \emph{scientifiquement impossible} à l'heure actuelle ; seule la \textbf{prévision} statistique (aléa à long terme) et la \textbf{prévention} (réduction vulnérabilité) sont réalisables. Écrire « les scientifiques prédisent les séismes » est une \textbf{faute scientifique éliminatoire}.
\item \textbf{Écrire que « le manteau est fondu » ou « liquide ».} Le manteau supérieur est \textbf{solide} (comportement visqueux/ductile à l'échelle géologique) ; seule une \emph{fusion partielle} (quelques pourcents, 5--20\%) des péridotites produit le magma, par \textbf{décompression} (dorsale) ou par \textbf{hydratation} (subduction).
\item \textbf{Inverser les mécanismes de fusion.} \textbf{Décompression} = zones d'accrétion (dorsales, rifts) $\rightarrow$ basaltes (fluides) ; \textbf{Abaissement du solidus par l'eau} = zones de subduction $\rightarrow$ andésites/dacites (visqueux, explosif). Ne pas les intervertir.
\item \textbf{Oublier le facteur 2 dans les vitesses d'expansion.} La vitesse calculée à partir de l'âge et de la distance à l'axe ($v = d/t$) est une \emph{demi-vitesse} (vitesse d'une plaque par rapport à l'axe) ; l'\textbf{écartement total} des deux plaques (vitesse relative) vaut le \textbf{double}. Toujours préciser l'unité (\si{cm/an}) et convertir correctement km en cm ($1\ \text{km} = 10^5\ \text{cm}$) et Ma en ans ($1\ \text{Ma} = 10^6\ \text{ans}$).
\item \textbf{Identifier une frontière sans citer d'indices chiffrés.} Toute conclusion (subduction, collision, coulissage) doit être \textbf{justifiée par des données du document} : profondeur de la fosse, profondeur des foyers sismiques (plan de Bénioff), type de roches volcaniques (basalte vs andésite), présence/absence de volcanisme. Une affirmation sans preuve documentaire n'est pas valorisée.
\end{enumerate}
\end{attention}

\begin{savaistu}[Le Cameroun face aux risques naturels]
Le Cameroun, bien que situé sur une plaque stable (plaque africaine), n'est pas exempt de risques :
\begin{itemize}
\item \textbf{Volcanisme} : La \textbf{ligne volcanique du Cameroun} (Mont Cameroun, Manengouba, Bamboutos, lacs Nyos, Monoun, Barombi Mbo) est une chaîne de volcans actifs (point chaud). Le Mont Cameroun a érupté en 1999, 2000, 2012 (coulées de lave effusives). Les lacs de cratère (Nyos, Monoun) présentent un risque limnique (\ce{CO2}).
\item \textbf{Séismes} : Faible à modéré. Quelques séismes ressentis dans l'Adamaoua et le Nord (failles intraplaques réactivées), ex. séisme de 1984 à Ngaoundéré (M $\approx$ 5).
\item \textbf{Glissements de terrain} : Fréquents dans les zones de relief (Ouest, Nord-Ouest, Sud-Ouest) lors des fortes pluies, aggravés par la déforestation et l'habitat précaire sur pentes.
\item \textbf{Inondations} : Récurrentes dans le Logone-et-Chari (Nord, plaine d'inondation du Logone), à Douala (urbanisation zones marécageuses, drainage insuffisant) et Yaoundé.
\item \textbf{Gestion} : Le \textbf{MINATD} (Ministère de l'Administration Territoriale et de la Décentralisation) et la \textbf{Protection Civile} coordonnent la prévention. Le \textbf{Centre National de Surveillance Volcanologique et Sismologique (CNSVS)} de l'Université de Buéa surveille le Mont Cameroun et les lacs. Le dégazage du lac Nyos est une réussite internationale de prévention technique.
\end{itemize}
\end{savaistu}

\newpage

\coursec{Exercices}

\courssub{Je vérifie mes connaissances}

\begin{exercice}[QCM : les fonds océaniques]
Pour chaque affirmation, choisis la (ou les) bonne(s) réponse(s) et justifie brièvement ton choix.
\begin{enumerate}
\item La dorsale médio-océanique est :
\begin{enumerate}
\item une chaîne de montagnes immergée au centre des océans ;
\item une fosse profonde située au contact d'un continent ;
\item une zone d'accrétion de nouvelle croûte océanique ;
\item une zone de subduction de la lithosphère.
\end{enumerate}
\item Le rift axial d'une dorsale est :
\begin{enumerate}
\item une vallée étroite creusée au sommet de la dorsale ;
\item le lieu de sortie des magmas basaltiques ;
\item une structure formée par effondrement entre deux failles normales ;
\item une zone de compression de la lithosphère.
\end{enumerate}
\item La fosse océanique marque :
\begin{enumerate}
\item une limite divergente de plaques ;
\item une limite convergente avec subduction ;
\item une limite transformante (coulissage) ;
\item l'axe d'expansion océanique.
\end{enumerate}
\item La faille transformante est une limite de plaque où :
\begin{enumerate}
\item deux plaques s'éloignent l'une de l'autre ;
\item une plaque plonge sous une autre ;
\item deux plaques coulissent horizontalement l'une contre l'autre ;
\item le magma remonte en permanence.
\end{enumerate}
\end{enumerate}
\end{exercice}

\begin{exercice}[Vrai ou faux, avec justification]
Indique si chaque affirmation est vraie ou fausse, puis justifie ta réponse en une ou deux phrases.
\begin{enumerate}
\item Le plancher océanique est d'autant plus ancien qu'on s'éloigne de l'axe de la dorsale.
\item La fusion des péridotites sous les dorsales résulte d'une augmentation de la température du manteau.
\item En zone de subduction, la fusion des péridotites du manteau est favorisée par l'apport d'eau libérée par la plaque plongeante.
\item Les séismes et les éruptions volcaniques se répartissent au hasard à la surface du globe.
\item La prévention d'une catastrophe naturelle consiste à annoncer la date et le lieu exacts de son déclenchement.
\end{enumerate}
\end{exercice}

\begin{exercice}[Définitions]
Définis en deux ou trois phrases chacun des termes suivants, en employant le vocabulaire scientifique exact : \cle{dorsale océanique} ; \cle{rift} ; \cle{subduction} ; \cle{tsunami} ; \cle{aléa} ; \cle{vulnérabilité} ; \cle{prévention}.
\end{exercice}

\begin{exercice}[Calcul direct : taux d'expansion]
Sur une dorsale océanique, un basalte situé à \SI{100}{km} de l'axe du rift a un âge de 5 millions d'années (Ma).
\begin{enumerate}
\item Calcule la vitesse d'éloignement de ce basalte par rapport à l'axe (demi-taux d'expansion), en centimètres par an.
\item En déduis le taux d'expansion total de l'océan à cette dorsale.
\item À quelle distance de l'axe se trouverait un basalte âgé de 12 Ma, en supposant le taux constant ?
\end{enumerate}
\end{exercice}

\courssub{Je m'entraîne}

\begin{exercice}[Exploitation d'un profil bathymétrique]
Un profil bathymétrique transocéanique, tracé perpendiculairement à une dorsale atlantique, fait apparaître de l'ouest vers l'est : une plaine abyssale vers \SI{-5000}{m}, une élévation large atteignant \SI{-2500}{m} avec, à son sommet, une dépression axiale de \SI{1000}{m} de profondeur relative, puis une nouvelle plaine abyssale symétrique.
\begin{enumerate}
\item Nomme chacune des trois structures repérées sur ce profil.
\item Représente ce profil par un schéma légendé en indiquant l'axe de la dorsale.
\item Justifie, en une phrase, la symétrie du profil de part et d'autre de l'axe.
\end{enumerate}
\end{exercice}

\begin{exercice}[Taux d'expansion : tableau et courbe]
Le tableau ci-dessous donne l'âge des basaltes du plancher océanique en fonction de leur distance à l'axe d'une dorsale du Pacifique.
\begin{center}
\begin{tabular}{|l|c|c|c|c|c|}
\hline
\rowcolor{popblueL} Distance à l'axe (km) & 0 & 200 & 400 & 600 & 800 \\
\hline
Âge des basaltes (Ma) & 0 & 2 & 4 & 6 & 8 \\
\hline
\end{tabular}
\end{center}
\begin{enumerate}
\item Trace la courbe donnant l'âge des basaltes en fonction de la distance à l'axe (échelles au choix, axes légendés avec leurs unités).
\item Détermine graphiquement puis par le calcul le demi-taux d'expansion, en cm/an.
\item Le plancher de cet océan mesure \SI{12000}{km} de large. Estime l'âge des basaltes les plus anciens de cet océan.
\item Explique pourquoi aucune croûte océanique très ancienne (plus de 200 Ma) n'est connue sur Terre.
\end{enumerate}
\end{exercice}

\begin{exercice}[Fusion des péridotites : lecture d'un diagramme]
Sur un diagramme pression-température sont tracés : le solidus des péridotites anhydres, le solidus des péridotites hydratées (décalé d'environ \SI{400}{\celsius} vers les basses températures) et deux géothermes : un géotherme de dorsale (remontée adiabatique) et un géotherme de subduction.
\begin{enumerate}
\item Rappelle ce qu'est le solidus d'une roche.
\item Explique, à partir du diagramme, le mécanisme de fusion des péridotites sous une dorsale (zone d'accrétion).
\item Explique le mécanisme de fusion des péridotites en zone de subduction, en précisant l'origine de l'eau.
\item Compare la nature des magmas produits dans les deux contextes (basaltes fluides / magmas visqueux) et le type de volcanisme associé.
\end{enumerate}
\end{exercice}

\begin{exercice}[Classer : prévision, prédiction ou prévention]
Classe chacune des actions suivantes dans la catégorie qui lui correspond (\cle{prévision}, \cle{prédiction} ou \cle{prévention}), en justifiant ton choix :
\begin{enumerate}
\item Construction de digues et de canaux d'évacuation des eaux dans les quartiers bas de Douala.
\item Surveillance sismométrique continue du Mont Cameroun par l'observatoire volcanologique.
\item Établissement de cartes des zones à risque d'inondation interdisant toute nouvelle construction.
\item Annonce, quelques heures à l'avance, de l'arrivée probable d'un tsunami sur une côte.
\item Organisation d'exercices d'évacuation dans les écoles d'une zone sismique.
\item Mesure de la déformation du sol et des émissions de gaz sur un volcan actif pour anticiper une éruption.
\end{enumerate}
\end{exercice}

\courssub{Je me prépare au Bac}

\begin{exercice}[Type Bac : le plancher océanique atlantique]
\textbf{Document 1.} Un profil bathymétrique est réalisé entre les côtes du Brésil et celles de l'Afrique de l'Ouest. Il montre, des continents vers le centre de l'océan : un plateau continental, un talus continental, une plaine abyssale vers \SI{-5000}{m}, puis une élévation culminant vers \SI{-2500}{m}, creusée d'une vallée axiale. Le profil est globalement symétrique de part et d'autre de cette élévation.

\textbf{Document 2.} Des datations des basaltes prélevés de part et d'autre de l'axe donnent les résultats suivants :
\begin{center}
\begin{tabular}{|l|c|c|c|c|}
\hline
\rowcolor{popblueL} Distance à l'axe (km) & 0 & 300 & 600 & 900 \\
\hline
Âge (Ma) & 0 & 10 & 20 & 30 \\
\hline
\end{tabular}
\end{center}
Les âges mesurés à égale distance de l'axe, de chaque côté, sont identiques.

\textbf{Document 3.} Les anomalies magnétiques enregistrées dans les basaltes forment des bandes parallèles à l'axe, symétriques de part et d'autre de celui-ci.

\begin{enumerate}
\item À partir du document 1, identifie et nomme les structures des fonds océaniques traversées par ce profil, et représente-les par un schéma légendé.
\item À partir du document 2, calcule le demi-taux d'expansion de l'Atlantique à cette latitude, en cm/an, puis le taux d'expansion total.
\item Montre, en exploitant les documents 2 et 3, que la dorsale est le lieu de formation continue d'une nouvelle croûte océanique qui s'écarte ensuite de part et d'autre de l'axe.
\item Décris en quelques lignes le fonctionnement du rift axial qui permet cette expansion (nature des failles, remontée du magma, formation du plancher).
\end{enumerate}
\end{exercice}

\begin{exercice}[Type Bac : séismes et volcans, indices des limites de plaques]
\textbf{Document 1.} La carte d'une région du globe montre la répartition des foyers des séismes enregistrés pendant 30 ans : les foyers superficiels (moins de \SI{70}{km}) se situent au niveau d'une fosse océanique profonde de \SI{8000}{m} ; les foyers intermédiaires et profonds (jusqu'à \SI{600}{km}) s'enfoncent en biais sous une chaîne de volcans continentale, selon un plan incliné d'environ 30 degrés.

\textbf{Document 2.} Les laves des volcans de cette chaîne sont visqueuses, riches en silice, et les éruptions y sont explosives.

\textbf{Document 3.} Rappel expérimental : le solidus des péridotites hydratées se situe environ \SI{400}{\celsius} en dessous de celui des péridotites anhydres ; à \SI{100}{km} de profondeur, les hydrates minéraux des roches basaltiques se déstabilisent et libèrent de l'eau.

\begin{enumerate}
\item À partir du document 1, détermine le type de limite de plaque qui existe dans cette région. Justifie ta réponse en exploitant la profondeur et la disposition des foyers sismiques.
\item Nomme le plan incliné matérialisé par les foyers profonds et précise ce qu'il représente.
\item En exploitant les documents 2 et 3, explique l'origine du magma des volcans de cette région et justifie le caractère explosif du volcanisme.
\item Résume en quelques phrases pourquoi cette région cumule plusieurs aléas naturels (séismes, volcanisme, tsunamis éventuels) et cite deux mesures de prévention adaptées.
\end{enumerate}
\end{exercice}

\begin{exercice}[Type Bac : la catastrophe du lac Nyos et la gestion des risques au Cameroun]
\textbf{Document 1.} Le 21 août 1986, le lac Nyos, lac de cratère situé dans la région du Nord-Ouest du Cameroun, a brutalement libéré un nuage de dioxyde de carbone (\ce{CO2}) estimé à plus de \SI{1}{km^3}. Ce gaz, plus dense que l'air, s'est répandu dans les vallées environnantes, provoquant la mort par asphyxie d'environ \num{1700} personnes et de milliers de têtes de bétail, jusqu'à \SI{25}{km} du lac.

\textbf{Document 2.} Le \ce{CO2} provient d'une source magmatique profonde ; il s'accumule lentement dissous dans les eaux profondes du lac, stratifiées et stables. Une perturbation (glissement de terrain, refroidissement des eaux de surface) peut provoquer la remontée brutale de ces eaux chargées en gaz : c'est l'éruption limnique.

\textbf{Document 3.} Depuis 2001, des tuyaux de dégazage pompent en continu l'eau profonde du lac Nyos, qui jaillit en fontaine et libère le \ce{CO2} progressivement. Des capteurs mesurent en permanence la concentration en gaz dissous, et les populations riveraines ont été sensibilisées et partiellement relogées en hauteur.

\begin{enumerate}
\item À partir des documents 1 et 2, identifie l'aléa naturel responsable de la catastrophe du lac Nyos et décris son mécanisme.
\item Explique pourquoi la vulnérabilité des populations riveraines était élevée (localisation des villages, méconnaissance du phénomène).
\item En exploitant le document 3, classe les mesures mises en œuvre selon qu'elles relèvent de la prévision, de la prédiction ou de la prévention, en justifiant chaque classement.
\item Rédige un court texte de sensibilisation (10 à 15 lignes) destiné aux habitants d'un quartier de Douala exposé aux inondations, en distinguant clairement ce qu'ils peuvent faire avant, pendant et après une crue.
\end{enumerate}
\end{exercice}

\newpage

\coursec{Corrigés des exercices}

\begin{corrige}[Exercice 1 — QCM : les fonds océaniques]
\begin{enumerate}
\item \textbf{Bonnes réponses : a et c.} La dorsale médio-océanique est une chaîne de montagnes immergée (elle culmine vers \SI{-2500}{m} alors que les plaines abyssales voisines sont vers \SI{-5000}{m}) et c'est une \cle{zone d'accrétion} : la nouvelle croûte océanique s'y forme en permanence. Les réponses b et d décrivent une fosse de subduction, qui est une limite convergente et non divergente.
\item \textbf{Bonnes réponses : a, b et c.} Le rift axial est une vallée étroite (quelques kilomètres de large, environ \SI{1000}{m} de profondeur relative) creusée au sommet de la dorsale ; il résulte de l'\cle{effondrement} d'un compartiment entre deux \cle{failles normales} (failles d'extension) ; c'est par ses fissures que sortent les magmas basaltiques. La réponse d est fausse : le rift est une zone d'\emph{extension}, non de compression.
\item \textbf{Bonne réponse : b.} La fosse océanique (jusqu'à \SI{11000}{m} de profondeur) marque une limite \cle{convergente} où la lithosphère océanique plonge dans le manteau : c'est la \cle{subduction}.
\item \textbf{Bonne réponse : c.} La faille transformante est une limite où deux plaques \cle{coulissent horizontalement} l'une contre l'autre, sans création ni destruction de croûte ; la friction y provoque des séismes superficiels.
\end{enumerate}
\end{corrige}

\begin{corrige}[Exercice 2 — Vrai ou faux, avec justification]
\begin{enumerate}
\item \textbf{Vrai.} La croûte océanique se forme à l'axe de la dorsale puis s'en écarte progressivement : l'âge des basaltes augmente donc régulièrement avec la distance à l'axe, de façon symétrique.
\item \textbf{Faux.} Sous les dorsales, la fusion des péridotites résulte d'une \cle{diminution de pression} (décompression) lors de la remontée du manteau, à température quasi constante : c'est la \emph{fusion par décompression adiabatique}, et non une augmentation de température.
\item \textbf{Vrai.} En subduction, la plaque plongeante déshydrate ses minéraux hydratés vers \SI{100}{km} de profondeur ; l'eau libérée s'infiltre dans le manteau sus-jacent et abaisse le solidus des péridotites d'environ \SI{400}{\celsius} : c'est la \emph{fusion par hydratation}.
\item \textbf{Faux.} Séismes et volcans se concentrent en ceintures étroites qui matérialisent les \cle{limites de plaques} (dorsales, fosses de subduction, failles transformantes) : leur répartition est un indice majeur de la tectonique des plaques.
\item \textbf{Faux.} Annoncer une date et un lieu est la \cle{prédiction} (et elle reste impossible pour les séismes). La \cle{prévention} consiste à réduire la vulnérabilité : constructions parasismiques, digues, plans d'évacuation, éducation des populations.
\end{enumerate}
\end{corrige}

\begin{corrige}[Exercice 3 — Définitions]
\begin{itemize}
\item \textbf{Dorsale océanique} : chaîne de montagnes sous-marine située au centre des océans, culminant vers \SI{-2500}{m}, correspondant à une limite divergente de plaques. C'est une zone d'accrétion où la remontée de magma basaltique fabrique en permanence une nouvelle croûte océanique.
\item \textbf{Rift} : vallée axiale étroite et profonde (environ \SI{1000}{m}) creusée au sommet d'une dorsale, formée par effondrement d'un compartiment entre deux failles normales. C'est le lieu de sortie des magmas et l'axe de l'expansion océanique.
\item \textbf{Subduction} : plongée d'une plaque lithosphérique (généralement océanique, dense) sous une autre plaque dans l'asthénosphère, au niveau d'une fosse océanique. Elle s'accompagne de séismes profonds et d'un volcanisme explosif.
\item \textbf{Tsunami} : train d'ondes gravitationnelles océaniques déclenché par une perturbation brutale du fond marin (séisme de subduction, glissement, éruption). Quasi imperceptible en pleine mer, il s'amplifie à l'approche des côtes en vagues dévastatrices.
\item \textbf{Aléa} : phénomène naturel potentiellement dangereux (séisme, éruption, inondation), caractérisé par sa nature, son intensité et sa probabilité d'occurrence dans une région donnée.
\item \textbf{Vulnérabilité} : degré de fragilité des enjeux humains, matériels et environnementaux exposés à un aléa (densité de population, qualité des constructions, niveau de préparation).
\item \textbf{Prévention} : ensemble des mesures prises \emph{avant} la catastrophe pour en réduire les effets : aménagement du territoire, normes de construction, digues, éducation et exercices d'évacuation. Le risque résulte de la combinaison : risque = aléa $\times$ vulnérabilité.
\end{itemize}
\end{corrige}

\begin{corrige}[Exercice 4 — Calcul direct : taux d'expansion]
\begin{enumerate}
\item Le demi-taux d'expansion est le rapport de la distance à l'axe sur l'âge du basalte :
\[ v = \frac{d}{t} = \frac{\SI{100}{km}}{\SI{5}{Ma}}. \]
Conversion : $\SI{100}{km} = 100 \times 10^{5}\ \text{cm} = 10^{7}\ \text{cm}$ et $\SI{5}{Ma} = 5 \times 10^{6}\ \text{ans}$.
\[ v = \frac{10^{7}}{5 \times 10^{6}} = \SI{2}{cm/an}. \]
\item Le taux d'expansion total compte l'écartement des deux flancs de la dorsale :
\[ v_{\text{total}} = 2 \times v = 2 \times 2 = \SI{4}{cm/an}. \]
\item À taux constant, $d = v \times t$ :
\[ d = \SI{2}{cm/an} \times 12 \times 10^{6}\ \text{ans} = 2{,}4 \times 10^{7}\ \text{cm} = \SI{240}{km}. \]
\end{enumerate}
\end{corrige}

\begin{attention}[Points de vigilance — Exercice 4]
Ne pas oublier le facteur 2 entre demi-taux et taux total ; convertir les kilomètres en centimètres \emph{avant} le calcul ($\SI{1}{km} = 10^{5}\ \text{cm}$) ; écrire les unités à chaque étape.
\end{attention}

\begin{corrige}[Exercice 5 — Exploitation d'un profil bathymétrique]
\begin{enumerate}
\item Les trois structures sont, de l'ouest vers l'est : la \cle{plaine abyssale} (vers \SI{-5000}{m}), la \cle{dorsale} (élévation large culminant vers \SI{-2500}{m}) et, à son sommet, le \cle{rift axial} (dépression de \SI{1000}{m} de profondeur relative), puis à nouveau la plaine abyssale.
\item Schéma légendé du profil :
\begin{popfigure}
\begin{tikzpicture}[scale=1]
\draw[thick, popblue] (-6,0) -- (-2.2,0) .. controls (-1.6,0) and (-1.3,2.2) .. (-0.55,2.5) -- (-0.25,1.5) -- (0.25,1.5) -- (0.55,2.5) .. controls (1.3,2.2) and (1.6,0) .. (2.2,0) -- (6,0);
\draw[dashed, popmuted] (-6,2.5) -- (6,2.5) node[right, popdark]{\scriptsize niveau marin (0 m)};
\draw[<->, popink] (5.4,2.5) -- (5.4,0) node[midway, right]{\scriptsize \SI{-5000}{m}};
\draw[<->, popink] (1.9,2.5) -- (1.9,2.42);
\node[popdark] at (1.35,2.9) {\scriptsize \SI{-2500}{m}};
\draw[->, poporange, thick] (0,3.4) -- (0,1.7);
\node[poporange] at (0,3.7) {\scriptsize axe de la dorsale};
\node[popdark] at (-4,0.35) {\scriptsize plaine abyssale};
\node[popdark] at (4,0.35) {\scriptsize plaine abyssale};
\node[popdark] at (-1.35,1.7) {\scriptsize dorsale};
\node[popdark] at (0,1.1) {\scriptsize rift};
\end{tikzpicture}
\legende{Profil bathymétrique transocéanique perpendiculaire à une dorsale : 1. plaine abyssale (\SI{-5000}{m}) ; 2. dorsale (\SI{-2500}{m}) ; 3. rift axial (dépression de \SI{1000}{m}) ; 4. axe de la dorsale (traité vertical). Échelle verticale fortement exagérée.}
\end{popfigure}
\item La symétrie s'explique par le fait que la croûte océanique se forme à l'axe puis s'écarte \emph{de part et d'autre} à la même vitesse : les deux flancs subissent le même refroidissement, la même contraction thermique et le même enfoncement isostatique en s'éloignant de l'axe.
\end{enumerate}
\end{corrige}

\begin{corrige}[Exercice 6 — Taux d'expansion : tableau et courbe]
\begin{enumerate}
\item Courbe âge = $f$(distance) : les points s'alignent sur une droite passant par l'origine.
\begin{popfigure}
\begin{tikzpicture}
\begin{axis}[width=0.85\linewidth, height=6cm, xlabel={Distance à l'axe (km)}, ylabel={Âge des basaltes (Ma)}, xmin=0, xmax=850, ymin=0, ymax=9, grid=both, grid style={popline!40}, legend pos=north west]
\addplot[popcurve, domain=0:800, samples=50]{x/100};
\addplot[only marks, mark=*, poporange] coordinates {(0,0) (200,2) (400,4) (600,6) (800,8)};
\legend{droite d'ajustement, points expérimentaux}
\end{axis}
\end{tikzpicture}
\legende{Âge des basaltes en fonction de la distance à l'axe de la dorsale : l'alignement des points montre un taux d'expansion constant.}
\end{popfigure}
\item \textbf{Graphiquement}, la pente de la droite vaut par exemple $\dfrac{8\ \text{Ma}}{800\ \text{km}}$. \textbf{Par le calcul} :
\[ v = \frac{d}{t} = \frac{\SI{200}{km}}{\SI{2}{Ma}} = \frac{200 \times 10^{5}\ \text{cm}}{2 \times 10^{6}\ \text{ans}} = \frac{2 \times 10^{7}}{2 \times 10^{6}} = \SI{10}{cm/an}. \]
Le demi-taux d'expansion est de \SI{10}{cm/an} (taux total : \SI{20}{cm/an}, valeur rapide typique du Pacifique).
\item Les basaltes les plus anciens sont à mi-largeur de l'axe, soit $d = \SI{6000}{km}$ :
\[ t = \frac{d}{v} = \frac{6000 \times 10^{5}\ \text{cm}}{\SI{10}{cm/an}} = 6 \times 10^{7}\ \text{ans} = \SI{60}{Ma}. \]
\item Aucune croûte océanique de plus de 200 Ma n'existe car la lithosphère océanique, devenue froide et dense en vieillissant, finit toujours par \cle{plonger en subduction} et être recyclée dans le manteau : les océans se créent et se détruisent en permanence, contrairement aux continents, plus légers, qui s'accumulent.
\end{enumerate}
\end{corrige}

\begin{corrige}[Exercice 7 — Fusion des péridotites : lecture d'un diagramme]
\begin{enumerate}
\item Le \cle{solidus} d'une roche est la courbe du diagramme pression-température qui sépare le domaine où la roche est totalement solide du domaine où elle commence à fondre : c'est la température de début de fusion pour chaque pression.
\item \textbf{Sous une dorsale} : le manteau chaud remonte par convection ; au cours de cette remontée quasi adiabatique, la température reste presque constante mais la \cle{pression diminue fortement}. Le géotherme de dorsale coupe alors le solidus des péridotites \emph{anhydres} : c'est la \textbf{fusion par décompression}. Le taux de fusion partielle est élevé (20 à \SI{30}{\percent}).
\item \textbf{En subduction} : la plaque plongeante transporte des minéraux hydratés (amphiboles, serpentines) qui se déstabilisent vers \SI{100}{km} de profondeur et libèrent de l'eau. Cette eau monte dans le coin mantellique et abaisse le solidus des péridotites d'environ \SI{400}{\celsius} : le géotherme local, pourtant froid, atteint le solidus \emph{hydraté} : c'est la \textbf{fusion par hydratation} (ou par abaissement du point de fusion).
\item \textbf{Comparaison} : sous les dorsales, la fusion partielle importante produit des \cle{basaltes fluides}, pauvres en silice, donnant un volcanisme effusif (épanchements, pillow-lavas). En subduction, les magmas s'enrichissent en silice et en eau, deviennent \cle{visqueux} ; les gaz ne peuvent s'échapper facilement : le volcanisme est \emph{explosif} et dangereux (nuées ardentes, panaches de cendres).
\end{enumerate}
\end{corrige}

\begin{corrige}[Exercice 8 — Classer : prévision, prédiction ou prévention]
\begin{enumerate}
\item \textbf{Prévention} : les digues et canaux sont des aménagements permanents qui réduisent la vulnérabilité des quartiers avant toute crue.
\item \textbf{Prévision} : la surveillance continue (sismomètres) vise à détecter les signes avant-coureurs afin d'anticiper l'événement, sans en fixer la date.
\item \textbf{Prévention} : la carte des zones à risque et l'interdiction de construire relèvent de l'aménagement du territoire, qui réduit l'exposition des populations.
\item \textbf{Prédiction} : l'alerte tsunami annonce un événement imminent, localisé et daté (à quelques heures près), grâce au décalage entre l'onde sismique et l'onde océanique.
\item \textbf{Prévention} : les exercices d'évacuation préparent les populations (éducation au risque) pour réduire la vulnérabilité.
\item \textbf{Prévision} : la mesure de la déformation du sol et des émissions de gaz est une surveillance des paramètres précurseurs d'une éruption, sans annonce d'une date certaine.
\end{enumerate}
\end{corrige}

\begin{attention}[Distinction à maîtriser — Exercice 8]
Prévision = surveillance et anticipation sans date précise ; prédiction = annonce d'un événement daté et localisé (rarement possible, sauf tsunami et parfois éruptions) ; prévention = réduction durable de la vulnérabilité.
\end{attention}

\begin{corrige}[Exercice 9 — Type Bac : le plancher océanique atlantique]
\begin{enumerate}
\item Le profil traverse, des continents vers le centre : le \cle{plateau continental}, le \cle{talus continental}, la \cle{plaine abyssale} (\SI{-5000}{m}), puis la \cle{dorsale médio-atlantique} (\SI{-2500}{m}) creusée du \cle{rift axial}.
\begin{popfigure}
\begin{tikzpicture}[scale=1]
\draw[thick, popblue] (-6,2.2) -- (-5,2.2) -- (-4.2,0) -- (-2.2,0) .. controls (-1.5,0) and (-1.2,2.0) .. (-0.5,2.3) -- (-0.2,1.4) -- (0.2,1.4) -- (0.5,2.3) .. controls (1.2,2.0) and (1.5,0) .. (2.2,0) -- (4.2,0) -- (5,2.2) -- (6,2.2);
\draw[dashed, popmuted] (-6,2.5) -- (6,2.5) node[right, popdark]{\scriptsize niveau marin};
\node[popdark] at (-5.5,2.55) {\scriptsize plateau continental};
\node[popdark] at (-4.6,1.15) {\scriptsize talus};
\node[popdark] at (-3.2,0.3) {\scriptsize plaine abyssale};
\node[popdark] at (0,1.0) {\scriptsize rift};
\node[popdark] at (-1.3,1.75) {\scriptsize dorsale};
\draw[->, poporange, thick] (0,3.3) -- (0,1.6);
\node[poporange] at (0,3.6) {\scriptsize axe};
\end{tikzpicture}
\legende{Profil bathymétrique de l'Atlantique : 1. plateau continental ; 2. talus continental ; 3. plaine abyssale (\SI{-5000}{m}) ; 4. dorsale (\SI{-2500}{m}) ; 5. rift axial. Échelle verticale exagérée.}
\end{popfigure}
\item Demi-taux d'expansion :
\[ v = \frac{d}{t} = \frac{\SI{300}{km}}{\SI{10}{Ma}} = \frac{300 \times 10^{5}\ \text{cm}}{10 \times 10^{6}\ \text{ans}} = \frac{3 \times 10^{7}}{10^{7}} = \SI{3}{cm/an}. \]
Taux d'expansion total : $v_{\text{total}} = 2 \times 3 = \SI{6}{cm/an}$.
\item \textbf{Formation continue d'une croûte qui s'écarte} : au document 2, l'âge des basaltes est nul à l'axe et augmente régulièrement avec la distance (10 Ma à \SI{300}{km}, 30 Ma à \SI{900}{km}), de façon \emph{identique des deux côtés} : la croûte se forme donc à l'axe puis s'en éloigne symétriquement. Au document 3, les anomalies magnétiques forment des bandes parallèles et symétriques : chaque inversion du champ magnétique terrestre a été enregistrée simultanément dans la croûte en formation à l'axe, puis les bandes ont été écartées de part et d'autre. La dorsale est donc bien une zone d'\cle{accrétion} continue.
\item \textbf{Fonctionnement du rift} : l'écartement des plaques fracture la lithosphère selon des \cle{failles normales} (extension) ; le compartiment central s'effondre, formant la vallée du rift. La décompression du manteau sous-jacent provoque la fusion partielle des péridotites ; le magma basaltique remonte par les fissures, se solidifie en basaltes et gabbros, comblant le vide créé par l'écartement : ainsi naît le nouveau plancher océanique, qui s'éloigne ensuite de l'axe.
\end{enumerate}
\textbf{Barème indicatif (sur 10)} : identification des structures et schéma légendé : 3 pts ; calcul du demi-taux avec conversions et taux total : 2 pts ; exploitation croisée des documents 2 et 3 (symétrie des âges et des anomalies) : 3 pts ; fonctionnement du rift (failles normales, décompression, magma) : 2 pts.
\end{corrige}

\begin{attention}[Points de vigilance — Exercice 9]
Citer explicitement la \emph{symétrie} des âges et des anomalies magnétiques (c'est l'argument central) ; distinguer demi-taux et taux total ; préciser que la fusion sous la dorsale est une fusion par \emph{décompression} et non par élévation de température.
\end{attention}

\begin{corrige}[Exercice 10 — Type Bac : séismes et volcans, indices des limites de plaques]
\begin{enumerate}
\item Il s'agit d'une limite \cle{convergente avec subduction}. Justification : la fosse océanique profonde (\SI{8000}{m}) marque l'entrée en subduction ; les foyers sismiques s'enfoncent en biais sous le continent, des séismes superficiels (moins de \SI{70}{km}) jusqu'à \SI{600}{km} de profondeur, selon un plan incliné : cette distribution matérialise la plaque plongeante, siège des contraintes et des ruptures.
\item Ce plan incliné est le \cle{plan de Wadati-Benioff} : il matérialise la lithosphère océanique plongeant dans l'asthénosphère ; la profondeur croissante des foyers indique le sens de plongée de la plaque.
\item \textbf{Origine du magma} : vers \SI{100}{km} de profondeur, les minéraux hydratés de la croûte basaltique plongeante se déstabilisent et libèrent de l'eau (document 3) ; cette eau monte dans le manteau sus-jacent et abaisse le solidus des péridotites d'environ \SI{400}{\celsius}, provoquant leur fusion partielle par hydratation. Le magma produit s'enrichit en silice et en eau au cours de sa remontée : il devient \cle{visqueux} (document 2). Les gaz dissous ne peuvent s'échapper d'un magma visqueux : la pression augmente jusqu'à la rupture brutale, d'où des éruptions \emph{explosives}.
\item \textbf{Cumul des aléas} : la convergence des plaques provoque des séismes puissants (blocage puis rupture brutale le long du plan de subduction) ; la fusion du manteau hydraté alimente un volcanisme explosif ; enfin, les séismes sous-marins de subduction, en déplaçant brutalement le fond océanique, peuvent déclencher des \cle{tsunamis}. \textbf{Mesures de prévention} : normes de construction parasismiques et exercices d'évacuation ; cartographie des zones à risque volcanique et plans d'évacuation ; système d'alerte aux tsunamis et signalétique des itinéraires de repli en hauteur (deux mesures suffisent).
\end{enumerate}
\textbf{Barème indicatif (sur 10)} : type de limite justifié par la fosse et la disposition des foyers : 2,5 pts ; plan de Wadati-Benioff : 1,5 pt ; origine du magma (déshydratation, abaissement du solidus) et explosivité (viscosité, gaz piégés) : 3,5 pts ; cumul des aléas et deux mesures de prévention : 2,5 pts.
\end{corrige}

\begin{attention}[Points de vigilance — Exercice 10]
Ne pas se contenter de dire « il y a des séismes profonds » : relier explicitement la profondeur croissante des foyers à la plongée de la plaque ; citer la valeur de \SI{400}{\celsius} d'abaissement du solidus ; pour l'explosivité, enchaîner logiquement : silice $\Rightarrow$ viscosité $\Rightarrow$ gaz piégés $\Rightarrow$ explosion.
\end{attention}

\begin{corrige}[Exercice 11 — Type Bac : la catastrophe du lac Nyos et la gestion des risques au Cameroun]
\begin{enumerate}
\item \textbf{Aléa} : l'\cle{éruption limnique} (dégazage brutal d'un lac de cratère). \textbf{Mécanisme} : le \ce{CO2} d'origine magmatique s'accumule lentement, dissous, dans les eaux profondes du lac, maintenues stratifiées et stables ; lorsqu'une perturbation (glissement de terrain, refroidissement des eaux superficielles) déstabilise cette stratification, les eaux profondes remontent brutalement, la pression diminue et le gaz se déssature instantanément, comme dans une bouteille d'eau gazeuse qu'on débouche : un nuage de \ce{CO2}, plus dense que l'air, se répand alors au ras du sol dans les vallées et asphyxie les êtres vivants.
\item \textbf{Vulnérabilité élevée} : les villages étaient installés dans les vallées et à proximité immédiate du lac, c'est-à-dire précisément là où le gaz dense s'accumule ; le phénomène était inconnu des populations (aucune mémoire, aucune consigne) ; la catastrophe s'est produite de nuit, sans signe avant-coureur perceptible, le \ce{CO2} étant incolore et inodore ; l'isolement des villages a retardé les secours.
\item \textbf{Classement des mesures (document 3)} :
\begin{itemize}
\item \textbf{Prévention} : le dégazage continu par tuyaux, qui supprime durablement l'aléa en évacuant progressivement le \ce{CO2} ; la sensibilisation et le relogement des populations en hauteur, qui réduisent leur exposition.
\item \textbf{Prévision} : les capteurs mesurant en permanence la concentration en gaz dissous, qui surveillent l'évolution du phénomène pour anticiper tout danger, sans prétendre dater un événement.
\item \textbf{Prédiction} : aucune mesure du document 3 n'en relève — et c'est logique : le dégazage a supprimé l'essentiel de l'aléa, rendant inutile l'annonce d'un événement daté.
\end{itemize}
\item \textbf{Texte de sensibilisation (exemple attendu)} : « Habitants des quartiers bas de Douala, les pluies diluviennes peuvent transformer nos rues en torrents. \emph{Avant la crue} : nettoyez régulièrement les caniveaux devant vos maisons, ne jetez jamais d'ordures dans les canaux, surélevez vos prises électriques et conservez documents et médicaments en hauteur ; repérez le point de rassemblement et l'itinéraire vers les zones hautes. \emph{Pendant la crue} : coupez l'électricité, ne traversez jamais une eau qui monte, même à pied ou en voiture, éloignez les enfants des canaux et montez à l'étage ou gagnez les zones hautes ; écoutez la radio pour les consignes. \emph{Après la crue} : ne touchez pas aux fils électriques tombés, faites bouillir l'eau de boisson, nettoyez et désinfectez les habitations, signalez les blessés et les dégâts aux autorités. La prévention est l'affaire de tous : un quartier propre et organisé est un quartier protégé. »
\end{enumerate}
\textbf{Barème indicatif (sur 10)} : aléa et mécanisme de l'éruption limnique : 2,5 pts ; analyse de la vulnérabilité : 2 pts ; classement justifié prévision/prévention : 2,5 pts ; texte de sensibilisation structuré avant/pendant/après : 3 pts.
\end{corrige}

\begin{attention}[Points de vigilance — Exercice 11]
Expliquer le rôle de la \emph{stratification} du lac et de la \emph{densité} du \ce{CO2} (ce sont les deux clés du mécanisme) ; dans le classement, justifier chaque affectation par la définition des trois notions ; dans le texte de sensibilisation, distinguer explicitement les trois phases (avant, pendant, après) et employer un ton adapté au public visé.
\end{attention}

\newpage

\coursec{Fiche bilan}

\begin{popfigure}
\centering
\begin{center}\tcbox[colback=popblue,colframe=popblue,coltext=white,arc=9pt,boxsep=4pt,left=6pt,right=6pt]{\bfseries\small Réduction risques catastrophes naturelles}\end{center}
\vspace{-2pt}
\begin{tcbraster}[raster columns=3,raster equal height=rows,raster column skip=6pt,raster row skip=6pt,enhanced,blankest]
\begin{tcolorbox}[enhanced,colback=white,colframe=popblue,boxrule=0.9pt,arc=3pt,title={Morphologie fonds océaniques},fonttitle=\bfseries\scriptsize,coltitle=white,colbacktitle=popblue,left=4pt,right=4pt,top=2pt,bottom=2pt,boxsep=1pt,toptitle=1.5pt,bottomtitle=1.5pt,halign=left]
\begin{itemize}[nosep,leftmargin=*,itemsep=1pt,font=\scriptsize]
\item Dorsales, fosses
\item Bathymétrie
\end{itemize}
\end{tcolorbox}
\begin{tcolorbox}[enhanced,colback=white,colframe=poporange,boxrule=0.9pt,arc=3pt,title={Rift \& Expansion},fonttitle=\bfseries\scriptsize,coltitle=white,colbacktitle=poporange,left=4pt,right=4pt,top=2pt,bottom=2pt,boxsep=1pt,toptitle=1.5pt,bottomtitle=1.5pt,halign=left]
\begin{itemize}[nosep,leftmargin=*,itemsep=1pt,font=\scriptsize]
\item Accrétion, MORB
\item Anomalies magnétiques
\end{itemize}
\end{tcolorbox}
\begin{tcolorbox}[enhanced,colback=white,colframe=popgreen,boxrule=0.9pt,arc=3pt,title={Magmas \& Péridotites},fonttitle=\bfseries\scriptsize,coltitle=white,colbacktitle=popgreen,left=4pt,right=4pt,top=2pt,bottom=2pt,boxsep=1pt,toptitle=1.5pt,bottomtitle=1.5pt,halign=left]
\begin{itemize}[nosep,leftmargin=*,itemsep=1pt,font=\scriptsize]
\item Fusion partielle
\item Rôle de l'eau
\end{itemize}
\end{tcolorbox}
\begin{tcolorbox}[enhanced,colback=white,colframe=poppurple,boxrule=0.9pt,arc=3pt,title={Convergence \& Coulissage},fonttitle=\bfseries\scriptsize,coltitle=white,colbacktitle=poppurple,left=4pt,right=4pt,top=2pt,bottom=2pt,boxsep=1pt,toptitle=1.5pt,bottomtitle=1.5pt,halign=left]
\begin{itemize}[nosep,leftmargin=*,itemsep=1pt,font=\scriptsize]
\item Subduction, collision
\item Transformante
\end{itemize}
\end{tcolorbox}
\begin{tcolorbox}[enhanced,colback=white,colframe=poppink,boxrule=0.9pt,arc=3pt,title={Catastrophes associées},fonttitle=\bfseries\scriptsize,coltitle=white,colbacktitle=poppink,left=4pt,right=4pt,top=2pt,bottom=2pt,boxsep=1pt,toptitle=1.5pt,bottomtitle=1.5pt,halign=left]
\begin{itemize}[nosep,leftmargin=*,itemsep=1pt,font=\scriptsize]
\item Séismes, tsunamis
\item Volcanisme, glissements
\end{itemize}
\end{tcolorbox}
\begin{tcolorbox}[enhanced,colback=white,colframe=popgold,boxrule=0.9pt,arc=3pt,title={Gestion des risques},fonttitle=\bfseries\scriptsize,coltitle=white,colbacktitle=popgold,left=4pt,right=4pt,top=2pt,bottom=2pt,boxsep=1pt,toptitle=1.5pt,bottomtitle=1.5pt,halign=left]
\begin{itemize}[nosep,leftmargin=*,itemsep=1pt,font=\scriptsize]
\item Prévision, prédiction
\item Prévention, sensibilisation
\end{itemize}
\end{tcolorbox}
\end{tcbraster}

\legende{Carte mentale récapitulative de la leçon 15 : structures océaniques, dynamique des plaques, genèse des magmas, risques géologiques et stratégies de gestion.}
\end{popfigure}

\begin{bilan}
\courssub{Morphologie des fonds océaniques \& Rift}
\begin{itemize}
    \item \cle{Dorsale médio-océanique} : chaîne sous-marine continue ($\sim 60\,000$ km), axe \cle{rift} (vallée d'effondrement, failles normales), forte chaleur, hydrothermalisme (\emph{black smokers}).
    \item \cle{Plancher océanique} : croûte basaltique (MORB) + gabbros + péridotites (mantle), épaisseur $\sim 7$ km. Âge croissant de l'axe vers les continents (datation magnétique).
    \item \cle{Expansion} : séparation des plaques ($\sim 1$--$15$ cm/an), création de croûte neuve par fusion adiabatique de la péridotite (décompression).
    \item \cle{Anomalies magnétiques} : symétriques par rapport à l'axe $\rightarrow$ preuve de l'expansion (inversions du champ magnétique enregistrées dans la magnétite des basaltes).
\end{itemize}

\courssub{Naissance des magmas : données expérimentales (fusion péridotite)}
\begin{center}
\begin{tabularx}{\linewidth}{|l|X|X|}\hline
\rowcolor{popblueL}
\textbf{Contexte} & \textbf{Conditions P-T-X} & \textbf{Magma produit} \\\hline
Zone d'accrétion (rift) & Décompression adiabatique ($P \downarrow$, $T$ cst), \textbf{sec} & \textbf{MORB} (basalte tholéiitique), basique, pauvre en eau \\\hline
Zone de subduction & $T \uparrow$ (friction), apport \textbf{eau} (déshydratation plaque plongeante), $P$ élevé & \textbf{Andésites, rhyolites} (calco-alcalins), acides, riches en eau, explosifs \\\hline
Point chaud (panache) & $T \uparrow$ (panache mantellique), $P \downarrow$ & Basaltes alcalins (OIB), \emph{hotspot} (ex. Mont Cameroun) \\\hline
\end{tabularx}
\end{center}
\emph{Expérience clé} : Fusion partielle de péridotite sèche $\rightarrow$ basalte à $\sim 1\,300$ °C / 1,5 GPa. Avec $0,1$--$0,5$ \% $H_2O$ $\rightarrow$ solidus abaissé de $\sim 200$ °C $\rightarrow$ magma andésitique.

\courssub{Limites de plaques convergentes \& de coulissage}
\begin{itemize}
    \item \cle{Subduction océanique-océanique} : arc insulaire (ex. Japon), fosses profondes ($> 10$ km), sismicité de Benioff (foyers profonds $< 700$ km), volcanisme explosif (andésites).
    \item \cle{Subduction océanique-continentale} : chaîne de montagnes côtière (Andes), volcanisme andésitique/rhyolitique, séismes profonds, raccourcissement crustal.
    \item \cle{Collision continentale} : pas de subduction (faible densité), épaississement crustal (Himalaya), séismes superficiels profonds, pas de volcanisme actif.
    \item \cle{Faille transformante (coulissage)} : plaques glissent latéralement (ex. faille de San Andreas, transformantes dorsales), séismes \textbf{superficiels} ($< 20$ km), pas de volcanisme, morphologie : vallées rectilignes, décaltage dorsales.
\end{itemize}

\courssub{Catastrophes naturelles associées}
\begin{center}
\begin{tabularx}{\linewidth}{|l|X|X|}\hline
\rowcolor{poporangeL}
\textbf{Phénomène} & \textbf{Mécanisme principal} & \textbf{Exemple camerounais / régional} \\\hline
Séisme & Rupture brutale faille (libération contrainte élastique), ondes P, S, Rayleigh & Mont Cameroun (volcano-tectoniques), faille de Sanaga \\
Tsunami & Déplacement vertical fond marin (séisme subduction, glissement sous-marin) & Risque golfe de Guinée (séismes atlantiques) \\
Volcanisme & Remontée magma (point chaud, subduction, rift) & \textbf{Mont Cameroun} (point chaud, éruptions 1999, 2000), Lac Nyos (dégazage $CO_2$ 1986) \\
Glissement de terrain & Instabilité pente (pluie, séisme, déforestation), saturation eau & Monts Mandara, quartiers pentus Yaoundé/Douala \\
Inondation & Pluies intenses, bassin versant imperméabilisé, montée nappes & Plaines de la Logone, quartiers bas Douala (Wouri) \\\hline
\end{tabularx}
\end{center}

\courssub{Gestion des risques : Prévision -- Prédiction -- Prévention}
\begin{methode}[Distinction terminologique rigoureuse]
\begin{enumerate}
    \item \textbf{Prévision} : estimation probabiliste à long/moyen terme (\emph{où} et \emph{quelle intensité} ?). Ex : carte d'aléa sismique, zone de récurrence volcanique.
    \item \textbf{Prédiction} : annonce précise à court terme (\emph{quand}, \emph{où}, \emph{quelle magnitude} ?). \textbf{Actuellement impossible pour les séismes} (pas de précurseur fiable universel). Possible pour certains volcans (déformation, gaz, sismicité) et crues (modèles hydro-météorologiques).
    \item \textbf{Prévention} : mesures pour réduire la \cle{vulnérabilité} (enjeux exposés) : urbanisation (PLU, zones inconstructibles), construction parasismique (Eurocode 8 / normes locales), education/simulation (ex. exercices tsunami), surveillance (réseaux sismiques, GPS, gaz volcaniques), reboisement pentes.
\end{enumerate}
\end{methode}
\textbf{Équation risque} : $R = A \times V \times E$ (Aléa $\times$ Vulnérabilité $\times$ Enjeu). Réduire $R$ $\equiv$ agir sur $V$ et $E$ (car $A$ souvent immuable).
\end{bilan}

\begin{aretenir}[Checklist avant le Bac]
\begin{itemize}
    \item $\square$ Définir et localiser sur une carte : dorsale, rift, fosse, arc insulaire, faille transformante, point chaud.
    \item $\square$ Expliquer l'expansion océanique : rôle de la convection mantellique, fusion adiabatique, création croûte océanique, anomalies magnétiques symétriques.
    \item $\square$ Utiliser les courbes solidus/liquidus péridotite (sec vs humide) pour justifier la nature des magmas (MORB vs andésites) selon le géodynamisme.
    \item $\square$ Comparer les 3 types de convergence (océan/océan, océan/continent, continent/continent) : structures, sismicité, volcanisme, topographie.
    \item $\square$ Caractériser une faille transformante : cinématique, sismicité superficielle, absence volcanisme, morphologie.
    \item $\square$ Lier chaque catastrophe (séisme, tsunami, éruption, glissement, inondation) à son contexte géodynamique précis (type de limite de plaque, déclencheur).
    \item $\square$ Différencier sans ambiguïté \textbf{prévision} (probabilité long terme), \textbf{prédiction} (certitude court terme) et \textbf{prévention} (réduction vulnérabilité).
    \item $\square$ Citer 3 mesures de prévention sismique (règles construction, microzonage, formation) et 3 mesures volcaniques (surveillance, plan d'évacuation, exclusion zones).
    \item $\square$ Appliquer la formule $R = A \times V \times E$ sur un cas concret camerounais (ex. risque éruption Mont Cameroun / risque glissement Monts Mandara / risque inondation Douala).
    \item $\square$ Analyser un document bathymétrique (profil, carte) : identifier structures, sens expansion, âge relatif croûte.
    \item $\square$ Construire un argumentaire pour un outil de sensibilisation (affiche, maquette, vidéo) ciblant une population locale sur un risque identifié.
\end{itemize}
\end{aretenir}

\findecours

\end{document}
